%% Generated by Sphinx.
\def\sphinxdocclass{report}
\documentclass[letterpaper,10pt,english]{sphinxmanual}
\ifdefined\pdfpxdimen
   \let\sphinxpxdimen\pdfpxdimen\else\newdimen\sphinxpxdimen
\fi \sphinxpxdimen=.75bp\relax
\ifdefined\pdfimageresolution
    \pdfimageresolution= \numexpr \dimexpr1in\relax/\sphinxpxdimen\relax
\fi
%% let collapsible pdf bookmarks panel have high depth per default
\PassOptionsToPackage{bookmarksdepth=5}{hyperref}

\PassOptionsToPackage{booktabs}{sphinx}
\PassOptionsToPackage{colorrows}{sphinx}

\PassOptionsToPackage{warn}{textcomp}
\usepackage[utf8]{inputenc}
\ifdefined\DeclareUnicodeCharacter
% support both utf8 and utf8x syntaxes
  \ifdefined\DeclareUnicodeCharacterAsOptional
    \def\sphinxDUC#1{\DeclareUnicodeCharacter{"#1}}
  \else
    \let\sphinxDUC\DeclareUnicodeCharacter
  \fi
  \sphinxDUC{00A0}{\nobreakspace}
  \sphinxDUC{2500}{\sphinxunichar{2500}}
  \sphinxDUC{2502}{\sphinxunichar{2502}}
  \sphinxDUC{2514}{\sphinxunichar{2514}}
  \sphinxDUC{251C}{\sphinxunichar{251C}}
  \sphinxDUC{2572}{\textbackslash}
\fi
\usepackage{cmap}
\usepackage[T1]{fontenc}
\usepackage{amsmath,amssymb,amstext}
\usepackage{babel}



\usepackage{tgtermes}
\usepackage{tgheros}
\renewcommand{\ttdefault}{txtt}



\usepackage[Bjarne]{fncychap}
\usepackage[,numfigreset=1,mathnumfig]{sphinx}

\fvset{fontsize=auto}
\usepackage{geometry}


% Include hyperref last.
\usepackage{hyperref}
% Fix anchor placement for figures with captions.
\usepackage{hypcap}% it must be loaded after hyperref.
% Set up styles of URL: it should be placed after hyperref.
\urlstyle{same}


\usepackage{sphinxmessages}
\setcounter{tocdepth}{0}


\addto\captionsenglish{\renewcommand{\bibname}{References}}
    \graphicspath{ {./_video_thumbnail/}{./} }
\newcommand{\sphinxcontribyoutube}[3]{\begin{quote}\begin{center}\fbox{\url{#1#2#3}}\end{center}\end{quote}}


\title{CLiC Documentation}
\date{Jul 15, 2026}
\release{v2.3.1}
\author{Michaela Mahlberg}
\newcommand{\sphinxlogo}{\sphinxincludegraphics{clic-logo.pdf}\par}
\renewcommand{\releasename}{Release}
\makeindex
\begin{document}

\ifdefined\shorthandoff
  \ifnum\catcode`\=\string=\active\shorthandoff{=}\fi
  \ifnum\catcode`\"=\active\shorthandoff{"}\fi
\fi

\pagestyle{empty}
\sphinxmaketitle
\pagestyle{plain}
\sphinxtableofcontents
\pagestyle{normal}
\phantomsection\label{\detokenize{index::doc}}


\sphinxstepscope


\chapter{Introduction}
\label{\detokenize{introduction:introduction}}\label{\detokenize{introduction::doc}}
\sphinxAtStartPar
Welcome to CLiC. CLiC is a web app to support the computer\sphinxhyphen{}assisted
study of narrative fiction. It started off with the
\sphinxhref{https://www.birmingham.ac.uk/schools/edacs/departments/englishlanguage/research/projects/clic/index.aspx}{CLiC Dickens}
project and now contains a range of fictional texts from children’s
literature to German novels (all out of copyright).
\sphinxcontribyoutube{https://youtu.be/}{yv6OP3W2NDE}{}

\sphinxAtStartPar
This guide covers information on how to use CLiC, developer features,
information on the corpora and legal \& technical details. For inspiration
on what to do with CLiC, the
\sphinxhref{https://blog.bham.ac.uk/clic-dickens/}{CLiC blog} provides plenty
of examples. If you are a creative writer, you might want to try out the
\sphinxhref{https://blog.bham.ac.uk/clic-dickens/clic-creative/}{CLiC Creative}
resources.


\section{Citing CLiC}
\label{\detokenize{introduction:citing-clic}}\label{\detokenize{introduction:id1}}
\sphinxAtStartPar
When you use CLiC in your work, please cite CLiC like this:
\begin{quote}

\sphinxAtStartPar
Mahlberg, M., Stockwell, P., Wiegand, V. and Lentin, J. (2020)
CLiC 2.3. \sphinxstyleemphasis{Corpus Linguistics in Context}, available at:
\sphinxhref{https://clic-fiction.com}{clic\sphinxhyphen{}fiction.com} {[}Accessed: DATE{]}
\end{quote}


\section{Navigating the interface}
\label{\detokenize{introduction:navigating-the-interface}}\label{\detokenize{introduction:id2}}
\sphinxAtStartPar
The CLiC landing page shows a rotating summary of the CLiC corpora.
The sidebar menu on the right provides access to all analysis functions.
You can open and close it by clicking the menu icon in the top right
corner.

\begin{figure}[htbp]
\centering
\capstart

\noindent\sphinxincludegraphics{{figure-landing-page-2_0_0_19C_menu}.png}
\caption{Open and close the sidebar menu by clicking the menu icon in
the top right corner}\label{\detokenize{introduction:id3}}\label{\detokenize{introduction:figure-landing-page-menu}}\end{figure}

\sphinxAtStartPar
The sidebar contains the following analysis tabs:
\begin{itemize}
\item {} 
\sphinxAtStartPar
\sphinxstylestrong{Concordance} \textendash{} search for words and phrases

\item {} 
\sphinxAtStartPar
\sphinxstylestrong{FlexiConc} \textendash{} flexible reproducible concordance algorithms

\item {} 
\sphinxAtStartPar
\sphinxstylestrong{Subsets} \textendash{} explore text categories

\item {} 
\sphinxAtStartPar
\sphinxstylestrong{Clusters} \textendash{} find repeated sequences

\item {} 
\sphinxAtStartPar
\sphinxstylestrong{Keywords} \textendash{} compare corpora

\item {} 
\sphinxAtStartPar
\sphinxstylestrong{Counts} \textendash{} get word and chapter counts

\item {} 
\sphinxAtStartPar
\sphinxstylestrong{Texts} \textendash{} read full texts

\end{itemize}

\sphinxAtStartPar
Each tab is documented in its own section of this guide. Clicking the
CLiC logo returns you to the landing page.

\sphinxAtStartPar
The top row of the sidebar shows buttons common to all tabs:
\begin{itemize}
\item {} 
\sphinxAtStartPar
\sphinxstylestrong{Load} \textendash{} upload a previously exported CLiC CSV file to restore
your settings and tag annotations. The CSV can only be reimported
if it has not been modified; we recommend keeping the original copy
for reloading and saving any manual changes in a separate version.
Note that Load will replace any existing tags in CLiC with those
from the file (unlike Merge).

\item {} 
\sphinxAtStartPar
\sphinxstylestrong{Merge} \textendash{} add tags from a CSV file to any pre\sphinxhyphen{}existing tags. This
is useful when combining annotations from multiple researchers to
check how far tag sets overlap or differ.

\item {} 
\sphinxAtStartPar
\sphinxstylestrong{Save} \textendash{} save your results, settings and annotated tags as a CSV
file. The file includes a shareable link that replicates your search
settings.

\item {} 
\sphinxAtStartPar
\sphinxstylestrong{Clear} \textendash{} reset the search settings and any tag columns (identical
to clicking the CLiC logo).

\item {} 
\sphinxAtStartPar
\sphinxstylestrong{Help} \textendash{} open this user guide.

\end{itemize}

\sphinxAtStartPar
For both Load and Merge, the results must be compatible, i.e. based on
the same node word.


\section{Printing output}
\label{\detokenize{introduction:printing-output}}
\sphinxAtStartPar
To print CLiC output with colours and highlighting preserved, use your
browser’s print menu and ensure that background graphics are enabled
(in Chrome: More settings \textgreater{} Background graphics). Landscape format
tends to work best. You can print to a PDF file or directly to a
printer.

\begin{figure}[htbp]
\centering
\capstart

\noindent\sphinxincludegraphics{{figure-analysis-common-printing-settings}.png}
\caption{Enable background graphics in your browser’s print settings
to preserve colours}\label{\detokenize{introduction:id4}}\label{\detokenize{introduction:figure-printing-settings}}\end{figure}


\section{System requirements}
\label{\detokenize{introduction:system-requirements}}
\sphinxAtStartPar
CLiC runs in your web browser (you will need an internet connection).
We recommend recent versions of Chrome, Firefox or Safari. If the CLiC
homepage displays a red error message, please check whether your browser
needs updating. CLiC should also work on most Apple and Android mobile
devices.


\section{Acknowledgements}
\label{\detokenize{introduction:acknowledgements}}
\sphinxAtStartPar
The first CLiC project between the University of Birmingham and the
University of Nottingham was funded by the Arts and Humanities Research
Council (grant references AH/P504634/1 \& AH/K005146/1).

\sphinxAtStartPar
The ChiLit corpus was compiled as part of the GLARE project. GLARE was
funded by the European Commission within Marie Sklodowska\sphinxhyphen{}Curie Actions
(reference number EU 749521).

\sphinxAtStartPar
The latest FlexiConc functionality in CLiC was funded by an AHRC / DFG
project (grant references AH/X002047/1 \& AH/X002047/2).

\sphinxAtStartPar
Current CLiC development is supported by the Alexander von Humboldt
Foundation.

\sphinxAtStartPar
The user guide is a work in progress. Please get in touch via
\sphinxhref{mailto:dhss-kontakt@fau.de}{dhss\sphinxhyphen{}kontakt@fau.de} if you have further questions or suggestions for
improvement.

\sphinxstepscope


\chapter{The CLiC corpora}
\label{\detokenize{cliccorpora:the-clic-corpora}}\label{\detokenize{cliccorpora::doc}}
\begin{sphinxShadowBox}
\begin{itemize}
\item {} 
\sphinxAtStartPar
\phantomsection\label{\detokenize{cliccorpora:id11}}{\hyperref[\detokenize{cliccorpora:available-corpora}]{\sphinxcrossref{Available corpora}}}

\item {} 
\sphinxAtStartPar
\phantomsection\label{\detokenize{cliccorpora:id12}}{\hyperref[\detokenize{cliccorpora:selecting-texts-and-corpora}]{\sphinxcrossref{Selecting texts and corpora}}}

\item {} 
\sphinxAtStartPar
\phantomsection\label{\detokenize{cliccorpora:id13}}{\hyperref[\detokenize{cliccorpora:textual-subsets}]{\sphinxcrossref{Textual subsets}}}

\item {} 
\sphinxAtStartPar
\phantomsection\label{\detokenize{cliccorpora:id14}}{\hyperref[\detokenize{cliccorpora:versioning-and-reproducibility}]{\sphinxcrossref{Versioning and reproducibility}}}

\item {} 
\sphinxAtStartPar
\phantomsection\label{\detokenize{cliccorpora:id15}}{\hyperref[\detokenize{cliccorpora:how-to-cite-corpora}]{\sphinxcrossref{How to cite corpora}}}

\end{itemize}
\end{sphinxShadowBox}

\sphinxAtStartPar
CLiC contains several pre\sphinxhyphen{}selected corpora of novels, which can be
selected individually and combined freely for analysis in any of the
CLiC tools.


\section{Available corpora}
\label{\detokenize{cliccorpora:available-corpora}}\begin{itemize}
\item {} 
\sphinxAtStartPar
\sphinxstylestrong{DNov} (15 texts) \textendash{} Dickens’s Novels. The complete novels of
Charles Dickens.

\item {} 
\sphinxAtStartPar
\sphinxstylestrong{19C} (29 texts) \textendash{} 19th Century Reference Corpus. A selection of
19th century novels by other authors, originally compiled as a reference corpus
in comparisons with Dickens.

\item {} 
\sphinxAtStartPar
\sphinxstylestrong{ChiLit} (71 texts) \textendash{} 19th Century Children’s Literature Corpus.
Compiled by Anna Cermakova for the GLARE project.

\item {} 
\sphinxAtStartPar
\sphinxstylestrong{ArTs} (31 texts) \textendash{} Additional Requested Texts. Includes
additional GCSE and A\sphinxhyphen{}Level titles and other individually requested
texts. This corpus was not designed to be analysed as a whole, but
rather to make individual texts available in CLiC.

\item {} 
\sphinxAtStartPar
\sphinxstylestrong{DE19} (33 texts) \textendash{} German 19th Century Reference Corpus. Compiled in the \sphinxstyleemphasis{Reading Concordances in the 21st Century} project.

\end{itemize}

\sphinxAtStartPar
An overview of the texts in each corpus (with word counts) can also be
generated dynamically using the {\hyperref[\detokenize{clicanalysis/counts:counts}]{\sphinxcrossref{\DUrole{std,std-ref}{Counts}}}} tab. For a full list of
titles, see {\hyperref[\detokenize{appendices:a-list-of-texts-available-in-clic}]{\sphinxcrossref{\DUrole{std,std-ref}{Appendix: Texts in CLiC}}}} and the \sphinxcite{references:github-corpora} repository.


\section{Selecting texts and corpora}
\label{\detokenize{cliccorpora:selecting-texts-and-corpora}}
\sphinxAtStartPar
The texts can be selected individually and combined freely for analysis
in any of the CLiC tools. You can also select one of the pre\sphinxhyphen{}selected
corpora listed above. For details on how to select texts and corpora,
see the {\hyperref[\detokenize{introduction:navigating-the-interface}]{\sphinxcrossref{\DUrole{std,std-ref}{Introduction}}}} and the
documentation for each analysis tab.

\sphinxAtStartPar
You can also select an author\sphinxhyphen{}based corpus from the drop\sphinxhyphen{}down. For
example, typing \sphinxstyleemphasis{austen} into the textbox (which is not case\sphinxhyphen{}sensitive)
gives you the option of selecting all books by Jane Austen at once, as
illustrated in \hyperref[\detokenize{cliccorpora:figure-corpora-authorbased}]{Fig.\@ \ref{\detokenize{cliccorpora:figure-corpora-authorbased}}}.

\begin{figure}[htbp]
\centering
\capstart

\noindent\sphinxincludegraphics{{figure-corpora-authorbased}.png}
\caption{Example of creating an author\sphinxhyphen{}based corpus:
selecting all of Jane Austen’s novels}\label{\detokenize{cliccorpora:id6}}\label{\detokenize{cliccorpora:figure-corpora-authorbased}}\end{figure}


\section{Textual subsets}
\label{\detokenize{cliccorpora:textual-subsets}}
\sphinxAtStartPar
The CLiC corpora have been marked up to distinguish between several
textual subsets of novels. The example from \sphinxstyleemphasis{Great Expectations} below
illustrates the subsets:

\begin{sphinxVerbatim}[commandchars=\\\{\}]
\PYGZdq{}And on what evidence, Pip,\PYGZdq{} asked Mr. Jaggers, very coolly, as he
paused with his handkerchief half way to his nose,\PYGZdq{}does Provis make
this claim?\PYGZdq{}

\PYGZdq{}He does not make it,\PYGZdq{} said I, \PYGZdq{}and has never made it, and has no
knowledge or belief that his daughter is in existence.\PYGZdq{}

For once, the powerful pocket\PYGZhy{}handkerchief failed. My reply was so
unexpected that Mr. Jaggers put the handkerchief back into his pocket
without completing the usual performance, folded his arms, and looked
with stern attention at me, though with an immovable face.
\end{sphinxVerbatim}

\sphinxAtStartPar
{[}\sphinxstyleemphasis{Great Expectations}, Chapter 51{]}
\begin{itemize}
\item {} 
\sphinxAtStartPar
\sphinxstylestrong{quotes} \textendash{} any text in quotation marks, i.e. mostly character
speech but also thoughts or songs that might appear in quotes.

\item {} 
\sphinxAtStartPar
\sphinxstylestrong{non\sphinxhyphen{}quotes} \textendash{} narration. A special case of non\sphinxhyphen{}quotes are
\sphinxstylestrong{suspensions}, which represent narratorial interruptions of
character speech that do not end with sentence\sphinxhyphen{}final punctuation.
Suspensions are further divided by length:
\begin{itemize}
\item {} 
\sphinxAtStartPar
\sphinxstylestrong{short suspensions} have a length of up to four words.

\item {} 
\sphinxAtStartPar
\sphinxstylestrong{long suspensions} have a length of five or more words.

\end{itemize}

\end{itemize}

\sphinxAtStartPar
\hyperref[\detokenize{cliccorpora:figure-corpora-markupsubsets}]{Fig.\@ \ref{\detokenize{cliccorpora:figure-corpora-markupsubsets}}} shows how the subsets are marked
up in the chapter views, which can be retrieved from the ‘in bk.’ (in
book) button in concordances (see the {\hyperref[\detokenize{clicanalysis/concordance:concordance}]{\sphinxcrossref{\DUrole{std,std-ref}{Concordance}}}} section) and
the {\hyperref[\detokenize{clicanalysis/texts:texts}]{\sphinxcrossref{\DUrole{std,std-ref}{Texts}}}} tab. The “in book” view also contains a legend of the
markup layers (see \hyperref[\detokenize{cliccorpora:figure-corpora-highlight-subsets}]{Fig.\@ \ref{\detokenize{cliccorpora:figure-corpora-highlight-subsets}}}), which
you can individually select and deselect.

\begin{figure}[htbp]
\centering
\capstart

\noindent\sphinxincludegraphics{{figure-corpora-markupsubsets}.png}
\caption{Chapter view of the example above, exemplifying the mark\sphinxhyphen{}up of
subsets}\label{\detokenize{cliccorpora:id7}}\label{\detokenize{cliccorpora:figure-corpora-markupsubsets}}\end{figure}

\begin{figure}[htbp]
\centering
\capstart

\noindent\sphinxincludegraphics{{figure-corpora-highlight-subsets}.png}
\caption{Example of highlighting subsets}\label{\detokenize{cliccorpora:id8}}\label{\detokenize{cliccorpora:figure-corpora-highlight-subsets}}\end{figure}

\sphinxAtStartPar
The rationale behind the division of the subsets can be found in
\sphinxcite{references:mahlberg-et-al-2016}. The tagging procedure for the most recently
added corpora differs in technical implementation \textendash{} see
\sphinxcode{\sphinxupquote{clic.region}} for details.


\section{Versioning and reproducibility}
\label{\detokenize{cliccorpora:versioning-and-reproducibility}}
\sphinxAtStartPar
The procedure for retrieving, cleaning and importing texts is described in detail in the \sphinxcite{references:github-corpora} repository. Every
change to the repository is marked with a commit number (a sequence of
characters and numbers), as shown in
\hyperref[\detokenize{cliccorpora:figure-version-numbers-github}]{Fig.\@ \ref{\detokenize{cliccorpora:figure-version-numbers-github}}}.

\sphinxAtStartPar
From CLiC 2.0.0 onwards, the corpus commit number is also displayed in
the CLiC interface (see \hyperref[\detokenize{cliccorpora:figure-version-numbers}]{Fig.\@ \ref{\detokenize{cliccorpora:figure-version-numbers}}}) to indicate
which version of the corpora is being used. \sphinxstylestrong{We recommend that users
record both the CLiC version number} (displayed under the CLiC logo,
circled in red in \hyperref[\detokenize{cliccorpora:figure-version-numbers}]{Fig.\@ \ref{\detokenize{cliccorpora:figure-version-numbers}}}) \sphinxstylestrong{and the corpus
version} (circled in green) \sphinxstylestrong{when saving results.} This ensures that
any later differences in search results can be traced to specific
changes in the interface or the corpora.

\begin{figure}[htbp]
\centering
\capstart

\noindent\sphinxincludegraphics{{figure-version-numbers-github}.png}
\caption{The commit version number in the GitHub corpora repository}\label{\detokenize{cliccorpora:id9}}\label{\detokenize{cliccorpora:figure-version-numbers-github}}\end{figure}

\begin{figure}[htbp]
\centering
\capstart

\noindent\sphinxincludegraphics{{figure-version-numbers}.png}
\caption{The CLiC release and corpora version numbers in the CLiC interface}\label{\detokenize{cliccorpora:id10}}\label{\detokenize{cliccorpora:figure-version-numbers}}\end{figure}

\sphinxAtStartPar
The corpora repository also contains the full text of the corpus files
after any manual cleaning changes. The repository history lists all
changes and links to the original versions as downloaded from Project
Gutenberg for the ArTs corpus \sphinxcite{references:github-corpora-initial-arts}
(note that ArTs used to be called “Other” in the repository history) and the
ChiLit corpus \sphinxcite{references:github-corpora-initial-chilit}.


\section{How to cite corpora}
\label{\detokenize{cliccorpora:how-to-cite-corpora}}
\sphinxAtStartPar
When reporting results from CLiC, we recommend citing both the CLiC
version and the corpus version (see {\hyperref[\detokenize{cliccorpora:versioning-and-reproducibility}]{\sphinxcrossref{\DUrole{std,std-ref}{Versioning and reproducibility}}}} for where to find these). A citation
might look like this:
\begin{quote}

\sphinxAtStartPar
We searched for X in {[}corpus/corpora{]} using the Concordance function
in CLiC 2.3 (Mahlberg et al. 2020; corpus version {[}number{]}),
accessed {[}date{]}.
\end{quote}

\sphinxAtStartPar
For the general CLiC citation, see the {\hyperref[\detokenize{introduction:citing-clic}]{\sphinxcrossref{\DUrole{std,std-ref}{Introduction}}}}.

\sphinxstepscope


\chapter{Concordance}
\label{\detokenize{clicanalysis/concordance:concordance}}\label{\detokenize{clicanalysis/concordance:id1}}\label{\detokenize{clicanalysis/concordance::doc}}
\sphinxAtStartPar
Clicking onto the \sphinxstylestrong{‘Concordance’} tab will take you to the concordance
view. In order to create a concordance, you will need to select a corpus
to search in (see {\hyperref[\detokenize{cliccorpora:the-clic-corpora}]{\sphinxcrossref{\DUrole{std,std-ref}{The CLiC corpora}}}}).


\section{Search the corpora}
\label{\detokenize{clicanalysis/concordance:search-the-corpora}}
\sphinxAtStartPar
This is where you select a corpus to search in. The
selection is very flexible and lets you pick a pre\sphinxhyphen{}defined corpus (see {\hyperref[\detokenize{cliccorpora:the-clic-corpora}]{\sphinxcrossref{\DUrole{std,std-ref}{The CLiC corpora}}}})
or choose your own subcorpus \textendash{} with any of the books available in CLiC.


\section{Only in subsets}
\label{\detokenize{clicanalysis/concordance:only-in-subsets}}
\sphinxAtStartPar
Here you can decide whether you want to search through ‘all text’ \textendash{} the
whole book(s) \textendash{} or just one of the subsets: ‘short suspensions’, ‘long
suspensions’, ‘quotes’ and ‘non\sphinxhyphen{}quotes’ (see {\hyperref[\detokenize{cliccorpora:the-clic-corpora}]{\sphinxcrossref{\DUrole{std,std-ref}{The CLiC corpora}}}}).


\section{Search for terms}
\label{\detokenize{clicanalysis/concordance:search-for-terms}}
\sphinxAtStartPar
This is the fundamental parameter of the concordance search \textendash{} it lets
you determine the node word or phrase that forms the basis of the
concordance.

\sphinxAtStartPar
The tokenisation from CLiC 2.0 onwards is based on unicode standard rules
(i.e. Unicode word boundaries implemented with the \sphinxcite{references:icu} library), used
both for queries and importing books.

\sphinxAtStartPar
We consider a boundary mark to be a word\sphinxhyphen{}boundary if…
\begin{itemize}
\item {} 
\sphinxAtStartPar
The \sphinxcite{references:icu} library describes it as at the end of a word, e.g. \sphinxcode{\sphinxupquote{jump}} or number, e.g. \sphinxcode{\sphinxupquote{32.3}}.

\item {} 
\sphinxAtStartPar
It is a single hyphen character surrounded by alpha\sphinxhyphen{}numeric characters.

\item {} 
\sphinxAtStartPar
It is an apostrophe preceded with \sphinxcode{\sphinxupquote{s}}, e.g. \sphinxcode{\sphinxupquote{3 days\textquotesingle{} work}}.

\item {} 
\sphinxAtStartPar
It is one of a whitelist of words preceded with an apostrophe, e.g. \sphinxcode{\sphinxupquote{\textquotesingle{}tis}}.

\end{itemize}

\sphinxAtStartPar
CLiC 2.0 and onwards supports \sphinxstylestrong{wildcards}:
\begin{itemize}
\item {} 
\sphinxAtStartPar
\sphinxcode{\sphinxupquote{*}} means “zero or more characters”, for example:

\sphinxAtStartPar
Placing * at the end of a the sequence \sphinxcode{\sphinxupquote{can}} serves as a placeholder for
any sequence of characters (or zero) and therefore retrieves all instances of
words starting with this sequence, including \sphinxcode{\sphinxupquote{can}}, \sphinxcode{\sphinxupquote{cannot}} and \sphinxcode{\sphinxupquote{can\textquotesingle{}t}}
(but also \sphinxcode{\sphinxupquote{candle}}, \sphinxcode{\sphinxupquote{candles}}, \sphinxcode{\sphinxupquote{candlestick}} etc.)

\sphinxAtStartPar
\sphinxcode{\sphinxupquote{*}} in \sphinxcode{\sphinxupquote{with * hands}} serves as a placeholder for any word token
between \sphinxcode{\sphinxupquote{with}} and \sphinxcode{\sphinxupquote{hands}}, retrieving sequences like \sphinxcode{\sphinxupquote{with her hands}},
\sphinxcode{\sphinxupquote{with his hands}}, \sphinxcode{\sphinxupquote{with their hands}}, \sphinxcode{\sphinxupquote{with both hands}},
\sphinxcode{\sphinxupquote{with clean hands}} etc.

\item {} 
\sphinxAtStartPar
\sphinxcode{\sphinxupquote{?}} means “one”

\end{itemize}

\sphinxAtStartPar
The search will only retrieve valid tokens according to the rules above.
This means that the search will ignore punctuation in your search query,
and searching for a punctuation sign alone will not retrieve any results.
If your research focuses on punctuation markers you can evade this issue
by using the filter function in the subset tab: Go to the subset tab,
select the relevant subset, for example non\sphinxhyphen{}quotes, and filter the rows
to the punctuation marker of interest.

\sphinxAtStartPar
Two hyphens separate words: for example, \sphinxstyleemphasis{Char\textendash{}lotte} in Oliver
Twist (OT.c6.p20) “Oliver’s gone mad! Char\textendash{}lotte!” counts as two
tokens.

\sphinxAtStartPar
For the detailed technical documentation and more examples see \sphinxcode{\sphinxupquote{clic.tokenizer}}.


\section{‘Whole phrase’ or ‘Any word’}
\label{\detokenize{clicanalysis/concordance:whole-phrase-or-any-word}}
\sphinxAtStartPar
When you have entered several terms, you need to specify whether it is
to be searched as one phrase (equivalent to using double quotes in a
search engine, e.g. \sphinxstyleemphasis{dense fog}) or any of the words individually
(\sphinxstyleemphasis{dense} and \sphinxstyleemphasis{fog}).


\section{Co\sphinxhyphen{}text}
\label{\detokenize{clicanalysis/concordance:co-text}}
\sphinxAtStartPar
The maximum number of words in the co\sphinxhyphen{}text is set at 10 on either side
in a concordance (depending on the length of the words and the size of
the screen you might see fewer). You can see the full chapter view by
clicking on \sphinxstylestrong{‘in bk.’ (in book) button} at the end of any row (see
\hyperref[\detokenize{clicanalysis/concordance:figure-analysis-concordance-cotext-inbookbutton}]{Fig.\@ \ref{\detokenize{clicanalysis/concordance:figure-analysis-concordance-cotext-inbookbutton}}} and \hyperref[\detokenize{clicanalysis/concordance:figure-analysis-concordance-cotext-inbookchapter}]{Fig.\@ \ref{\detokenize{clicanalysis/concordance:figure-analysis-concordance-cotext-inbookchapter}}}).

\begin{figure}[htbp]
\centering
\capstart

\noindent\sphinxincludegraphics{{figure-analysis-concordance-cotext-inbookbutton}.png}
\caption{The ‘in bk.’ (in book) button leads to the book view
of the occurrence}\label{\detokenize{clicanalysis/concordance:id8}}\label{\detokenize{clicanalysis/concordance:figure-analysis-concordance-cotext-inbookbutton}}\end{figure}

\begin{figure}[htbp]
\centering
\capstart

\noindent\sphinxincludegraphics{{figure-analysis-concordance-cotext-inbookchapter}.png}
\caption{The ‘in bk.’ view shows the relevant section in the whole book.}\label{\detokenize{clicanalysis/concordance:id9}}\label{\detokenize{clicanalysis/concordance:figure-analysis-concordance-cotext-inbookchapter}}
\begin{sphinxlegend}
\sphinxAtStartPar
The preface shown here is a very short chapter. (Note that all authorial
text occurring before the official first chapter, is counted as
‘chapter 0’ in CLiC). This preface contains no quotes or suspensions;
compare to the subset markup in the chapter view shown above.
\end{sphinxlegend}
\end{figure}


\section{Results}
\label{\detokenize{clicanalysis/concordance:results}}
\sphinxAtStartPar
These options allow you to adjust the way the concordance output is
displayed.


\subsection{Filter rows}
\label{\detokenize{clicanalysis/concordance:filter-rows}}
\sphinxAtStartPar
This filter option lets you filter the concordance output by the rows
that contain a particular sequence of letters (both in the node and
co\sphinxhyphen{}text). For example, searching for \sphinxcode{\sphinxupquote{hands}} in \sphinxstyleemphasis{Oliver Twist} yields 124
results; when we use the option \sphinxstylestrong{‘filter rows’} and search for
\sphinxcode{\sphinxupquote{pockets}}, this is filtered down to 8 results as illustrated in \hyperref[\detokenize{clicanalysis/concordance:figure-analysis-concordance-results-filter}]{Fig.\@ \ref{\detokenize{clicanalysis/concordance:figure-analysis-concordance-results-filter}}}.

\begin{figure}[htbp]
\centering
\capstart

\noindent\sphinxincludegraphics{{figure-analysis-concordance-results-filter}.png}
\caption{Concordance of \sphinxcode{\sphinxupquote{hands}} in \sphinxstyleemphasis{Oliver Twist} filtered down to
\sphinxcode{\sphinxupquote{pockets}} in the co\sphinxhyphen{}text}\label{\detokenize{clicanalysis/concordance:id10}}\label{\detokenize{clicanalysis/concordance:figure-analysis-concordance-results-filter}}\end{figure}

\sphinxAtStartPar
Note that the filter, when searching for character sequences does not
necessarily search for complete words: for example, filtering a
concordance of \sphinxcode{\sphinxupquote{head}} in \sphinxstyleemphasis{Oliver Twist} for \sphinxcode{\sphinxupquote{eat}} yields both
occurrences of the verb \sphinxcode{\sphinxupquote{eat}}, and the instance \sphinxcode{\sphinxupquote{threatened}}, which
contains the same sequence of letters (see \hyperref[\detokenize{clicanalysis/concordance:figure-analysis-concordance-results-filtersequence}]{Fig.\@ \ref{\detokenize{clicanalysis/concordance:figure-analysis-concordance-results-filtersequence}}}).

\sphinxAtStartPar
The filter function is cruder than the KWICGrouper; it can be usefully
applied to filter down a large set of results before you do a more
fine\sphinxhyphen{}grained categorisation. You might want to filter down the results
to rows containing similar word forms. For example, filtering for \sphinxcode{\sphinxupquote{girl}}
will also retrieve rows containing \sphinxcode{\sphinxupquote{girlish}} and \sphinxcode{\sphinxupquote{girls}}. Moreover,
unlike the main concordance search and the KWICGrouper, the filter lets
you search for particular types of punctuation (e.g. round brackets used
in suspensions).

\begin{figure}[htbp]
\centering
\capstart

\noindent\sphinxincludegraphics{{figure-analysis-concordance-results-filtersequence}.png}
\caption{Filtering for the letter sequence \sphinxcode{\sphinxupquote{eat}} returns forms of
the verb \sphinxcode{\sphinxupquote{eat}} and other words containing the sequence}\label{\detokenize{clicanalysis/concordance:id11}}\label{\detokenize{clicanalysis/concordance:figure-analysis-concordance-results-filtersequence}}\end{figure}


\subsection{View as}
\label{\detokenize{clicanalysis/concordance:view-as}}
\sphinxAtStartPar
From CLiC 2.0 onwards there are three options to view the concordance results:
\begin{enumerate}
\sphinxsetlistlabels{\arabic}{enumi}{enumii}{}{.}%
\item {} 
\sphinxAtStartPar
Basic results: concordance lines + book short title; link to “in bk.” view

\item {} 
\sphinxAtStartPar
Full metadata: concordance lines + book short title; chapter, paragraph \&
sentence numbers; link to “in bk.” view

\item {} 
\sphinxAtStartPar
Distribution plot: overview of matching lines per book

\end{enumerate}

\sphinxAtStartPar
The default view is 1. and 2. gives more information on the same view. View 3.
is completely different: it does not show the text in concordance lines but plots
the distribution of matching concordance lines across the searched books. Note that
if a book in the searched corpus has zero matches it will not be shown in
the distribution plot (hence, \hyperref[\detokenize{clicanalysis/concordance:figure-distribution-plot-workhouse}]{Fig.\@ \ref{\detokenize{clicanalysis/concordance:figure-distribution-plot-workhouse}}} only shows
10 out of the 15 books in the DNov corpus).

\begin{figure}[htbp]
\centering
\capstart

\noindent\sphinxincludegraphics{{figure_distribution_plot_workhouse}.png}
\caption{Distribution plot of \sphinxtitleref{workhouse} in DNov}\label{\detokenize{clicanalysis/concordance:id12}}\label{\detokenize{clicanalysis/concordance:figure-distribution-plot-workhouse}}\end{figure}

\sphinxAtStartPar
Like the concordance view, the distribution plot view can show
“KWICGrouped lines” i.e. lines that contain particular patterns in the
proximity of the nodes (see the {\hyperref[\detokenize{clicanalysis/subsets:kwicgrouper}]{\sphinxcrossref{\DUrole{std,std-ref}{KWICGrouper}}}} section below for a
detailed example). Once the KWICGrouper is activated by selecting word
types in KWICGrouper box in the menu bar, the corresponding instances
in the distribution plot are coloured.
\hyperref[\detokenize{clicanalysis/concordance:figure-distribution-plot-workhouse-kwicgroup}]{Fig.\@ \ref{\detokenize{clicanalysis/concordance:figure-distribution-plot-workhouse-kwicgroup}}} shows coloured
lines containing references to children (\sphinxtitleref{boy}, \sphinxtitleref{chance\sphinxhyphen{}child},
\sphinxtitleref{charity\sphinxhyphen{}boy}, \sphinxtitleref{child}, \sphinxtitleref{children}, \sphinxtitleref{girl} and \sphinxtitleref{orphan}). Light green
represents one match in the concordance line, a darker green two and
purple three matches. (Not displayed here: a line with four matches
would be shown in pink.)

\begin{figure}[htbp]
\centering
\capstart

\noindent\sphinxincludegraphics{{figure_distribution_plot_workhouse_KWICGroup}.png}
\caption{Distribution plot of \sphinxtitleref{workhouse} in DNov, KWICGrouped for references to children}\label{\detokenize{clicanalysis/concordance:id13}}\label{\detokenize{clicanalysis/concordance:figure-distribution-plot-workhouse-kwicgroup}}\end{figure}


\subsection{Basic sorting}
\label{\detokenize{clicanalysis/concordance:basic-sorting}}
\sphinxAtStartPar
The concordance lines can be sorted by any of the columns in the
concordance by clicking on the header, which will then be marked with
dark arrows. For example, by clicking on \sphinxstylestrong{‘Left’} the lines will be
sorted by the first word to the left of the node and by clicking on
\sphinxstylestrong{‘Right’} by the first word on the right. If you have the metadata
columns activated you can also sort by these, for example to sort all
entries by chapter. Similarly, if you have created your own tags (see
the section {\hyperref[\detokenize{clicanalysis/subsets:manage-tag-columns}]{\sphinxcrossref{\DUrole{std,std-ref}{Manage tag columns}}}} below), you can sort for lines with a
particular tag. Clicking on the same header a second time will reverse
the order of sorting.

\sphinxAtStartPar
Note that you can create a \sphinxstylestrong{“sorting sequence”} by clicking on various
headers while pressing the \sphinxstylestrong{shift key}. For example, you could sort a
concordance first by the words on the right and then by book, as
illustrated in \hyperref[\detokenize{clicanalysis/concordance:figure-analysis-concordance-sorting-fireplacecombined}]{Fig.\@ \ref{\detokenize{clicanalysis/concordance:figure-analysis-concordance-sorting-fireplacecombined}}},
which shows a concordance of \sphinxstyleemphasis{fireplace} sorted first by book \textendash{} so that
results from \sphinxstyleemphasis{Barnaby Rudge (BR)} come first \textendash{} and then ordered by the
co\sphinxhyphen{}text on the right.

\begin{figure}[htbp]
\centering
\capstart

\noindent\sphinxincludegraphics{{figure-analysis-concordance-sorting-fireplacecombined}.png}
\caption{Concordance of fireplace in DNov (Dickens’s Novels) \textendash{}
first ordered by book, then by the first word on the right}\label{\detokenize{clicanalysis/concordance:id14}}\label{\detokenize{clicanalysis/concordance:figure-analysis-concordance-sorting-fireplacecombined}}\end{figure}


\subsection{KWICGrouper}
\label{\detokenize{clicanalysis/concordance:kwicgrouper}}\label{\detokenize{clicanalysis/concordance:id4}}
\sphinxAtStartPar
The KWICGrouper is a tool that allows you to quickly group the
concordance lines according to patterns that you find as you go through
the concordance. For a basic introduction to the KWICGrouper
functionality (in the CLiC 1.5 interface) you can watch our KWICGrouper
video tutorial from May 2017 \sphinxcite{references:clic-kwicgrouper-video}.

\sphinxAtStartPar
The idea of the KWICGrouper is that you look for patterns as you search
for particular words. Any matching lines will be highlighted and moved
to the top of the screen. Among the matching lines we further
distinguish between the lines based on how many matches they contain. A
line with one match is highlighted in light green, lines with two
matches are coloured in a darker green, those with three in purple and,
finally, those with four in pink. (For lines with more matches than
these, the colours will repeat.) The KWICGrouper gives you two options:
\begin{itemize}
\item {} 
\sphinxAtStartPar
\sphinxstylestrong{‘Search in span’:} Set the span for the KWICGrouper search. By
dragging the slider you can adjust the number of words that will be
searched to the left and right of the search term. The maximum (and
default) span is 5 positions to either side.

\item {} 
\sphinxAtStartPar
\sphinxstylestrong{‘Search for types’:} Choose one or more words to search for in the
span. This is currently limited to single words, but there is no
limit on how many words you add.

\end{itemize}

\sphinxAtStartPar
The total number of matching rows will be displayed at the top; the
process is illustrated in \hyperref[\detokenize{clicanalysis/concordance:figure-analysis-concordance-kwicgrouper-fireplain}]{Fig.\@ \ref{\detokenize{clicanalysis/concordance:figure-analysis-concordance-kwicgrouper-fireplain}}}
and \hyperref[\detokenize{clicanalysis/concordance:figure-analysis-concordance-kwicgrouper-firetypes}]{Fig.\@ \ref{\detokenize{clicanalysis/concordance:figure-analysis-concordance-kwicgrouper-firetypes}}}, which
show 1. the plain concordance lines as returned when searching for
\sphinxstyleemphasis{fire} in Dickens’s novels and 2. the process of choosing types (forms
of words) from co\sphinxhyphen{}text surrounding \sphinxstyleemphasis{fire} in the concordance in order
to group the concordance lines.

\begin{figure}[htbp]
\centering
\capstart

\noindent\sphinxincludegraphics{{figure-analysis-concordance-kwicgrouper-fireplain}.png}
\caption{The first concordance lines of fire in DNov (Dickens’s
Novels) with the default sorting by ‘in bk’}\label{\detokenize{clicanalysis/concordance:id15}}\label{\detokenize{clicanalysis/concordance:figure-analysis-concordance-kwicgrouper-fireplain}}\end{figure}

\begin{figure}[htbp]
\centering
\capstart

\noindent\sphinxincludegraphics{{figure-analysis-concordance-kwicgrouper-firetypes}.png}
\caption{Selecting types related to sitting from the KWICGrouper
to group the concordance lines}\label{\detokenize{clicanalysis/concordance:id16}}\label{\detokenize{clicanalysis/concordance:figure-analysis-concordance-kwicgrouper-firetypes}}\end{figure}

\sphinxAtStartPar
The dropdown only contains those word forms that actually appear
around the node term in the specified search span. Therefore, while
\sphinxstyleemphasis{sitiwation} is listed here, it wouldn’t be listed if we had searched
for another node term or used other books; it only appears once in this
set in the following example context:
\begin{quote}

\sphinxAtStartPar
I don’t take no pride out on it, Sammy,’ replied Mr. Weller, poking
the fire vehemently, ‘it’s a horrid \sphinxstylestrong{sitiwation}. I’m actiwally
drove out o’ house and home by it.The breath was scarcely out o’ your
poor mother\sphinxhyphen{}in\sphinxhyphen{}law’s body, ven vun old ‘ooman sends me a pot o’ jam,
and another a pot o’ jelly, and another brews a blessed large jug o’
camomile\sphinxhyphen{}tea, vich she brings in vith her own hands.’

\sphinxAtStartPar
\sphinxstyleemphasis{{[}Pickwick Papers, Chapter LI.{]}}
\end{quote}

\begin{figure}[htbp]
\centering
\capstart

\noindent\sphinxincludegraphics{{figure-analysis-concordance-kwicgrouper-fireresults}.png}
\caption{The resulting ‘KWICGrouped’ concordance lines: the
selected types are listed in the search box on the right; and in the
case of this example it is suitable to restrict the search span to
only the left side of the node}\label{\detokenize{clicanalysis/concordance:id17}}\label{\detokenize{clicanalysis/concordance:figure-analysis-concordance-kwicgrouper-fireresults}}\end{figure}

\sphinxAtStartPar
The KWICGrouper only searches through a number of words to the left and
right of the node term, as specified by the search span.
\hyperref[\detokenize{clicanalysis/concordance:figure-analysis-concordance-kwicgrouper-fireresults}]{Fig.\@ \ref{\detokenize{clicanalysis/concordance:figure-analysis-concordance-kwicgrouper-fireresults}}} shows
the resulting concordance lines according to the KWICGrouper settings
after manually choosing types related to the action of sitting. Apart
from the selected search types the search span has also been restricted
to the left side so that clearer patterns of sitting by the fire become
visible.

\begin{figure}[htbp]
\centering
\capstart

\noindent\sphinxincludegraphics{{figure-analysis-concordance-kwicgrouper-fireback}.png}
\caption{The first lines of fire co\sphinxhyphen{}occurring with back (i.e. one
KWICGrouper match) are highlighted and moved to the top}\label{\detokenize{clicanalysis/concordance:id18}}\label{\detokenize{clicanalysis/concordance:figure-analysis-concordance-kwicgrouper-fireback}}\end{figure}

\sphinxAtStartPar
Apart from looking for characters sitting by the fire, it might also be
of interest to look for characters standing by the fire. We have shown
in our previous work (see Chapter 6 of \sphinxcite{references:mahlberg-2013}) that the cluster
\sphinxstyleemphasis{with his back to the fire} is prominent in Dickens’s and 19th century
novels by other writers.
\hyperref[\detokenize{clicanalysis/concordance:figure-analysis-concordance-kwicgrouper-fireback}]{Fig.\@ \ref{\detokenize{clicanalysis/concordance:figure-analysis-concordance-kwicgrouper-fireback}}} shows the
first concordance lines of \sphinxstyleemphasis{fire} with \sphinxstyleemphasis{back} on the left (sorted to the
left).

\sphinxAtStartPar
The output from the KWICGrouper lists at the top of the screen the
number of lines that contain any number of matches. In this case
there are only lines with one match, but no lines with more than one
match. The message “36 entries with 1 KWIC match” means that 36 lines
contain both \sphinxstyleemphasis{fire} and \sphinxstyleemphasis{back}. This function becomes useful when we now
look for gendered pronouns. As shown in
\hyperref[\detokenize{clicanalysis/concordance:figure-analysis-concordance-kwicgrouper-firebackhis}]{Fig.\@ \ref{\detokenize{clicanalysis/concordance:figure-analysis-concordance-kwicgrouper-firebackhis}}}, there
are 27 lines in which \sphinxstyleemphasis{fire} co\sphinxhyphen{}occurs with both \sphinxstyleemphasis{back} and \sphinxstyleemphasis{his}. Most
of these occurrences appear in the pattern \sphinxstyleemphasis{with his back to the fire},
as becomes obvious when we reverse the sorting on the left so that
\sphinxstyleemphasis{his} occurs at the top in the first position to the left of \sphinxstyleemphasis{fire} \textendash{} the L1
position. On the other hand, as we can see from
\hyperref[\detokenize{clicanalysis/concordance:figure-analysis-concordance-kwicgrouper-firebackher}]{Fig.\@ \ref{\detokenize{clicanalysis/concordance:figure-analysis-concordance-kwicgrouper-firebackher}}},
Dickens’s novels contain only one instance of \sphinxstyleemphasis{fire} co\sphinxhyphen{}occurring with
\sphinxstyleemphasis{back} and \sphinxstyleemphasis{her} (with \sphinxstyleemphasis{her back to the fire}).

\begin{figure}[htbp]
\centering
\capstart

\noindent\sphinxincludegraphics{{figure-analysis-concordance-kwicgrouper-firebackhis}.png}
\caption{The 27 lines with two matches (here, back and his) are
highlighted in a darker green}\label{\detokenize{clicanalysis/concordance:id19}}\label{\detokenize{clicanalysis/concordance:figure-analysis-concordance-kwicgrouper-firebackhis}}\end{figure}

\begin{figure}[htbp]
\centering
\capstart

\noindent\sphinxincludegraphics{{figure-analysis-concordance-kwicgrouper-firebackher}.png}
\caption{Only one line contains both back and her; it is
highlighted and shown above single match lines}\label{\detokenize{clicanalysis/concordance:id20}}\label{\detokenize{clicanalysis/concordance:figure-analysis-concordance-kwicgrouper-firebackher}}\end{figure}


\subsection{Manage tag columns}
\label{\detokenize{clicanalysis/concordance:manage-tag-columns}}\label{\detokenize{clicanalysis/concordance:id7}}
\sphinxAtStartPar
Once you have identified lines with patterns of interest, you might want
to place these into one or more categories. CLiC provides a flexible
tagging system for this.
\hyperref[\detokenize{clicanalysis/concordance:figure-analysis-concordance-tagcolumns-dream}]{Fig.\@ \ref{\detokenize{clicanalysis/concordance:figure-analysis-concordance-tagcolumns-dream}}} illustrates
what a tagged concordance can look like. The tags are
user\sphinxhyphen{}defined so you can create tags that are relevant to your project.
In this case, occurrences of \sphinxstyleemphasis{dream} in \sphinxstyleemphasis{Oliver Twist} have been tagged
according to who is dreaming.

\begin{figure}[htbp]
\centering
\capstart

\noindent\sphinxincludegraphics{{figure-analysis-concordance-tagcolumns-dream}.png}
\caption{Tagged concordance lines of dream in Oliver Twist}\label{\detokenize{clicanalysis/concordance:id21}}\label{\detokenize{clicanalysis/concordance:figure-analysis-concordance-tagcolumns-dream}}\end{figure}

\sphinxAtStartPar
In order to tag the lines, click on \sphinxstylestrong{‘manage tag columns’} (shown in
the bottom right corner of \hyperref[\detokenize{clicanalysis/concordance:figure-analysis-concordance-tagcolumns-dream}]{Fig.\@ \ref{\detokenize{clicanalysis/concordance:figure-analysis-concordance-tagcolumns-dream}}}) and create
your own tag(s) through the \sphinxstylestrong{‘Add new’} option (see
\hyperref[\detokenize{clicanalysis/concordance:figure-analysis-concordance-tagcolumns-menu}]{Fig.\@ \ref{\detokenize{clicanalysis/concordance:figure-analysis-concordance-tagcolumns-menu}}}). You can
rename a tag by selecting it from the \sphinxstylestrong{‘Tag columns’} list and
renaming it in the text box. Once you have created your tag(s), you can
click \sphinxstylestrong{‘Back’} to return to the menu. Now you can select the relevant
concordance lines by clicking on them and you will see that the sidebar
contains the list of your tags. Once one or more lines are selected you
can click the tick next to the relevant tag in order to tag the line
(see \hyperref[\detokenize{clicanalysis/concordance:figure-analysis-concordance-tagcolumns-selectline}]{Fig.\@ \ref{\detokenize{clicanalysis/concordance:figure-analysis-concordance-tagcolumns-selectline}}}).
An extra column will appear for each tag and you can sort on these
columns as mentioned in the {\hyperref[\detokenize{clicanalysis/concordance:basic-sorting}]{\sphinxcrossref{\DUrole{std,std-ref}{Basic sorting}}}} section above. Selected and tagged
rows will be automatically deselected when you click on (i.e. select) a
new row.

\begin{figure}[htbp]
\centering
\capstart

\noindent\sphinxincludegraphics{{figure-analysis-concordance-tagcolumns-menu}.png}
\caption{The menu for adding and renaming tags}\label{\detokenize{clicanalysis/concordance:id22}}\label{\detokenize{clicanalysis/concordance:figure-analysis-concordance-tagcolumns-menu}}\end{figure}

\begin{figure}[htbp]
\centering
\capstart

\noindent\sphinxincludegraphics{{figure-analysis-concordance-tagcolumns-selectline}.png}
\caption{Select a line (by clicking on it) in order to apply an
existing tag; once tagged, the tick in the sidebar will appear green
for the selected line.
A tick will also be added to the tag column in
the concordance itself.}\label{\detokenize{clicanalysis/concordance:id23}}\label{\detokenize{clicanalysis/concordance:figure-analysis-concordance-tagcolumns-selectline}}\end{figure}

\sphinxstepscope


\chapter{FlexiConc}
\label{\detokenize{clicanalysis/flexiconc:flexiconc}}\label{\detokenize{clicanalysis/flexiconc:id1}}\label{\detokenize{clicanalysis/flexiconc::doc}}
\sphinxAtStartPar
\sphinxstylestrong{‘FlexiConc’} is a Python library developed to support corpus linguists
by supporting the analysis of concordances. CLiC incorporates FlexiConc
functionality in a dedicated analysis tab.


\section{Starting FlexiConc}
\label{\detokenize{clicanalysis/flexiconc:starting-flexiconc}}
\sphinxAtStartPar
After selecting FlexiConc in the sidebar, the initial query window looks
identical to the one in the Concordance tab, with some additional buttons
underneath.

\begin{figure}[htbp]
\centering
\capstart

\noindent\sphinxincludegraphics{{flexiclic_query}.png}
\caption{Query in FlexiConc mode in CLiC}\label{\detokenize{clicanalysis/flexiconc:id2}}\label{\detokenize{clicanalysis/flexiconc:flexiclic-query}}\end{figure}

\sphinxAtStartPar
The query options (corpus selection, subsets, search terms, wildcards,
whole phrase/any word) work the same as in the {\hyperref[\detokenize{clicanalysis/concordance:concordance}]{\sphinxcrossref{\DUrole{std,std-ref}{Concordance}}}} tab.


\subsection{Search results}
\label{\detokenize{clicanalysis/flexiconc:search-results}}
\sphinxAtStartPar
In the screenshot above, we are searching the 19C corpus of 19th century
English novels. We focus on non\sphinxhyphen{}quotes \textendash{} all parts of the text that are
not part of a character’s direct speech. Within this subset, the search
terms are cheeks, cheek, neck, fingers, and ear. These are examples of
body\sphinxhyphen{}part nouns that appear in the mid\sphinxhyphen{}level frequency range in this
corpus. Thus, neither of them is extremely common or infrequent.
Searching for any word ensures that all terms are searched separately,
rather than searching for the rather nonsensical string “cheeks cheek
neck fingers ear” as a sequence. The resulting concordance view looks
like this:

\begin{figure}[htbp]
\centering
\capstart

\noindent\sphinxincludegraphics{{flexiclic_first_result}.png}
\caption{Initial concordance in FlexiConc mode}\label{\detokenize{clicanalysis/flexiconc:id3}}\label{\detokenize{clicanalysis/flexiconc:flexiclic-first-result}}\end{figure}


\section{Concordancing strategies}
\label{\detokenize{clicanalysis/flexiconc:concordancing-strategies}}
\sphinxAtStartPar
FlexiConc takes the query result as input and allows you to perform
different steps, which are operationalized as \sphinxstyleemphasis{algorithms}.

\sphinxAtStartPar
Each FlexiConc algorithm performs an operation that belongs to one of
three central categories:
\begin{enumerate}
\sphinxsetlistlabels{\arabic}{enumi}{enumii}{}{.}%
\item {} 
\sphinxAtStartPar
\sphinxstylestrong{Selecting} \textendash{} focus on specific subsets of concordance lines based
on a variety of criteria, including metadata categories and contextual
keywords.

\item {} 
\sphinxAtStartPar
\sphinxstylestrong{Ordering} \textendash{} arrange concordance lines by sorting or ranking them,
using numeric preference scores to prioritize those of interest.

\item {} 
\sphinxAtStartPar
\sphinxstylestrong{Grouping} \textendash{} organize lines into groups by applying explicit
partitioning criteria or through clustering based on similarity
measures.

\end{enumerate}


\section{Analysis tree}
\label{\detokenize{clicanalysis/flexiconc:analysis-tree}}
\sphinxAtStartPar
FlexiConc organizes the concordancing process in an analysis tree.
Algorithms can be applied sequentially in a hierarchical structure,
meaning that you can ‘branch off’ the analysis on any level. For
instance, you can apply a sorting algorithm to a subset, which results
in only that subset being sorted. Alternatively, you can apply the same
sorting algorithm to the concordance lines that you obtained the subset
from, which would lead to all lines within that view being sorted.

\sphinxAtStartPar
By default, your analysis has one branch, which is started when you run
the query, and is represented by the button with the number 1 next to
the tree symbol. Clicking on the tree symbol takes you to an overview of
the entire analysis tree:

\begin{figure}[htbp]
\centering
\capstart

\noindent\sphinxincludegraphics{{flexiclic_tree}.png}
\caption{Analysis tree}\label{\detokenize{clicanalysis/flexiconc:id4}}\label{\detokenize{clicanalysis/flexiconc:flexiclic-tree}}\end{figure}

\sphinxAtStartPar
In this example, we started from the concordance shown in the initial
query, which contains 1,575 lines. We applied a \sphinxstylestrong{random sort} \textendash{} which
is an arrangement node (in that it changes the order of lines). From
there, the analysis branched off in two directions:
\begin{enumerate}
\sphinxsetlistlabels{\arabic}{enumi}{enumii}{}{.}%
\item {} 
\sphinxAtStartPar
\sphinxcode{\sphinxupquote{Branch 1}} is a subset based on \sphinxstylestrong{select by a token\sphinxhyphen{}level
attribute}. In this case, the concordance view is limited to lines
containing the token \sphinxstyleemphasis{her} in the L1 position (= \sphinxstyleemphasis{offset \sphinxhyphen{}1}). We
then applied an \sphinxstyleemphasis{arrangement} node, \sphinxstylestrong{flat clustering by embeddings}.
This is a grouping algorithm that sorts concordance lines into
textually similar partitions.

\item {} 
\sphinxAtStartPar
\sphinxcode{\sphinxupquote{Branch 2}} is also a subset of the concordance. Like in branch 1,
it subsets the concordance by \sphinxstylestrong{select by a token\sphinxhyphen{}level attribute}.
However, in branch 1, the attribute was a specific token (\sphinxstyleemphasis{her}). In
branch 2, we instead use the attribute \sphinxstylestrong{pos\_} as a filtering
criterion, which is a part\sphinxhyphen{}of\sphinxhyphen{}speech tag that was added as an
annotation layer to the overall concordance in an \sphinxstyleemphasis{annotation} step.

\end{enumerate}

\sphinxAtStartPar
Like the selection for her in branch 1, this subset selects based on the
L1 position, this time, the selection criterion is that pos\_ has the
value \sphinxstyleemphasis{ADJ}. In other words, we are selecting lines where the node is
preceded by an adjective.

\sphinxAtStartPar
Clicking on the \sphinxstylestrong{numbers at the end of each branch} takes you to the
respective branch, where you can add more algorithms to it.


\subsection{Saving and loading the tree}
\label{\detokenize{clicanalysis/flexiconc:saving-and-loading-the-tree}}
\sphinxAtStartPar
While you are in the tree view, you will also see a \sphinxstylestrong{save and a load
button}, as seen below.

\begin{figure}[htbp]
\centering
\capstart

\noindent\sphinxincludegraphics{{flexiclic_save_tree}.png}
\caption{Save and load buttons for the analysis tree}\label{\detokenize{clicanalysis/flexiconc:id5}}\label{\detokenize{clicanalysis/flexiconc:flexiclic-save-tree}}\end{figure}

\sphinxAtStartPar
Clicking on \sphinxcode{\sphinxupquote{save to file}} initiates a download, where the entire tree
structure is stored as a JSON file.

\sphinxAtStartPar
You can \sphinxcode{\sphinxupquote{load}} this file back into FlexiConc to recreate your analysis
at any point. This also allows you to directly share your analysis with
other researchers or students.


\section{Adding algorithms}
\label{\detokenize{clicanalysis/flexiconc:adding-algorithms}}
\sphinxAtStartPar
When you are on a branch, clicking on the \sphinxcode{\sphinxupquote{add algorithm}} button shows
all available algorithms. You can scroll down or use the search bar to
type an algorithm name.

\begin{figure}[htbp]
\centering
\capstart

\noindent\sphinxincludegraphics{{flexiclic_algo}.png}
\caption{Adding an algorithm to the tree}\label{\detokenize{clicanalysis/flexiconc:id6}}\label{\detokenize{clicanalysis/flexiconc:flexiclic-algo}}\end{figure}
\begin{enumerate}
\sphinxsetlistlabels{\arabic}{enumi}{enumii}{}{.}%
\item {} 
\sphinxAtStartPar
Adding a new algorithm below an existing one will create a new node
on the same branch.

\item {} 
\sphinxAtStartPar
Clicking on the \sphinxcode{\sphinxupquote{plus sign +}} next to the tree symbol creates a new
branch on the top level.

\item {} 
\sphinxAtStartPar
By clicking on the \sphinxcode{\sphinxupquote{branch symbol}} at the bottom of a node, you
create a new branch below that node. In the example below, the new
branch would be initiated containing all steps up to, and including,
the algorithm \sphinxcode{\sphinxupquote{select by token\sphinxhyphen{}level string attribute}}, but not any
steps carried out below.

\end{enumerate}

\begin{figure}[htbp]
\centering
\capstart

\noindent\sphinxincludegraphics[width=0.400\linewidth]{{flexiclic_algo_options}.png}
\caption{Branching off from an existing branch}\label{\detokenize{clicanalysis/flexiconc:id7}}\label{\detokenize{clicanalysis/flexiconc:flexiclic-algo-options}}\end{figure}

\sphinxAtStartPar
As we have seen so far, the branches are numbered by default; and the
number automatically counts up in the order of branch creation. This is
a useful default, as the numbering then serves as a record of the order
in which your analysis gradually built up. However, there are cases
where the numbering isn’t enough to keep track of your results. You can
therefore create a named path for any given branch by using the menu
below the \sphinxcode{\sphinxupquote{Add algorithm}} button:

\begin{figure}[htbp]
\centering
\capstart

\noindent\sphinxincludegraphics{{flexiclic_named_path}.png}
\caption{Naming a path}\label{\detokenize{clicanalysis/flexiconc:id8}}\label{\detokenize{clicanalysis/flexiconc:flexiclic-named-path}}\end{figure}

\sphinxAtStartPar
Creating a named path will not overwrite your numbered path. Instead, it
creates a copy that is stored separately under the name that you chose,
and which you can access through the tree view just like a numbered path.
After creating a named path, you will stay on the numbered path you have
copied.

\sphinxAtStartPar
\sphinxstylestrong{Named paths are immutable}, i.e., you cannot append any algorithms to
them.

\sphinxAtStartPar
Named paths are also displayed in the tree view. Below, you can see that
\sphinxcode{\sphinxupquote{branch 2}} and \sphinxcode{\sphinxupquote{my\_path}} currently share the exact same sequence of
algorithm steps:

\begin{figure}[htbp]
\centering
\capstart

\noindent\sphinxincludegraphics{{flexiclic_named_path_tree}.png}
\caption{Named path in the tree view}\label{\detokenize{clicanalysis/flexiconc:id9}}\label{\detokenize{clicanalysis/flexiconc:flexiclic-named-path-tree}}\end{figure}


\section{Annotations}
\label{\detokenize{clicanalysis/flexiconc:annotations}}
\sphinxAtStartPar
Annotations are automatic analyses of the concordance data that are
added to the concordance through some external source of information.
Once added, you can use annotations in various algorithms. You can find
the annotations in the add annotation toolbar located right below the
query window.

\begin{figure}[htbp]
\centering
\capstart

\noindent\sphinxincludegraphics{{flexiclic_annotations}.png}
\caption{Adding annotation}\label{\detokenize{clicanalysis/flexiconc:id10}}\label{\detokenize{clicanalysis/flexiconc:flexiclic-annotations}}\end{figure}

\sphinxAtStartPar
Currently, two \sphinxstylestrong{types of annotations} are available:

\sphinxAtStartPar
\sphinxstylestrong{Similarity scores} calculate the similarity between concordance
lines. Three distinct types of similarity annotation are available:
sentence transformers, SpaCy embeddings, and TF\sphinxhyphen{}IDF score.

\sphinxAtStartPar
Each of these similarity measures offers a range of settings to specify
how the similarities should be calculated. The figure below shows the
options for \sphinxcode{\sphinxupquote{Annotate with SpaCy Embeddings}}: in particular, you need
to specify the \sphinxstylestrong{model} and, optionally, the \sphinxstylestrong{offset} if you want to
obtain similarities for a fixed context window.

\begin{figure}[htbp]
\centering
\capstart

\noindent\sphinxincludegraphics{{flexiclic_annotations_spacy}.png}
\caption{Options for adding annotation with SpaCy embeddings}\label{\detokenize{clicanalysis/flexiconc:id11}}\label{\detokenize{clicanalysis/flexiconc:flexiclic-annotations-spacy}}\end{figure}

\sphinxAtStartPar
Similarity measures can be used in the Flat clustering by embeddings
algorithm, which creates groups based on the similarity scores by
applying k\sphinxhyphen{}means or agglomerative clustering.

\sphinxAtStartPar
\sphinxstylestrong{Part\sphinxhyphen{}of\sphinxhyphen{}speech (POS) tags} are the second type of annotation,
provided by spaCy models. POS tags are labels that provide a grammatical
analysis of each token in the text, e.g. identifying words as nouns or
adjectives. In contrast to the similarity annotations, POS tags are not
comparisons that result in a score. Instead, they can be used in any
algorithm that operates on token\sphinxhyphen{}level attributes. In other words, you
can use these tags much in the same way that you might use the token
itself in a specific position. For example, the tree shown earlier
contains a select algorithm with the condition that the token left to
the node has to be an adjective.

\sphinxstepscope


\chapter{Subsets}
\label{\detokenize{clicanalysis/subsets:subsets}}\label{\detokenize{clicanalysis/subsets:id1}}\label{\detokenize{clicanalysis/subsets::doc}}
\sphinxAtStartPar
The Subsets tab can display the full subset of your choice for the
selected corpus. Therefore, you can retrieve all quotes or all long
suspensions, etc. in any of the books or pre\sphinxhyphen{}selected corpora for
further analysis. Note that we find this option most useful for the
smaller subsets, i.e. quotes and suspensions; if you select the whole
‘non\sphinxhyphen{}quotes’ subset the output may become unwieldy.


\section{Show subsets}
\label{\detokenize{clicanalysis/subsets:show-subsets}}
\sphinxAtStartPar
Click onto the dropdown \sphinxstylestrong{‘Show subsets’} (see \hyperref[\detokenize{clicanalysis/subsets:figure-analysis-subsets-show-options}]{Fig.\@ \ref{\detokenize{clicanalysis/subsets:figure-analysis-subsets-show-options}}}) to select a relevant
subset (short suspensions, long suspensions, quotes or non\sphinxhyphen{}quotes). You
will also need to choose a corpus.

\begin{figure}[htbp]
\centering
\capstart

\noindent\sphinxincludegraphics{{figure-analysis-subsets-show-options}.png}
\caption{The basic subset options}\label{\detokenize{clicanalysis/subsets:id2}}\label{\detokenize{clicanalysis/subsets:figure-analysis-subsets-show-options}}\end{figure}

\sphinxAtStartPar
\hyperref[\detokenize{clicanalysis/subsets:figure-analysis-subsets-show-longsuspensions}]{Fig.\@ \ref{\detokenize{clicanalysis/subsets:figure-analysis-subsets-show-longsuspensions}}} shows sample
lines from the subset of long suspensions in \sphinxstyleemphasis{Oliver Twist}. You can
then use the filter option to narrow down the lines and group them using
the KWICGrouper. For subsets, the “relative frequency” is not given in terms of
frequency per million words, as in the Concordance tab, but as the percentage of
total words in the corpus found in the selected subset.

\begin{figure}[htbp]
\centering
\capstart

\noindent\sphinxincludegraphics{{figure-analysis-subsets-show-longsuspensions}.png}
\caption{The first few lines from the subset of ‘long suspensions’
in Oliver Twist}\label{\detokenize{clicanalysis/subsets:id3}}\label{\detokenize{clicanalysis/subsets:figure-analysis-subsets-show-longsuspensions}}\end{figure}


\section{Results}
\label{\detokenize{clicanalysis/subsets:results}}
\sphinxAtStartPar
Like in the concordance tab, this allows you to adjust the way the
concordance output (‘table’) is displayed.


\subsection{Filter rows}
\label{\detokenize{clicanalysis/subsets:filter-rows}}
\sphinxAtStartPar
The filter option lets you filter the output by the rows that contain a
particular sequence of letters, as described in the {\hyperref[\detokenize{clicanalysis/subsets:filter-rows}]{\sphinxcrossref{\DUrole{std,std-ref}{Filter rows}}}}
subsection of the Concordance tab documentation. For example, you could filter
suspensions for particular speech verbs like \sphinxstyleemphasis{cried}
(\hyperref[\detokenize{clicanalysis/subsets:figure-analysis-subsets-results-filter-cried}]{Fig.\@ \ref{\detokenize{clicanalysis/subsets:figure-analysis-subsets-results-filter-cried}}}).

\begin{figure}[htbp]
\centering
\capstart

\noindent\sphinxincludegraphics{{figure-analysis-subsets-results-filter-cried}.png}
\caption{Filtering long suspensions in Oliver Twist for \sphinxstyleemphasis{cried}}\label{\detokenize{clicanalysis/subsets:id4}}\label{\detokenize{clicanalysis/subsets:figure-analysis-subsets-results-filter-cried}}\end{figure}

\begin{figure}[htbp]
\centering
\capstart

\noindent\sphinxincludegraphics{{figure-analysis-subsets-results-filter-cotext}.png}
\caption{Filtering the co\sphinxhyphen{}text of long suspensions for \sphinxstyleemphasis{perhaps} in
Oliver Twist}\label{\detokenize{clicanalysis/subsets:id5}}\label{\detokenize{clicanalysis/subsets:figure-analysis-subsets-results-filter-cotext}}\end{figure}

\sphinxAtStartPar
Note, however, that the filter will search through the whole row and
therefore also accounts for words in the context, not only in the subset
itself. For example, when searching through the subset of long
suspensions in \sphinxstyleemphasis{Oliver Twist} and filtering rows for \sphinxstyleemphasis{perhaps} the
results originate only from the co\sphinxhyphen{}text, as \sphinxstyleemphasis{perhaps} does not occur in
long suspensions (see \hyperref[\detokenize{clicanalysis/subsets:figure-analysis-subsets-results-filter-cotext}]{Fig.\@ \ref{\detokenize{clicanalysis/subsets:figure-analysis-subsets-results-filter-cotext}}}).


\subsection{View as}
\label{\detokenize{clicanalysis/subsets:view-as}}
\sphinxAtStartPar
Like the {\hyperref[\detokenize{clicanalysis/subsets:view-as}]{\sphinxcrossref{\DUrole{std,std-ref}{View as}}}} options for the Concordance tab, in Subsets you
can view the ‘Basic results’ (concordance lines; book short title; link
to ‘in bk.’ view) the ‘full metadata’ (+ chapter, paragraph \& sentence
numbers) or the ‘distribution plot’, which gives an overview of matching
lines per book.

\begin{figure}[htbp]
\centering
\capstart

\noindent\sphinxincludegraphics{{figure_distribution_plot_quotes}.png}
\caption{The distribution plot view in the Subsets tab}\label{\detokenize{clicanalysis/subsets:id6}}\label{\detokenize{clicanalysis/subsets:figure-distribution-plot-quotes}}\end{figure}

\sphinxAtStartPar
In the case of the Subsets tab, these lines obviously are not concordance
lines, but instances of the subset e.g. a quote or non\sphinxhyphen{}quote element or
a short/long suspension. When you then create a distribution plot of a
selected subset, you will therefore see how the subset is distributed
across a book or corpus. Note that this operation may take a moment to
load for a large corpus. \hyperref[\detokenize{clicanalysis/subsets:figure-distribution-plot-quotes}]{Fig.\@ \ref{\detokenize{clicanalysis/subsets:figure-distribution-plot-quotes}}} gives
an example of three particular books with rather distinct quote
distributions: whereas \sphinxstyleemphasis{Pride and Prejudice} contains a lot of dialogue
\textendash{} as you can see from the white quote subsets interspersed by grey
non\sphinxhyphen{}quotes \textendash{} both \sphinxstyleemphasis{The Time Machine} and \sphinxstyleemphasis{Heart of Darkness} contain
much longer quote chunks by a key character telling a story.


\subsection{KWICGrouper}
\label{\detokenize{clicanalysis/subsets:kwicgrouper}}
\sphinxAtStartPar
If you want to restrict your search to the subset itself, the
KWICGrouper is the better option; it will also highlight your search
terms, as described in the {\hyperref[\detokenize{clicanalysis/concordance:concordance}]{\sphinxcrossref{\DUrole{std,std-ref}{Concordance}}}} section. The Subset
KWICGrouper works like the Concordance KWICGrouper, with the exception
of its search span which operates only on the subset itself. See
\hyperref[\detokenize{clicanalysis/subsets:figure-analysis-subsets-kwicgrouper-criedscreamedsobbed}]{Fig.\@ \ref{\detokenize{clicanalysis/subsets:figure-analysis-subsets-kwicgrouper-criedscreamedsobbed}}}
for an illustration of the Subset KWICGrouper searching for lines with
\sphinxstyleemphasis{cried}, \sphinxstyleemphasis{screamed} and \sphinxstyleemphasis{sobbed}.

\begin{figure}[htbp]
\centering
\capstart

\noindent\sphinxincludegraphics{{figure-analysis-subsets-kwicgrouper-criedscreamedsobbed}.png}
\caption{The search span of the Subset KWICGrouper applies to the
subset; not to the co\sphinxhyphen{}text}\label{\detokenize{clicanalysis/subsets:id7}}\label{\detokenize{clicanalysis/subsets:figure-analysis-subsets-kwicgrouper-criedscreamedsobbed}}\end{figure}


\subsection{Manage tag columns}
\label{\detokenize{clicanalysis/subsets:manage-tag-columns}}
\sphinxAtStartPar
Just like in the Concordance tab (see {\hyperref[\detokenize{clicanalysis/concordance:concordance}]{\sphinxcrossref{\DUrole{std,std-ref}{Concordance}}}}), subset rows can be
annotated with user\sphinxhyphen{}defined tags.
\hyperref[\detokenize{clicanalysis/subsets:figure-analysis-subsets-tagcolumns-gender}]{Fig.\@ \ref{\detokenize{clicanalysis/subsets:figure-analysis-subsets-tagcolumns-gender}}} shows a
potential application of tagging subsets: long suspensions in the 19th
Century Children’s Literature (ChiLit) corpus containing \sphinxstyleemphasis{cried} are
tagged for whether the crying character is male or female. Note that
this screenshot just illustrates the technique; it does not represent
the actual gender distribution of \sphinxstyleemphasis{cried} in the ChiLit long
suspensions.

\begin{figure}[htbp]
\centering
\capstart

\noindent\sphinxincludegraphics{{figure-analysis-subsets-tagcolumns-gender}.png}
\caption{Tagging subsets \textendash{} here, long suspensions in ChiLit
containing \sphinxstyleemphasis{cried} are tagged for character gender}\label{\detokenize{clicanalysis/subsets:id8}}\label{\detokenize{clicanalysis/subsets:figure-analysis-subsets-tagcolumns-gender}}\end{figure}

\sphinxstepscope


\chapter{Clusters}
\label{\detokenize{clicanalysis/clusters:clusters}}\label{\detokenize{clicanalysis/clusters:id1}}\label{\detokenize{clicanalysis/clusters::doc}}
\sphinxAtStartPar
The output of the cluster tool generates frequency lists of single words
and ‘clusters’ (repeated sequences of words). Clusters are also called
‘n\sphinxhyphen{}grams’, where ‘n’ stands for the length of the phrase. If we choose a
‘1\sphinxhyphen{}gram’ (single word), we retrieve a simple word list. (In \sphinxstyleemphasis{Oliver
Twist}, for example, the top 10 words retrieved from this tool are \sphinxstyleemphasis{the,
and, to, of, a, he, in, his, that} \textendash{} all function words, as we would
generally expect.) From version 2.0 onwards, CLiC supports clusters of
length 1 (single words) up to 7 (\sphinxtitleref{i am very much obliged to you}), as
illustrated in \hyperref[\detokenize{clicanalysis/clusters:figure-analysis-clusters-ngrams}]{Fig.\@ \ref{\detokenize{clicanalysis/clusters:figure-analysis-clusters-ngrams}}}.

\begin{figure}[htbp]
\centering
\capstart

\noindent\sphinxincludegraphics{{figure-analysis-clusters-ngrams}.png}
\caption{Cluster options}\label{\detokenize{clicanalysis/clusters:id2}}\label{\detokenize{clicanalysis/clusters:figure-analysis-clusters-ngrams}}\end{figure}

\sphinxAtStartPar
As in the other tabs, you can restrict the search to a particular subset
(\sphinxstylestrong{‘Only in subsets: Select an Option’}) so that, for example, you can
create frequency lists for clusters in quotes (or any of the other
subsets). You can save the resulting list as a CSV file (for example for
use in a spreadsheet viewer) by clicking the \sphinxstylestrong{‘Save’} button at the
top. Note that the CLiC ‘Cluster’ tab will display words and clusters
with a minimum frequency of 5.

\sphinxstepscope


\chapter{Keywords}
\label{\detokenize{clicanalysis/keywords:keywords}}\label{\detokenize{clicanalysis/keywords:id1}}\label{\detokenize{clicanalysis/keywords::doc}}
\sphinxAtStartPar
The keywords tool finds words (and phrases) that are used significantly
more often in one corpus compared to another. CLiC incorporates the
keyword extraction formula reported by \sphinxcite{references:rayson-garside-2000}. Apart
from comparing single words, CLiC also allows you to compare clusters.
Whereas the cluster tab focuses only on one corpus, the Keywords
function can compare cluster lists. You have to make selections for the
following options (also see
\hyperref[\detokenize{clicanalysis/keywords:figure-analysis-keywords-settings}]{Fig.\@ \ref{\detokenize{clicanalysis/keywords:figure-analysis-keywords-settings}}}):
\begin{itemize}
\item {} 
\sphinxAtStartPar
\sphinxstylestrong{‘Target corpora’:} Choose the corpus/corpora that you are
interested in.
\begin{itemize}
\item {} 
\sphinxAtStartPar
‘within subset’: Specify which subset of the target corpus you
want to compare (or simply choose ‘all text’)

\end{itemize}

\item {} 
\sphinxAtStartPar
\sphinxstylestrong{‘Reference corpora’:} Choose the reference corpus to compare your
target corpus to.
\begin{itemize}
\item {} 
\sphinxAtStartPar
‘within subset’: Specify the subset for the reference corpus.

\end{itemize}

\item {} 
\sphinxAtStartPar
\sphinxstylestrong{‘n\sphinxhyphen{}gram’:} Do you want to compare single words (1\sphinxhyphen{}grams) or
phrases (2\sphinxhyphen{}grams up to 7\sphinxhyphen{}grams)

\end{itemize}

\begin{figure}[htbp]
\centering
\capstart

\noindent\sphinxincludegraphics{{figure-analysis-keywords-settings}.png}
\caption{The settings for the keywords tab require you to select
two sets of corpora for the keyword comparison \textendash{} target and reference
\textendash{} and their corresponding subsets}\label{\detokenize{clicanalysis/keywords:id4}}\label{\detokenize{clicanalysis/keywords:figure-analysis-keywords-settings}}\end{figure}

\begin{figure}[htbp]
\centering
\capstart

\noindent\sphinxincludegraphics{{figure-analysis-keywords-19thcentury}.png}
\caption{Key 5\sphinxhyphen{}word clusters in Oliver Twist ‘quotes’ compared to
‘quotes’ in the 19th Century Reference Corpus}\label{\detokenize{clicanalysis/keywords:id5}}\label{\detokenize{clicanalysis/keywords:figure-analysis-keywords-19thcentury}}\end{figure}

\sphinxAtStartPar
Note that you have to select a subset for each of the two corpora or
you’ll see the error message: “Please select a subset”. So, for example,
when comparing 5\sphinxhyphen{}grams in \sphinxstyleemphasis{Oliver Twist} (quotes) against the 19th
Century Reference Corpus (quotes), we retrieve the results displayed in
\hyperref[\detokenize{clicanalysis/keywords:figure-analysis-keywords-19thcentury}]{Fig.\@ \ref{\detokenize{clicanalysis/keywords:figure-analysis-keywords-19thcentury}}} (for a p\sphinxhyphen{}value of
0.0001). The keyword output is by default ordered by the log\sphinxhyphen{}likelihood
(LL) value, the ‘keyness’ statistic used here (for more details on the
calculation, please refer to \sphinxcite{references:rayson-garside-2000}).

\sphinxAtStartPar
The frequency threshold of 5 used for the cluster tab is not applied to
the keyword tab, so that all frequencies are compared. The keyword
output shows the top 3000 results (for most comparisons, you will yield
fewer results, though). Moreover, CLiC only generates so\sphinxhyphen{}called
‘positive keywords’: those that are overused in the target corpus
compared to the reference corpus. CLiC does not generate ‘negative’ or
‘underused’ keywords.

\sphinxstepscope


\chapter{Counts}
\label{\detokenize{clicanalysis/counts:counts}}\label{\detokenize{clicanalysis/counts:id1}}\label{\detokenize{clicanalysis/counts::doc}}
\sphinxAtStartPar
The counts tab lists information about individual books and all books in
a corpus:
\begin{itemize}
\item {} 
\sphinxAtStartPar
Book title

\item {} 
\sphinxAtStartPar
Number of chapters

\item {} 
\sphinxAtStartPar
Total number of words

\item {} 
\sphinxAtStartPar
Number of words in quotes

\item {} 
\sphinxAtStartPar
Number of words in non\sphinxhyphen{}quotes

\item {} 
\sphinxAtStartPar
Number of words in short suspensions

\item {} 
\sphinxAtStartPar
Number of words in long suspensions

\end{itemize}

\sphinxAtStartPar
You can create a counts table for a single book or a selection of books.
\hyperref[\detokenize{clicanalysis/counts:figure-analysis-counts}]{Fig.\@ \ref{\detokenize{clicanalysis/counts:figure-analysis-counts}}} shows the complete table for all books
in CLiC (this may take a while to load).

\begin{figure}[htbp]
\centering
\capstart

\noindent\sphinxincludegraphics{{figure-analysis-counts}.png}
\caption{Counts overview for all corpora in CLiC}\label{\detokenize{clicanalysis/counts:id2}}\label{\detokenize{clicanalysis/counts:figure-analysis-counts}}\end{figure}

\sphinxstepscope


\chapter{Texts}
\label{\detokenize{clicanalysis/texts:texts}}\label{\detokenize{clicanalysis/texts:id1}}\label{\detokenize{clicanalysis/texts::doc}}
\sphinxAtStartPar
The ‘Texts’ tab shows the full text of individual books and allows you
to navigate to particular chapters. Selections of the text can be easily
copied and pasted into other applications (e.g. Microsoft Word
documents).

\sphinxAtStartPar
By default the text view is plain, i.e. it won’t show any annotation
unless you tick the options. You can select the levels of annotation you
want to display, such as “Sentences”, and the “Quote” and “Non\sphinxhyphen{}quote”
subsets, etc. \hyperref[\detokenize{clicanalysis/texts:figure-analysis-texts}]{Fig.\@ \ref{\detokenize{clicanalysis/texts:figure-analysis-texts}}} illustrates the annotation
views for quotes, short suspensions and long suspensions in \sphinxstyleemphasis{Pride and
Prejudice}.

\begin{figure}[htbp]
\centering
\capstart

\noindent\sphinxincludegraphics{{figure-analysis-texts}.png}
\caption{View of Chapter 3 in \sphinxstyleemphasis{Pride and Prejudice}}\label{\detokenize{clicanalysis/texts:id2}}\label{\detokenize{clicanalysis/texts:figure-analysis-texts}}\end{figure}

\sphinxstepscope


\chapter{CLiC API}
\label{\detokenize{advanced/api_usage:clic-api}}\label{\detokenize{advanced/api_usage:api-usage}}\label{\detokenize{advanced/api_usage::doc}}
\begin{sphinxShadowBox}
\begin{itemize}
\item {} 
\sphinxAtStartPar
\phantomsection\label{\detokenize{advanced/api_usage:id1}}{\hyperref[\detokenize{advanced/api_usage:overview}]{\sphinxcrossref{Overview}}}

\item {} 
\sphinxAtStartPar
\phantomsection\label{\detokenize{advanced/api_usage:id2}}{\hyperref[\detokenize{advanced/api_usage:example-code}]{\sphinxcrossref{Example code}}}
\begin{itemize}
\item {} 
\sphinxAtStartPar
\phantomsection\label{\detokenize{advanced/api_usage:id3}}{\hyperref[\detokenize{advanced/api_usage:python-3}]{\sphinxcrossref{Python 3}}}

\item {} 
\sphinxAtStartPar
\phantomsection\label{\detokenize{advanced/api_usage:id4}}{\hyperref[\detokenize{advanced/api_usage:r}]{\sphinxcrossref{R}}}

\end{itemize}

\end{itemize}
\end{sphinxShadowBox}


\section{Overview}
\label{\detokenize{advanced/api_usage:overview}}
\sphinxAtStartPar
Any data available through the \sphinxhref{https://clic-fiction.com}{CLiC web interface}
is also available by directly calling the \sphinxstyleemphasis{CLiC API}. The CLiC API
returns a JSON representation of the CLiC data, which means that the
data can be retrieved directly using any programming language. To get you
started we have written some example code for both Python and R.

\sphinxAtStartPar
The sample code is for the \sphinxstyleemphasis{corpora} and \sphinxstyleemphasis{subset} endpoints. The corpora
endpoint is used to retrieve the list of resources available in CLiC.
This can then be used by your code to filter or select the resources you
want to fetch (see the example usage). In the example code the corpora
endpoint is called using the \sphinxcode{\sphinxupquote{get\_lookup()}} function which returns
content something like:

\begin{sphinxVerbatim}[commandchars=\\\{\}]
\PYG{o}{\PYGZgt{}} \PYG{n}{lookup}
     \PYG{n}{corpus}                      \PYG{n}{author}     \PYG{n+nb}{id}               \PYG{n}{title}
  \PYG{l+m+mi}{1}\PYG{p}{:} \PYG{n}{ChiLit}            \PYG{n}{Agnes} \PYG{n}{Strickland}  \PYG{n}{rival}   \PYG{n}{The} \PYG{n}{Rival} \PYG{n}{Crusoes}
  \PYG{l+m+mi}{2}\PYG{p}{:} \PYG{n}{ChiLit}                 \PYG{n}{Andrew} \PYG{n}{Lang} \PYG{n}{prigio}       \PYG{n}{Prince} \PYG{n}{Prigio}
  \PYG{l+m+mi}{3}\PYG{p}{:} \PYG{n}{ChiLit}           \PYG{n}{Ann} \PYG{n}{Fraser} \PYG{n}{Tytler}  \PYG{n}{leila}       \PYG{n}{Leila} \PYG{n}{at} \PYG{n}{Home}
 \PYG{o}{\PYGZhy{}}\PYG{o}{\PYGZhy{}}\PYG{o}{\PYGZhy{}}
\PYG{l+m+mi}{136}\PYG{p}{:}    \PYG{n}{ntc}              \PYG{n}{Wilkie} \PYG{n}{Collins}   \PYG{n}{arma}            \PYG{n}{Armadale}
\PYG{l+m+mi}{137}\PYG{p}{:}    \PYG{n}{ntc}              \PYG{n}{Wilkie} \PYG{n}{Collins} \PYG{n}{wwhite}  \PYG{n}{The} \PYG{n}{Woman} \PYG{o+ow}{in} \PYG{n}{White}
\PYG{l+m+mi}{138}\PYG{p}{:}    \PYG{n}{ntc} \PYG{n}{William} \PYG{n}{Makepeace} \PYG{n}{Thackeray} \PYG{n}{vanity}         \PYG{n}{Vanity} \PYG{n}{Fair}
\end{sphinxVerbatim}

\sphinxAtStartPar
The subset endpoint is used to retrieve tokenized text for all or parts
of one or more of the corpora. This endpoint is called ‘subset’ because
you can restrict the retrieved tokens to a specific subset of the whole
text, these subsets are: quoted text, non\sphinxhyphen{}quoted text, long suspensions
or short suspensions. In the example code the subset endpoint is called
using the \sphinxcode{\sphinxupquote{get\_tokens()}} function. The subset endpoint is documented
at \sphinxcode{\sphinxupquote{clic.subset}}.

\sphinxAtStartPar
The \sphinxstyleemphasis{cluster} endpoint is used to retrieve n\sphinxhyphen{}grams and their counts. In
the example code the clusters endpoint is called using the
\sphinxcode{\sphinxupquote{get\_clusters()}} function. The cluster endpoint is documented at
\sphinxcode{\sphinxupquote{clic.cluster}}.

\sphinxAtStartPar
We would be interested to hear about how you use the CLiC API and are
always happy to consider CLiC related guest posts for the
\sphinxhref{https://blog.bham.ac.uk/clic-dickens/}{CLiC blog}. To let us know
how you are using the CLiC API, to give us feedback, or if you need any
help, you can contact us at \sphinxhref{mailto:dhss-kontakt@fau.de}{dhss\sphinxhyphen{}kontakt@fau.de}.

\sphinxAtStartPar
To help us understand who is using the API, when writing code to access
the API please set the “User\sphinxhyphen{}Agent” header to something that identifies
you or your application.

\sphinxAtStartPar
The CLiC API uses the legacy names for the CLiC corpora. The following
table gives the correspondence between the corpora names as seen in the
CLiC web interface and those used by the CLiC API.


\begin{savenotes}\sphinxattablestart
\sphinxthistablewithglobalstyle
\centering
\begin{tabulary}{\linewidth}[t]{TTT}
\sphinxtoprule
\sphinxstyletheadfamily 
\sphinxAtStartPar
CLiC Web
&\sphinxstyletheadfamily 
\sphinxAtStartPar
CLiC API
&\sphinxstyletheadfamily 
\sphinxAtStartPar
Description
\\
\sphinxmidrule
\sphinxtableatstartofbodyhook
\sphinxAtStartPar
\sphinxcode{\sphinxupquote{DNov}}
&
\sphinxAtStartPar
\sphinxcode{\sphinxupquote{dickens}}
&
\sphinxAtStartPar
Dickens’s Novels
\\
\sphinxhline
\sphinxAtStartPar
\sphinxcode{\sphinxupquote{19C}}
&
\sphinxAtStartPar
\sphinxcode{\sphinxupquote{ntc}}
&
\sphinxAtStartPar
19th Century Reference Corpus
\\
\sphinxhline
\sphinxAtStartPar
\sphinxcode{\sphinxupquote{ChiLit}}
&
\sphinxAtStartPar
\sphinxcode{\sphinxupquote{ChiLit}}
&
\sphinxAtStartPar
19th Century Children’s Literature Corpus
\\
\sphinxhline
\sphinxAtStartPar
\sphinxcode{\sphinxupquote{ArTs}}
&
\sphinxAtStartPar
\sphinxcode{\sphinxupquote{Other}}
&
\sphinxAtStartPar
Additional Requested Texts
\\
\sphinxhline
\sphinxAtStartPar
\sphinxcode{\sphinxupquote{DE19}}
&
\sphinxAtStartPar
\sphinxcode{\sphinxupquote{DE19}}
&
\sphinxAtStartPar
German 19th Century Reference Corpus
\\
\sphinxbottomrule
\end{tabulary}
\sphinxtableafterendhook\par
\sphinxattableend\end{savenotes}


\section{Example code}
\label{\detokenize{advanced/api_usage:example-code}}

\subsection{Python 3}
\label{\detokenize{advanced/api_usage:python-3}}
\begin{sphinxVerbatim}[commandchars=\\\{\}]
\PYG{k+kn}{import} \PYG{n+nn}{json}
\PYG{k+kn}{from} \PYG{n+nn}{operator} \PYG{k+kn}{import} \PYG{n}{itemgetter}
\PYG{k+kn}{from} \PYG{n+nn}{collections} \PYG{k+kn}{import} \PYG{n}{OrderedDict}
\PYG{k+kn}{import} \PYG{n+nn}{requests}
\PYG{k+kn}{import} \PYG{n+nn}{pandas} \PYG{k}{as} \PYG{n+nn}{pd}

\PYG{n}{UA} \PYG{o}{=} \PYG{l+s+s2}{\PYGZdq{}}\PYG{l+s+s2}{CLiC API Example Python3 Code}\PYG{l+s+s2}{\PYGZdq{}}  \PYG{c+c1}{\PYGZsh{} user agent !! CHANGE ME !!}
\PYG{n}{HOSTNAME} \PYG{o}{=} \PYG{l+s+s2}{\PYGZdq{}}\PYG{l+s+s2}{clic\PYGZhy{}fiction.com}\PYG{l+s+s2}{\PYGZdq{}}

\PYG{k}{def} \PYG{n+nf}{api\PYGZus{}request}\PYG{p}{(}\PYG{n}{endpoint}\PYG{p}{,} \PYG{n}{query}\PYG{o}{=}\PYG{k+kc}{None}\PYG{p}{)}\PYG{p}{:}
\PYG{+w}{    }\PYG{l+s+sd}{\PYGZdq{}\PYGZdq{}\PYGZdq{}}
\PYG{l+s+sd}{    Makes the API requests.}
\PYG{l+s+sd}{    Returns the endpoint specific data structure as a python structure.}

\PYG{l+s+sd}{    \PYGZhy{} endpoint: see the inline docs in /server/clic}
\PYG{l+s+sd}{    \PYGZhy{} query: endpoint specific parameters as a querystring}
\PYG{l+s+sd}{    \PYGZdq{}\PYGZdq{}\PYGZdq{}}
    \PYG{k}{if} \PYG{n}{query} \PYG{o+ow}{is} \PYG{k+kc}{None}\PYG{p}{:}
        \PYG{n}{uri} \PYG{o}{=} \PYG{l+s+s1}{\PYGZsq{}}\PYG{l+s+s1}{https://}\PYG{l+s+si}{\PYGZpc{}s}\PYG{l+s+s1}{/api/}\PYG{l+s+si}{\PYGZpc{}s}\PYG{l+s+s1}{\PYGZsq{}} \PYG{o}{\PYGZpc{}} \PYG{p}{(}\PYG{n}{HOSTNAME}\PYG{p}{,} \PYG{n}{endpoint}\PYG{p}{)}
    \PYG{k}{else}\PYG{p}{:}
        \PYG{n}{uri} \PYG{o}{=} \PYG{l+s+s1}{\PYGZsq{}}\PYG{l+s+s1}{https://}\PYG{l+s+si}{\PYGZpc{}s}\PYG{l+s+s1}{/api/}\PYG{l+s+si}{\PYGZpc{}s}\PYG{l+s+s1}{?}\PYG{l+s+si}{\PYGZpc{}s}\PYG{l+s+s1}{\PYGZsq{}} \PYG{o}{\PYGZpc{}} \PYG{p}{(}\PYG{n}{HOSTNAME}\PYG{p}{,} \PYG{n}{endpoint}\PYG{p}{,} \PYG{n}{query}\PYG{p}{)}
    \PYG{n}{resp} \PYG{o}{=} \PYG{n}{requests}\PYG{o}{.}\PYG{n}{get}\PYG{p}{(}\PYG{n}{uri}\PYG{p}{,} \PYG{n}{headers}\PYG{o}{=}\PYG{p}{\PYGZob{}}\PYG{l+s+s1}{\PYGZsq{}}\PYG{l+s+s1}{User\PYGZhy{}agent}\PYG{l+s+s1}{\PYGZsq{}}\PYG{p}{:} \PYG{n}{UA}\PYG{p}{,} \PYG{l+s+s1}{\PYGZsq{}}\PYG{l+s+s1}{Accept}\PYG{l+s+s1}{\PYGZsq{}}\PYG{p}{:} \PYG{l+s+s1}{\PYGZsq{}}\PYG{l+s+s1}{application/json}\PYG{l+s+s1}{\PYGZsq{}}\PYG{p}{\PYGZcb{}}\PYG{p}{)}
    \PYG{k}{try}\PYG{p}{:}
        \PYG{n}{rv} \PYG{o}{=} \PYG{n}{resp}\PYG{o}{.}\PYG{n}{json}\PYG{p}{(}\PYG{p}{)}
    \PYG{k}{except} \PYG{n}{json}\PYG{o}{.}\PYG{n}{decoder}\PYG{o}{.}\PYG{n}{JSONDecodeError}\PYG{p}{:}
        \PYG{n+nb}{print}\PYG{p}{(}\PYG{l+s+s2}{\PYGZdq{}}\PYG{l+s+s2}{API request did not return valid JSON}\PYG{l+s+s2}{\PYGZdq{}}\PYG{p}{)}
    \PYG{k}{if} \PYG{n}{rv}\PYG{o}{.}\PYG{n}{get}\PYG{p}{(}\PYG{l+s+s1}{\PYGZsq{}}\PYG{l+s+s1}{error}\PYG{l+s+s1}{\PYGZsq{}}\PYG{p}{,} \PYG{k+kc}{False}\PYG{p}{)}\PYG{p}{:}
        \PYG{k}{raise} \PYG{n+ne}{ValueError}\PYG{p}{(}\PYG{l+s+s2}{\PYGZdq{}}\PYG{l+s+s2}{API returned error: }\PYG{l+s+s2}{\PYGZdq{}} \PYG{o}{+} \PYG{n}{rv}\PYG{p}{[}\PYG{l+s+s1}{\PYGZsq{}}\PYG{l+s+s1}{error}\PYG{l+s+s1}{\PYGZsq{}}\PYG{p}{]}\PYG{p}{[}\PYG{l+s+s1}{\PYGZsq{}}\PYG{l+s+s1}{message}\PYG{l+s+s1}{\PYGZsq{}}\PYG{p}{]}\PYG{p}{)}
    \PYG{k}{if} \PYG{n}{rv}\PYG{o}{.}\PYG{n}{get}\PYG{p}{(}\PYG{l+s+s1}{\PYGZsq{}}\PYG{l+s+s1}{warn}\PYG{l+s+s1}{\PYGZsq{}}\PYG{p}{,} \PYG{k+kc}{False}\PYG{p}{)}\PYG{p}{:}
        \PYG{n+nb}{print}\PYG{p}{(}\PYG{l+s+s2}{\PYGZdq{}}\PYG{l+s+s2}{API returned warning: }\PYG{l+s+s2}{\PYGZdq{}} \PYG{o}{+} \PYG{n}{rv}\PYG{p}{[}\PYG{l+s+s1}{\PYGZsq{}}\PYG{l+s+s1}{warn}\PYG{l+s+s1}{\PYGZsq{}}\PYG{p}{]}\PYG{p}{[}\PYG{l+s+s1}{\PYGZsq{}}\PYG{l+s+s1}{message}\PYG{l+s+s1}{\PYGZsq{}}\PYG{p}{]}\PYG{p}{)}
    \PYG{k}{if} \PYG{n}{rv}\PYG{o}{.}\PYG{n}{get}\PYG{p}{(}\PYG{l+s+s1}{\PYGZsq{}}\PYG{l+s+s1}{info}\PYG{l+s+s1}{\PYGZsq{}}\PYG{p}{,} \PYG{k+kc}{False}\PYG{p}{)}\PYG{p}{:}
        \PYG{n+nb}{print}\PYG{p}{(}\PYG{l+s+s2}{\PYGZdq{}}\PYG{l+s+s2}{API returned info: }\PYG{l+s+s2}{\PYGZdq{}} \PYG{o}{+} \PYG{n}{rv}\PYG{p}{[}\PYG{l+s+s1}{\PYGZsq{}}\PYG{l+s+s1}{info}\PYG{l+s+s1}{\PYGZsq{}}\PYG{p}{]}\PYG{p}{[}\PYG{l+s+s1}{\PYGZsq{}}\PYG{l+s+s1}{message}\PYG{l+s+s1}{\PYGZsq{}}\PYG{p}{]}\PYG{p}{)}
    \PYG{k}{return} \PYG{n}{rv}

\PYG{k}{def} \PYG{n+nf}{get\PYGZus{}lookup}\PYG{p}{(}\PYG{p}{)}\PYG{p}{:}
\PYG{+w}{    }\PYG{l+s+sd}{\PYGZdq{}\PYGZdq{}\PYGZdq{}}
\PYG{l+s+sd}{    Returns a pandas DataFrame listing the texts for each of the available corpora.}
\PYG{l+s+sd}{    \PYGZdq{}\PYGZdq{}\PYGZdq{}}
    \PYG{n}{rv} \PYG{o}{=} \PYG{n}{api\PYGZus{}request}\PYG{p}{(}\PYG{n}{endpoint}\PYG{o}{=}\PYG{l+s+s2}{\PYGZdq{}}\PYG{l+s+s2}{corpora}\PYG{l+s+s2}{\PYGZdq{}}\PYG{p}{)}
    \PYG{n}{d} \PYG{o}{=} \PYG{p}{[}\PYG{p}{]}
    \PYG{k}{for} \PYG{n}{corpus} \PYG{o+ow}{in} \PYG{n}{rv}\PYG{p}{[}\PYG{l+s+s1}{\PYGZsq{}}\PYG{l+s+s1}{corpora}\PYG{l+s+s1}{\PYGZsq{}}\PYG{p}{]}\PYG{p}{:}
        \PYG{n}{corpus\PYGZus{}id} \PYG{o}{=} \PYG{n}{corpus}\PYG{p}{[}\PYG{l+s+s1}{\PYGZsq{}}\PYG{l+s+s1}{id}\PYG{l+s+s1}{\PYGZsq{}}\PYG{p}{]}
        \PYG{k}{for} \PYG{n}{book} \PYG{o+ow}{in} \PYG{n}{corpus}\PYG{p}{[}\PYG{l+s+s1}{\PYGZsq{}}\PYG{l+s+s1}{children}\PYG{l+s+s1}{\PYGZsq{}}\PYG{p}{]}\PYG{p}{:}
            \PYG{n}{d}\PYG{o}{.}\PYG{n}{append}\PYG{p}{(}\PYG{p}{\PYGZob{}}\PYG{l+s+s1}{\PYGZsq{}}\PYG{l+s+s1}{corpus}\PYG{l+s+s1}{\PYGZsq{}} \PYG{p}{:} \PYG{n}{corpus\PYGZus{}id}\PYG{p}{,} \PYG{l+s+s1}{\PYGZsq{}}\PYG{l+s+s1}{author}\PYG{l+s+s1}{\PYGZsq{}} \PYG{p}{:} \PYG{n}{book}\PYG{p}{[}\PYG{l+s+s1}{\PYGZsq{}}\PYG{l+s+s1}{author}\PYG{l+s+s1}{\PYGZsq{}}\PYG{p}{]}\PYG{p}{,} \PYGZbs{}
                      \PYG{l+s+s1}{\PYGZsq{}}\PYG{l+s+s1}{shortname}\PYG{l+s+s1}{\PYGZsq{}} \PYG{p}{:} \PYG{n}{book}\PYG{p}{[}\PYG{l+s+s1}{\PYGZsq{}}\PYG{l+s+s1}{id}\PYG{l+s+s1}{\PYGZsq{}}\PYG{p}{]}\PYG{p}{,} \PYG{l+s+s1}{\PYGZsq{}}\PYG{l+s+s1}{title}\PYG{l+s+s1}{\PYGZsq{}}\PYG{p}{:} \PYG{n}{book}\PYG{p}{[}\PYG{l+s+s1}{\PYGZsq{}}\PYG{l+s+s1}{title}\PYG{l+s+s1}{\PYGZsq{}}\PYG{p}{]}\PYG{p}{\PYGZcb{}}\PYG{p}{)}
    \PYG{n}{df} \PYG{o}{=} \PYG{n}{pd}\PYG{o}{.}\PYG{n}{DataFrame}\PYG{p}{(}\PYG{n}{d}\PYG{p}{,} \PYG{n}{columns}\PYG{o}{=}\PYG{p}{[}\PYG{l+s+s1}{\PYGZsq{}}\PYG{l+s+s1}{corpus}\PYG{l+s+s1}{\PYGZsq{}}\PYG{p}{,} \PYG{l+s+s1}{\PYGZsq{}}\PYG{l+s+s1}{author}\PYG{l+s+s1}{\PYGZsq{}}\PYG{p}{,} \PYG{l+s+s1}{\PYGZsq{}}\PYG{l+s+s1}{shortname}\PYG{l+s+s1}{\PYGZsq{}}\PYG{p}{,} \PYG{l+s+s1}{\PYGZsq{}}\PYG{l+s+s1}{title}\PYG{l+s+s1}{\PYGZsq{}}\PYG{p}{]}\PYG{p}{)}
    \PYG{n}{df}\PYG{o}{.}\PYG{n}{sort\PYGZus{}values}\PYG{p}{(}\PYG{p}{[}\PYG{l+s+s1}{\PYGZsq{}}\PYG{l+s+s1}{corpus}\PYG{l+s+s1}{\PYGZsq{}}\PYG{p}{,} \PYG{l+s+s1}{\PYGZsq{}}\PYG{l+s+s1}{author}\PYG{l+s+s1}{\PYGZsq{}}\PYG{p}{,} \PYG{l+s+s1}{\PYGZsq{}}\PYG{l+s+s1}{title}\PYG{l+s+s1}{\PYGZsq{}}\PYG{p}{]}\PYG{p}{,} \PYG{n}{inplace}\PYG{o}{=}\PYG{k+kc}{True}\PYG{p}{,} \PYG{n}{ascending}\PYG{o}{=}\PYG{k+kc}{True}\PYG{p}{)}
    \PYG{n}{df}\PYG{o}{.}\PYG{n}{reset\PYGZus{}index}\PYG{p}{(}\PYG{n}{inplace}\PYG{o}{=}\PYG{k+kc}{True}\PYG{p}{,} \PYG{n}{drop}\PYG{o}{=}\PYG{k+kc}{True}\PYG{p}{)}
    \PYG{k}{return} \PYG{n}{df}

\PYG{k}{def} \PYG{n+nf}{get\PYGZus{}tokens}\PYG{p}{(}\PYG{n}{shortname}\PYG{p}{,} \PYG{n}{subset}\PYG{o}{=}\PYG{k+kc}{None}\PYG{p}{,} \PYG{n}{lowercase}\PYG{o}{=}\PYG{k+kc}{True}\PYG{p}{,} \PYG{n}{punctuation}\PYG{o}{=}\PYG{k+kc}{False}\PYG{p}{)}\PYG{p}{:}
\PYG{+w}{    }\PYG{l+s+sd}{\PYGZdq{}\PYGZdq{}\PYGZdq{}}
\PYG{l+s+sd}{    Fetches tokens using the \PYGZsq{}subset\PYGZsq{} endpoint.}
\PYG{l+s+sd}{    Returns a list of tokens.}

\PYG{l+s+sd}{    \PYGZhy{} shortname: can be any value from the \PYGZsq{}corpus\PYGZsq{} or \PYGZsq{}shortname\PYGZsq{} columns returned}
\PYG{l+s+sd}{          by get\PYGZus{}lookup() can be a string or a list of strings}
\PYG{l+s+sd}{    \PYGZhy{} subset: any one of \PYGZdq{}shortsus\PYGZdq{}, \PYGZdq{}longsus\PYGZdq{}, \PYGZdq{}nonquote\PYGZdq{}, \PYGZdq{}quote\PYGZdq{}}
\PYG{l+s+sd}{    \PYGZhy{} lowercase: boolean indicating if the tokens should be transformed to lower case}
\PYG{l+s+sd}{    \PYGZhy{} punctuation: boolean indicating if punctuation tokens should be included}
\PYG{l+s+sd}{    \PYGZdq{}\PYGZdq{}\PYGZdq{}}
    \PYG{k}{if} \PYG{n+nb}{isinstance}\PYG{p}{(}\PYG{n}{shortname}\PYG{p}{,} \PYG{n+nb}{str}\PYG{p}{)}\PYG{p}{:}
        \PYG{n}{shortname} \PYG{o}{=} \PYG{p}{[}\PYG{n}{shortname}\PYG{p}{]}
    \PYG{n}{query} \PYG{o}{=} \PYG{l+s+s1}{\PYGZsq{}}\PYG{l+s+s1}{\PYGZam{}}\PYG{l+s+s1}{\PYGZsq{}}\PYG{o}{.}\PYG{n}{join}\PYG{p}{(}\PYG{p}{[}\PYG{l+s+s2}{\PYGZdq{}}\PYG{l+s+s2}{corpora=}\PYG{l+s+si}{\PYGZpc{}s}\PYG{l+s+s2}{\PYGZdq{}} \PYG{o}{\PYGZpc{}} \PYG{n}{sn} \PYG{k}{for} \PYG{n}{sn} \PYG{o+ow}{in} \PYG{n}{shortname}\PYG{p}{]}\PYG{p}{)}
    \PYG{k}{if} \PYG{n}{subset} \PYG{o+ow}{is} \PYG{o+ow}{not} \PYG{k+kc}{None}\PYG{p}{:}
        \PYG{k}{if} \PYG{n}{subset} \PYG{o+ow}{not} \PYG{o+ow}{in} \PYG{p}{[}\PYG{l+s+s2}{\PYGZdq{}}\PYG{l+s+s2}{shortsus}\PYG{l+s+s2}{\PYGZdq{}}\PYG{p}{,} \PYG{l+s+s2}{\PYGZdq{}}\PYG{l+s+s2}{longsus}\PYG{l+s+s2}{\PYGZdq{}}\PYG{p}{,} \PYG{l+s+s2}{\PYGZdq{}}\PYG{l+s+s2}{nonquote}\PYG{l+s+s2}{\PYGZdq{}}\PYG{p}{,} \PYG{l+s+s2}{\PYGZdq{}}\PYG{l+s+s2}{quote}\PYG{l+s+s2}{\PYGZdq{}}\PYG{p}{]}\PYG{p}{:}
            \PYG{k}{raise} \PYG{n+ne}{ValueError}\PYG{p}{(}\PYG{l+s+s1}{\PYGZsq{}}\PYG{l+s+s1}{bad subset parameter: }\PYG{l+s+s1}{\PYGZdq{}}\PYG{l+s+si}{\PYGZpc{}s}\PYG{l+s+s1}{\PYGZdq{}}\PYG{l+s+s1}{\PYGZsq{}} \PYG{o}{\PYGZpc{}} \PYG{n}{subset}\PYG{p}{)}
        \PYG{n}{query} \PYG{o}{=} \PYG{n}{query} \PYG{o}{+} \PYG{l+s+s2}{\PYGZdq{}}\PYG{l+s+s2}{\PYGZam{}subset=}\PYG{l+s+si}{\PYGZpc{}s}\PYG{l+s+s2}{\PYGZdq{}} \PYG{o}{\PYGZpc{}} \PYG{n}{subset}
    \PYG{n}{rv} \PYG{o}{=} \PYG{n}{api\PYGZus{}request}\PYG{p}{(}\PYG{n}{endpoint}\PYG{o}{=}\PYG{l+s+s2}{\PYGZdq{}}\PYG{l+s+s2}{subset}\PYG{l+s+s2}{\PYGZdq{}}\PYG{p}{,} \PYG{n}{query}\PYG{o}{=}\PYG{n}{query}\PYG{p}{)}
    \PYG{k}{if} \PYG{n}{punctuation}\PYG{p}{:}
        \PYG{n}{tokens} \PYG{o}{=} \PYG{p}{[}\PYG{n}{j} \PYG{k}{for} \PYG{n}{i} \PYG{o+ow}{in} \PYG{n}{rv}\PYG{p}{[}\PYG{l+s+s1}{\PYGZsq{}}\PYG{l+s+s1}{data}\PYG{l+s+s1}{\PYGZsq{}}\PYG{p}{]} \PYG{k}{for} \PYG{n}{j} \PYG{o+ow}{in} \PYG{n}{i}\PYG{p}{[}\PYG{l+m+mi}{0}\PYG{p}{]}\PYG{p}{[}\PYG{p}{:}\PYG{o}{\PYGZhy{}}\PYG{l+m+mi}{1}\PYG{p}{]}\PYG{p}{]}
    \PYG{k}{else}\PYG{p}{:}
        \PYG{n}{tokens} \PYG{o}{=} \PYG{p}{[}\PYG{n}{j} \PYG{k}{for} \PYG{n}{i} \PYG{o+ow}{in} \PYG{n}{rv}\PYG{p}{[}\PYG{l+s+s1}{\PYGZsq{}}\PYG{l+s+s1}{data}\PYG{l+s+s1}{\PYGZsq{}}\PYG{p}{]} \PYG{k}{for} \PYG{n}{j} \PYG{o+ow}{in} \PYG{p}{[}\PYG{n}{i}\PYG{p}{[}\PYG{l+m+mi}{0}\PYG{p}{]}\PYG{p}{[}\PYG{p}{:}\PYG{o}{\PYGZhy{}}\PYG{l+m+mi}{1}\PYG{p}{]}\PYG{p}{[}\PYG{n}{k}\PYG{p}{]} \PYG{k}{for} \PYG{n}{k} \PYG{o+ow}{in} \PYG{n}{i}\PYG{p}{[}\PYG{l+m+mi}{0}\PYG{p}{]}\PYG{p}{[}\PYG{o}{\PYGZhy{}}\PYG{l+m+mi}{1}\PYG{p}{]}\PYG{p}{]}\PYG{p}{]}
    \PYG{k}{if} \PYG{n}{lowercase}\PYG{p}{:}
        \PYG{k}{return} \PYG{p}{[}\PYG{n}{i}\PYG{o}{.}\PYG{n}{lower}\PYG{p}{(}\PYG{p}{)} \PYG{k}{for} \PYG{n}{i} \PYG{o+ow}{in} \PYG{n}{tokens}\PYG{p}{]}
    \PYG{k}{return} \PYG{n}{tokens}

\PYG{k}{def} \PYG{n+nf}{get\PYGZus{}clusters}\PYG{p}{(}\PYG{n}{shortname}\PYG{p}{,} \PYG{n}{length}\PYG{p}{,} \PYG{n}{cutoff}\PYG{o}{=}\PYG{l+m+mi}{5}\PYG{p}{,} \PYG{n}{subset}\PYG{o}{=}\PYG{k+kc}{None}\PYG{p}{)}\PYG{p}{:}
\PYG{+w}{    }\PYG{l+s+sd}{\PYGZdq{}\PYGZdq{}\PYGZdq{}}
\PYG{l+s+sd}{    Fetches n\PYGZhy{}grams using the \PYGZsq{}cluster\PYGZsq{} endpoint.}
\PYG{l+s+sd}{    Returns a OrderedDict of clusters to counts.}

\PYG{l+s+sd}{    \PYGZhy{} shortname: can be any value from the \PYGZsq{}corpus\PYGZsq{} or \PYGZsq{}shortname\PYGZsq{} columns returned}
\PYG{l+s+sd}{          by get\PYGZus{}lookup() can be a string or a list of strings}
\PYG{l+s+sd}{    \PYGZhy{} length: cluster length to search for, one of 1/3/4/5 (NB: There is no 2)}
\PYG{l+s+sd}{    \PYGZhy{} cutoff: [default: 5] the cutoff frequency, if a cluster occurs less times}
\PYG{l+s+sd}{          than this it is not returned}
\PYG{l+s+sd}{    \PYGZhy{} subset: [optional] any one of \PYGZdq{}shortsus\PYGZdq{}, \PYGZdq{}longsus\PYGZdq{}, \PYGZdq{}nonquote\PYGZdq{}, \PYGZdq{}quote\PYGZdq{}}
\PYG{l+s+sd}{    \PYGZdq{}\PYGZdq{}\PYGZdq{}}
    \PYG{k}{if} \PYG{n+nb}{isinstance}\PYG{p}{(}\PYG{n}{shortname}\PYG{p}{,} \PYG{n+nb}{str}\PYG{p}{)}\PYG{p}{:}
        \PYG{n}{shortname} \PYG{o}{=} \PYG{p}{[}\PYG{n}{shortname}\PYG{p}{]}
    \PYG{n}{query} \PYG{o}{=} \PYG{l+s+s1}{\PYGZsq{}}\PYG{l+s+s1}{\PYGZam{}}\PYG{l+s+s1}{\PYGZsq{}}\PYG{o}{.}\PYG{n}{join}\PYG{p}{(}\PYG{p}{[}\PYG{l+s+s2}{\PYGZdq{}}\PYG{l+s+s2}{corpora=}\PYG{l+s+si}{\PYGZpc{}s}\PYG{l+s+s2}{\PYGZdq{}} \PYG{o}{\PYGZpc{}} \PYG{n}{sn} \PYG{k}{for} \PYG{n}{sn} \PYG{o+ow}{in} \PYG{n}{shortname}\PYG{p}{]}\PYG{p}{)}
    \PYG{k}{if} \PYG{n}{subset} \PYG{o+ow}{is} \PYG{o+ow}{not} \PYG{k+kc}{None}\PYG{p}{:}
        \PYG{k}{if} \PYG{n}{subset} \PYG{o+ow}{not} \PYG{o+ow}{in} \PYG{p}{[}\PYG{l+s+s2}{\PYGZdq{}}\PYG{l+s+s2}{shortsus}\PYG{l+s+s2}{\PYGZdq{}}\PYG{p}{,} \PYG{l+s+s2}{\PYGZdq{}}\PYG{l+s+s2}{longsus}\PYG{l+s+s2}{\PYGZdq{}}\PYG{p}{,} \PYG{l+s+s2}{\PYGZdq{}}\PYG{l+s+s2}{nonquote}\PYG{l+s+s2}{\PYGZdq{}}\PYG{p}{,} \PYG{l+s+s2}{\PYGZdq{}}\PYG{l+s+s2}{quote}\PYG{l+s+s2}{\PYGZdq{}}\PYG{p}{]}\PYG{p}{:}
            \PYG{k}{raise} \PYG{n+ne}{ValueError}\PYG{p}{(}\PYG{l+s+s1}{\PYGZsq{}}\PYG{l+s+s1}{bad subset parameter: }\PYG{l+s+s1}{\PYGZdq{}}\PYG{l+s+si}{\PYGZpc{}s}\PYG{l+s+s1}{\PYGZdq{}}\PYG{l+s+s1}{\PYGZsq{}} \PYG{o}{\PYGZpc{}} \PYG{n}{subset}\PYG{p}{)}
        \PYG{n}{query} \PYG{o}{=} \PYG{n}{query} \PYG{o}{+} \PYG{l+s+s2}{\PYGZdq{}}\PYG{l+s+s2}{\PYGZam{}subset=}\PYG{l+s+si}{\PYGZpc{}s}\PYG{l+s+s2}{\PYGZdq{}} \PYG{o}{\PYGZpc{}} \PYG{n}{subset}
    \PYG{n}{query} \PYG{o}{=} \PYG{n}{query} \PYG{o}{+} \PYG{l+s+s2}{\PYGZdq{}}\PYG{l+s+s2}{\PYGZam{}clusterlength=}\PYG{l+s+si}{\PYGZpc{}d}\PYG{l+s+s2}{\PYGZam{}cutoff=}\PYG{l+s+si}{\PYGZpc{}d}\PYG{l+s+s2}{\PYGZdq{}} \PYG{o}{\PYGZpc{}} \PYG{p}{(}\PYG{n}{length}\PYG{p}{,} \PYG{n}{cutoff}\PYG{p}{)}
    \PYG{n}{rv} \PYG{o}{=} \PYG{n}{api\PYGZus{}request}\PYG{p}{(}\PYG{n}{endpoint}\PYG{o}{=}\PYG{l+s+s2}{\PYGZdq{}}\PYG{l+s+s2}{cluster}\PYG{l+s+s2}{\PYGZdq{}}\PYG{p}{,} \PYG{n}{query}\PYG{o}{=}\PYG{n}{query}\PYG{p}{)}
    \PYG{n}{clusters} \PYG{o}{=} \PYG{n}{OrderedDict}\PYG{p}{(}\PYG{n+nb}{sorted}\PYG{p}{(}\PYG{n}{rv}\PYG{p}{[}\PYG{l+s+s1}{\PYGZsq{}}\PYG{l+s+s1}{data}\PYG{l+s+s1}{\PYGZsq{}}\PYG{p}{]}\PYG{p}{,} \PYG{n}{key}\PYG{o}{=}\PYG{n}{itemgetter}\PYG{p}{(}\PYG{l+m+mi}{1}\PYG{p}{)}\PYG{p}{,} \PYG{n}{reverse}\PYG{o}{=}\PYG{k+kc}{True}\PYG{p}{)}\PYG{p}{)}
    \PYG{k}{return} \PYG{n}{clusters}
\end{sphinxVerbatim}

\sphinxAtStartPar
Find out what texts are available:

\begin{sphinxVerbatim}[commandchars=\\\{\}]
\PYG{g+gp}{\PYGZgt{}\PYGZgt{}\PYGZgt{} }\PYG{n}{lookup} \PYG{o}{=} \PYG{n}{get\PYGZus{}lookup}\PYG{p}{(}\PYG{p}{)}
\PYG{g+gp}{\PYGZgt{}\PYGZgt{}\PYGZgt{} }\PYG{n}{lookup}\PYG{o}{.}\PYG{n}{head}\PYG{p}{(}\PYG{p}{)}
\PYG{g+go}{   corpus             author shortname                       title}
\PYG{g+go}{0  ChiLit   Agnes Strickland     rival           The Rival Crusoes}
\PYG{g+go}{1  ChiLit        Andrew Lang    prigio               Prince Prigio}
\PYG{g+go}{2  ChiLit  Ann Fraser Tytler     leila               Leila at Home}
\PYG{g+go}{3  ChiLit        Anna Sewell    beauty                Black Beauty}
\PYG{g+go}{4  ChiLit     Beatrix Potter     bunny  The Tale Of Benjamin Bunny}
\PYG{g+gp}{\PYGZgt{}\PYGZgt{}\PYGZgt{} }\PYG{n}{lookup}\PYG{o}{.}\PYG{n}{tail}\PYG{p}{(}\PYG{p}{)}
\PYG{g+go}{    corpus                       author shortname                          title}
\PYG{g+go}{133    ntc                 Thomas Hardy    native       The Return of the Native}
\PYG{g+go}{134    ntc               Wilkie Collins    Antoni  Antonina, or the Fall of Rome}
\PYG{g+go}{135    ntc               Wilkie Collins      arma                       Armadale}
\PYG{g+go}{136    ntc               Wilkie Collins    wwhite             The Woman in White}
\PYG{g+go}{137    ntc  William Makepeace Thackeray    vanity                    Vanity Fair}
\end{sphinxVerbatim}

\sphinxAtStartPar
Filter what is available:

\begin{sphinxVerbatim}[commandchars=\\\{\}]
\PYG{g+gp}{\PYGZgt{}\PYGZgt{}\PYGZgt{} }\PYG{n}{lookup}\PYG{p}{[}\PYG{n}{lookup}\PYG{p}{[}\PYG{l+s+s1}{\PYGZsq{}}\PYG{l+s+s1}{author}\PYG{l+s+s1}{\PYGZsq{}}\PYG{p}{]} \PYG{o}{==} \PYG{l+s+s2}{\PYGZdq{}}\PYG{l+s+s2}{Thomas Hardy}\PYG{l+s+s2}{\PYGZdq{}}\PYG{p}{]}
\PYG{g+go}{    corpus        author shortname                      title}
\PYG{g+go}{131    ntc  Thomas Hardy      Jude           Jude the Obscure}
\PYG{g+go}{132    ntc  Thomas Hardy      Tess  Tess of the D\PYGZsq{}Urbervilles}
\PYG{g+go}{133    ntc  Thomas Hardy    native   The Return of the Native}
\end{sphinxVerbatim}

\sphinxAtStartPar
Fetch the tokens for a specific text:

\begin{sphinxVerbatim}[commandchars=\\\{\}]
\PYG{g+gp}{\PYGZgt{}\PYGZgt{}\PYGZgt{} }\PYG{n}{tokens} \PYG{o}{=} \PYG{n}{get\PYGZus{}tokens}\PYG{p}{(}\PYG{n}{shortname}\PYG{o}{=}\PYG{l+s+s1}{\PYGZsq{}}\PYG{l+s+s1}{leila}\PYG{l+s+s1}{\PYGZsq{}}\PYG{p}{)}
\PYG{g+gp}{\PYGZgt{}\PYGZgt{}\PYGZgt{} }\PYG{n+nb}{len}\PYG{p}{(}\PYG{n}{tokens}\PYG{p}{)}
\PYG{g+go}{63026}
\PYG{g+gp}{\PYGZgt{}\PYGZgt{}\PYGZgt{} }\PYG{n}{tokens}\PYG{p}{[}\PYG{l+m+mi}{0}\PYG{p}{:}\PYG{l+m+mi}{9}\PYG{p}{]}
\PYG{g+go}{[\PYGZsq{}it\PYGZsq{}, \PYGZsq{}was\PYGZsq{}, \PYGZsq{}the\PYGZsq{}, \PYGZsq{}intention\PYGZsq{}, \PYGZsq{}of\PYGZsq{}, \PYGZsq{}the\PYGZsq{}, \PYGZsq{}writer\PYGZsq{}, \PYGZsq{}of\PYGZsq{}, \PYGZsq{}the\PYGZsq{}]}
\end{sphinxVerbatim}

\sphinxAtStartPar
Fetch the tokens for all quotes text in novels by Jane Austen:

\begin{sphinxVerbatim}[commandchars=\\\{\}]
\PYG{g+gp}{\PYGZgt{}\PYGZgt{}\PYGZgt{} }\PYG{n}{wanted} \PYG{o}{=} \PYG{p}{[}\PYG{n}{sn} \PYG{k}{for} \PYG{n}{sn} \PYG{o+ow}{in} \PYG{n}{lookup}\PYG{p}{[}\PYG{n}{lookup}\PYG{p}{[}\PYG{l+s+s1}{\PYGZsq{}}\PYG{l+s+s1}{author}\PYG{l+s+s1}{\PYGZsq{}}\PYG{p}{]} \PYG{o}{==} \PYG{l+s+s2}{\PYGZdq{}}\PYG{l+s+s2}{Jane Austen}\PYG{l+s+s2}{\PYGZdq{}}\PYG{p}{]}\PYG{p}{[}\PYG{l+s+s1}{\PYGZsq{}}\PYG{l+s+s1}{shortname}\PYG{l+s+s1}{\PYGZsq{}}\PYG{p}{]}\PYG{p}{]}
\PYG{g+gp}{\PYGZgt{}\PYGZgt{}\PYGZgt{} }\PYG{n}{wanted}
\PYG{g+go}{[\PYGZsq{}ladysusan\PYGZsq{}, \PYGZsq{}mansfield\PYGZsq{}, \PYGZsq{}northanger\PYGZsq{}, \PYGZsq{}sense\PYGZsq{}, \PYGZsq{}emma\PYGZsq{}, \PYGZsq{}persuasion\PYGZsq{}, \PYGZsq{}pride\PYGZsq{}]}

\PYG{g+gp}{\PYGZgt{}\PYGZgt{}\PYGZgt{} }\PYG{n}{austen\PYGZus{}quotes} \PYG{o}{=} \PYG{n}{get\PYGZus{}tokens}\PYG{p}{(}\PYG{n}{shortname}\PYG{o}{=}\PYG{n}{wanted}\PYG{p}{,} \PYG{n}{subset}\PYG{o}{=}\PYG{l+s+s2}{\PYGZdq{}}\PYG{l+s+s2}{quote}\PYG{l+s+s2}{\PYGZdq{}}\PYG{p}{)}
\PYG{g+gp}{\PYGZgt{}\PYGZgt{}\PYGZgt{} }\PYG{n+nb}{len}\PYG{p}{(}\PYG{n}{austen\PYGZus{}quotes}\PYG{p}{)}
\PYG{g+go}{307445}
\PYG{g+gp}{\PYGZgt{}\PYGZgt{}\PYGZgt{} }\PYG{n}{austen\PYGZus{}quotes}\PYG{p}{[}\PYG{l+m+mi}{0}\PYG{p}{:}\PYG{l+m+mi}{9}\PYG{p}{]}
\PYG{g+go}{[\PYGZsq{}poor\PYGZsq{}, \PYGZsq{}miss\PYGZsq{}, \PYGZsq{}taylor\PYGZsq{}, \PYGZsq{}i\PYGZsq{}, \PYGZsq{}wish\PYGZsq{}, \PYGZsq{}she\PYGZsq{}, \PYGZsq{}were\PYGZsq{}, \PYGZsq{}here\PYGZsq{}, \PYGZsq{}again\PYGZsq{}]}
\end{sphinxVerbatim}

\sphinxAtStartPar
Keep each text separate:

\begin{sphinxVerbatim}[commandchars=\\\{\}]
\PYG{g+gp}{\PYGZgt{}\PYGZgt{}\PYGZgt{} }\PYG{n}{austen\PYGZus{}quotes} \PYG{o}{=} \PYG{p}{\PYGZob{}}\PYG{p}{\PYGZcb{}}
\PYG{g+gp}{\PYGZgt{}\PYGZgt{}\PYGZgt{} }\PYG{k}{for} \PYG{n}{sn} \PYG{o+ow}{in} \PYG{n}{wanted}\PYG{p}{:}
\PYG{g+gp}{... }    \PYG{n}{austen\PYGZus{}quotes}\PYG{p}{[}\PYG{n}{sn}\PYG{p}{]} \PYG{o}{=} \PYG{n}{get\PYGZus{}tokens}\PYG{p}{(}\PYG{n}{shortname}\PYG{o}{=}\PYG{n}{sn}\PYG{p}{,} \PYG{n}{subset}\PYG{o}{=}\PYG{l+s+s2}{\PYGZdq{}}\PYG{l+s+s2}{quote}\PYG{l+s+s2}{\PYGZdq{}}\PYG{p}{)}
\PYG{g+gp}{...}
\PYG{g+gp}{\PYGZgt{}\PYGZgt{}\PYGZgt{} }\PYG{p}{\PYGZob{}}\PYG{n}{key}\PYG{p}{:}\PYG{n+nb}{len}\PYG{p}{(}\PYG{n}{value}\PYG{p}{)} \PYG{k}{for} \PYG{n}{key}\PYG{p}{,}\PYG{n}{value} \PYG{o+ow}{in} \PYG{n}{austen\PYGZus{}quotes}\PYG{o}{.}\PYG{n}{items}\PYG{p}{(}\PYG{p}{)}\PYG{p}{\PYGZcb{}}
\PYG{g+gp}{\PYGZgt{}\PYGZgt{}\PYGZgt{} }\PYG{n+nb}{print}\PYG{p}{(}\PYG{n}{json}\PYG{o}{.}\PYG{n}{dumps}\PYG{p}{(}\PYG{p}{\PYGZob{}}\PYG{n}{key}\PYG{p}{:}\PYG{n+nb}{len}\PYG{p}{(}\PYG{n}{value}\PYG{p}{)} \PYG{k}{for} \PYG{n}{key}\PYG{p}{,}\PYG{n}{value} \PYG{o+ow}{in} \PYG{n}{austen\PYGZus{}quotes}\PYG{o}{.}\PYG{n}{items}\PYG{p}{(}\PYG{p}{)}\PYG{p}{\PYGZcb{}}\PYG{p}{)}\PYG{p}{)}
\PYG{g+go}{\PYGZob{}}
\PYG{g+go}{  \PYGZdq{}ladysusan\PYGZdq{}: 2791,}
\PYG{g+go}{  \PYGZdq{}mansfield\PYGZdq{}: 62013,}
\PYG{g+go}{  \PYGZdq{}northanger\PYGZdq{}: 28937,}
\PYG{g+go}{  \PYGZdq{}sense\PYGZdq{}: 51744,}
\PYG{g+go}{  \PYGZdq{}emma\PYGZdq{}: 80319,}
\PYG{g+go}{  \PYGZdq{}persuasion\PYGZdq{}: 28653,}
\PYG{g+go}{  \PYGZdq{}pride\PYGZdq{}: 52988}
\PYG{g+go}{\PYGZcb{}}
\PYG{g+gp}{\PYGZgt{}\PYGZgt{}\PYGZgt{} }\PYG{n}{austen\PYGZus{}quotes}\PYG{p}{[}\PYG{l+s+s1}{\PYGZsq{}}\PYG{l+s+s1}{emma}\PYG{l+s+s1}{\PYGZsq{}}\PYG{p}{]}\PYG{p}{[}\PYG{l+m+mi}{0}\PYG{p}{:}\PYG{l+m+mi}{9}\PYG{p}{]}
\PYG{g+go}{[\PYGZsq{}poor\PYGZsq{}, \PYGZsq{}miss\PYGZsq{}, \PYGZsq{}taylor\PYGZsq{}, \PYGZsq{}i\PYGZsq{}, \PYGZsq{}wish\PYGZsq{}, \PYGZsq{}she\PYGZsq{}, \PYGZsq{}were\PYGZsq{}, \PYGZsq{}here\PYGZsq{}, \PYGZsq{}again\PYGZsq{}]}
\end{sphinxVerbatim}

\sphinxAtStartPar
And now lets get some clusters for the Jane Austen novels:

\begin{sphinxVerbatim}[commandchars=\\\{\}]
\PYG{g+gp}{\PYGZgt{}\PYGZgt{}\PYGZgt{} }\PYG{n}{austen\PYGZus{}clusters} \PYG{o}{=} \PYG{n}{get\PYGZus{}clusters}\PYG{p}{(}\PYG{n}{shortname}\PYG{o}{=}\PYG{n}{wanted}\PYG{p}{,} \PYG{n}{length}\PYG{o}{=}\PYG{l+m+mi}{5}\PYG{p}{,} \PYG{n}{cutoff}\PYG{o}{=}\PYG{l+m+mi}{5}\PYG{p}{,} \PYG{n}{subset}\PYG{o}{=}\PYG{l+s+s2}{\PYGZdq{}}\PYG{l+s+s2}{quote}\PYG{l+s+s2}{\PYGZdq{}}\PYG{p}{)}
\PYG{g+gp}{\PYGZgt{}\PYGZgt{}\PYGZgt{} }\PYG{n+nb}{print}\PYG{p}{(}\PYG{n}{json}\PYG{o}{.}\PYG{n}{dumps}\PYG{p}{(}\PYG{n}{austen\PYGZus{}clusters}\PYG{p}{,} \PYG{n}{indent}\PYG{o}{=}\PYG{l+m+mi}{2}\PYG{p}{)}\PYG{p}{)}
\PYG{g+go}{\PYGZob{}}
\PYG{g+go}{  \PYGZdq{}i do not know what\PYGZdq{}: 26,}
\PYG{g+go}{  \PYGZdq{}i am sure you will\PYGZdq{}: 16,}
\PYG{g+go}{  \PYGZdq{}i do not know that\PYGZdq{}: 16,}
\PYG{g+go}{  \PYGZdq{}i do not mean to\PYGZdq{}: 16,}
\PYG{g+go}{  \PYGZdq{}and i am sure i\PYGZdq{}: 16,}
\PYG{g+go}{  \PYGZdq{}i have no doubt of\PYGZdq{}: 14,}
\PYG{g+go}{  \PYGZdq{}i do not think i\PYGZdq{}: 14,}
\PYG{g+go}{  \PYGZdq{}i am sure i should\PYGZdq{}: 13,}
\PYG{g+go}{  \PYGZdq{}i am sure i do\PYGZdq{}: 11,}
\PYG{g+go}{  \PYGZdq{}i do not pretend to\PYGZdq{}: 11,}
\PYG{g+go}{  ...}
\end{sphinxVerbatim}


\subsection{R}
\label{\detokenize{advanced/api_usage:r}}
\sphinxAtStartPar
Functions to access the CLiC API from R are available in the
\sphinxhref{https://github.com/mahlberg-lab/clicclient}{clicclient R package}.
This package is still under development to work with CLiC 2.0. If you
would like to contribute to this package, please get in touch at
\sphinxhref{mailto:dhss-kontakt@fau.de}{dhss\sphinxhyphen{}kontakt@fau.de}.

\sphinxstepscope


\chapter{Uploading new texts}
\label{\detokenize{advanced/adding_corpora:uploading-new-texts}}\label{\detokenize{advanced/adding_corpora:adding-corpora}}\label{\detokenize{advanced/adding_corpora::doc}}
\sphinxAtStartPar
The raw texts for CLiC to process are assumed to be in the corpora
repository:
\begin{quote}

\sphinxAtStartPar
\sphinxurl{https://github.com/mahlberg-lab/corpora}
\end{quote}

\sphinxAtStartPar
If you haven’t done so already, add the texts to this repository,
following the instructions there.

\sphinxAtStartPar
For each new book, first check the region annotation using, for
example:

\begin{sphinxVerbatim}[commandchars=\\\{\}]
\PYG{o}{.}\PYG{o}{/}\PYG{n}{server}\PYG{o}{/}\PYG{n+nb}{bin}\PYG{o}{/}\PYG{n}{region\PYGZus{}preview} \PYG{n}{corpora}\PYG{o}{/}\PYG{n}{ChiLit}\PYG{o}{/}\PYG{n}{alone}\PYG{o}{.}\PYG{n}{txt}
\end{sphinxVerbatim}

\sphinxAtStartPar
This will show the text file with all regions marked up with different
colours. Inspect the output to see if regions are marked up correctly by
default. If the scripts have made mistakes, you have 2 options:
\begin{itemize}
\item {} 
\sphinxAtStartPar
Fix any systematic problems in the region taggers in
\sphinxcode{\sphinxupquote{server/clic/region}}

\item {} 
\sphinxAtStartPar
Export the tagged regions to a \sphinxcode{\sphinxupquote{*.regions.csv}} file using
\sphinxcode{\sphinxupquote{server/bin/region\_export}}, any regions edited here will override
the tagger.

\end{itemize}

\sphinxAtStartPar
Once you are happy with the output, it needs to be imported into the
CLiC database. Do not perform the following on a live CLiC server, as
CLiC will not be available during the process. Instead, follow the
instructions on dumping / restoring a database in
{\hyperref[\detokenize{advanced/production:managing-a-production-instance-of-clic}]{\sphinxcrossref{\DUrole{std,std-ref}{Managing a production instance of CLiC}}}}.

\sphinxAtStartPar
\sphinxcode{\sphinxupquote{import\_corpora\_repo}} will import all given texts into the database.
Ideally (re\sphinxhyphen{})import the entire corpora repository with:

\begin{sphinxVerbatim}[commandchars=\\\{\}]
\PYG{o}{.}\PYG{o}{/}\PYG{n}{server}\PYG{o}{/}\PYG{n+nb}{bin}\PYG{o}{/}\PYG{n}{import\PYGZus{}corpora\PYGZus{}repo} \PYG{n}{corpora}\PYG{o}{/}\PYG{o}{*}\PYG{o}{/}\PYG{o}{*}\PYG{o}{.}\PYG{n}{txt}
\end{sphinxVerbatim}

\sphinxAtStartPar
…however this takes some time. It is also possible to import
individual books:

\begin{sphinxVerbatim}[commandchars=\\\{\}]
\PYG{o}{.}\PYG{o}{/}\PYG{n}{server}\PYG{o}{/}\PYG{n+nb}{bin}\PYG{o}{/}\PYG{n}{import\PYGZus{}corpora\PYGZus{}repo} \PYG{n}{corpora}\PYG{o}{/}\PYG{n}{ChiLit}\PYG{o}{/}\PYG{n}{alone}\PYG{o}{.}\PYG{n}{txt}
\end{sphinxVerbatim}

\sphinxAtStartPar
Restart CLiC with \sphinxcode{\sphinxupquote{./install.sh}} to ensure that caches are flushed.


\section{What if my book is not appearing in CLiC?}
\label{\detokenize{advanced/adding_corpora:what-if-my-book-is-not-appearing-in-clic}}
\sphinxAtStartPar
First, make sure that the \sphinxcode{\sphinxupquote{corpora.bib}} file has an entry for your new
book. Without an entry here, CLiC will not know to add it to the corpus,
and thus it will not appear in the drop\sphinxhyphen{}down. Re\sphinxhyphen{}import the book once
this has been done.

\sphinxstepscope


\chapter{Developing the system}
\label{\detokenize{advanced/development:developing-the-system}}\label{\detokenize{advanced/development::doc}}
\sphinxAtStartPar
Set\sphinxhyphen{}up of the development environment is very similar to a production
environment, see \sphinxurl{https://github.com/mahlberg-lab/clic/blob/HEAD/README.rst} for details on how to do this.

\sphinxAtStartPar
The main difference is to use \sphinxcode{\sphinxupquote{PROJECT\_MODE= development}} (or not include any
PROJECT\_MODE, development is default), in your \sphinxcode{\sphinxupquote{local\sphinxhyphen{}conf.mk}}. This will stop
CLiC from being started automatically so you can run the server in debug mode.
It will also disable NGINX’s server cache so you won’t get stale responses.

\sphinxAtStartPar
To start the API server in debug mode:

\begin{sphinxVerbatim}[commandchars=\\\{\}]
\PYG{n}{make} \PYG{n}{start}
\end{sphinxVerbatim}

\sphinxAtStartPar
To log database queries that the API makes:

\begin{sphinxVerbatim}[commandchars=\\\{\}]
\PYG{n}{QUERY\PYGZus{}LOG}\PYG{o}{=}\PYG{n}{yes} \PYG{n}{make} \PYG{n}{start}
\PYG{n}{QUERY\PYGZus{}LOG}\PYG{o}{=}\PYG{n}{explain} \PYG{n}{make} \PYG{n}{start}  \PYG{c+c1}{\PYGZsh{} Also include explain plan}
\end{sphinxVerbatim}


\section{Client\sphinxhyphen{}side development}
\label{\detokenize{advanced/development:client-side-development}}
\sphinxAtStartPar
To run tests, lint code and compile:

\begin{sphinxVerbatim}[commandchars=\\\{\}]
\PYG{n}{make} \PYG{o}{\PYGZhy{}}\PYG{n}{C} \PYG{n}{client}
\end{sphinxVerbatim}

\sphinxAtStartPar
…you will need to do this after any change to the HTML/CSS/JS source.


\section{Server\sphinxhyphen{}side development}
\label{\detokenize{advanced/development:server-side-development}}
\sphinxAtStartPar
To run unit tests:

\begin{sphinxVerbatim}[commandchars=\\\{\}]
\PYG{n}{make} \PYG{o}{\PYGZhy{}}\PYG{n}{C} \PYG{n}{server} \PYG{n}{test}
\end{sphinxVerbatim}


\section{Coverage reports}
\label{\detokenize{advanced/development:coverage-reports}}
\sphinxAtStartPar
Coverage reports can be built for both client and server:

\begin{sphinxVerbatim}[commandchars=\\\{\}]
\PYG{n}{make} \PYG{n}{coverage}
\end{sphinxVerbatim}

\sphinxAtStartPar
…the reports will be available through the reports at the end.


\section{Branching and publishing to live}
\label{\detokenize{advanced/development:branching-and-publishing-to-live}}
\sphinxAtStartPar
The CLiC project has several long\sphinxhyphen{}lived branches:
\begin{description}
\sphinxlineitem{test}
\sphinxAtStartPar
The default branch, the test server will run this version of the CLiC code.
Documentation updates should be done directly in this branch.
Code updates should be done in a feature branch and merged with a pull request.

\sphinxlineitem{2.x (e.g. \sphinxcode{\sphinxupquote{2.0}}, \sphinxcode{\sphinxupquote{2.1}}, …)}
\sphinxAtStartPar
Production branch for major versions of CLiC, the production server will be running one of these branches.
Generally, new releases should be done with a pull request from the \sphinxcode{\sphinxupquote{development}} branch.
Any bug\sphinxhyphen{}fix changes should be done with a pull request directly to the \sphinxcode{\sphinxupquote{2.x}} branch.

\end{description}


\section{Exercise API integration tests}
\label{\detokenize{advanced/development:exercise-api-integration-tests}}
\sphinxAtStartPar
This repository contains some canned API calls and output that can be run against
any server. For example:

\begin{sphinxVerbatim}[commandchars=\\\{\}]
\PYG{o}{.}\PYG{o}{/}\PYG{n}{exercise}\PYG{o}{.}\PYG{n}{sh} \PYG{n}{http}\PYG{p}{:}\PYG{o}{/}\PYG{o}{/}\PYG{n}{cal}\PYG{o}{\PYGZhy{}}\PYG{n}{n}\PYG{o}{\PYGZhy{}}\PYG{n}{clic}\PYG{o}{\PYGZhy{}}\PYG{l+m+mf}{01.}\PYG{n}{bham}\PYG{o}{.}\PYG{n}{ac}\PYG{o}{.}\PYG{n}{uk} \PYG{n}{exercise\PYGZus{}outputs}\PYG{o}{/}\PYG{o}{*}
\end{sphinxVerbatim}


\section{Favico regeneration}
\label{\detokenize{advanced/development:favico-regeneration}}
\sphinxAtStartPar
Upload \sphinxcode{\sphinxupquote{./client/www/logos/clic\_icon.svg}} to \sphinxurl{http://cthedot.de/icongen/}, and place the results into
\sphinxcode{\sphinxupquote{\textasciigrave{}client/www/index.html}} and \sphinxcode{\sphinxupquote{client/www/iconx}} as appropriate.


\section{Preparing a release}
\label{\detokenize{advanced/development:preparing-a-release}}
\sphinxAtStartPar
Steps for making a release:

\begin{sphinxVerbatim}[commandchars=\\\{\}]
VERSION=\PYGZdq{}2.0.0\PYGZhy{}beta5\PYGZdq{}
echo \PYGZhy{}e \PYGZdq{}\PYGZsh{}\PYGZsh{} \PYGZdl{}\PYGZob{}VERSION\PYGZcb{} (\PYGZdl{}(date +\PYGZsq{}\PYGZpc{}Y\PYGZhy{}\PYGZpc{}m\PYGZhy{}\PYGZpc{}d\PYGZsq{}))\PYGZbs{}n\PYGZdq{} | cat \PYGZhy{} CHANGELOG.md \PYGZgt{} CHANGELOG.md.n
mv CHANGELOG.md.n CHANGELOG.md
git commit \PYGZhy{}m \PYGZdq{}CHANGELOG: Release \PYGZdl{}\PYGZob{}VERSION\PYGZcb{}\PYGZdq{} CHANGELOG.md
git tag \PYGZhy{}am \PYGZdq{}Release \PYGZdl{}\PYGZob{}VERSION\PYGZcb{}\PYGZdq{} v\PYGZdl{}\PYGZob{}VERSION\PYGZcb{} HEAD
git push \PYGZam{}\PYGZam{} git push \PYGZhy{}\PYGZhy{}tags
\end{sphinxVerbatim}


\section{Documentation}
\label{\detokenize{advanced/development:documentation}}
\sphinxAtStartPar
readthedocs is used for our documentation. In order for newly released versions of CLiC (as managed by tags, e.g. v2.1.2) to be included in readthedocs please make sure to follow the above instructions in ‘Preparing a release’ when updating CLiC on the production server.

\sphinxAtStartPar
You can access the settings for CLiC’s readthedocs by going to: \sphinxurl{https://readthedocs.org/projects/clic/}

\sphinxAtStartPar
On this page, there should be a line for each CLiC branch / major.minor version, and ‘latest’. When a new branch is started in CLiC the corresponding branch needs to be added to readthedocs, otherwise there will be a 404 error when selecting the “help” link. If a new branch is missing, go to the above page and…

\sphinxAtStartPar
\# “Add version”
\# Search for the branch number in the box
\# Turn on “Active”
\# Click “Update version”

\sphinxAtStartPar
You should also change the default branch by…

\sphinxAtStartPar
\# “Settings” in top\sphinxhyphen{}right corner
\# Select new default branch in drop\sphinxhyphen{}down

\sphinxAtStartPar
If you need maintainer access, please request from an existing maintainer.

\sphinxAtStartPar
Sphinx will also be run locally for testing, and if it proves useful in future.

\sphinxAtStartPar
To re\sphinxhyphen{}build documentation:

\begin{sphinxVerbatim}[commandchars=\\\{\}]
\PYG{n}{make} \PYG{o}{\PYGZhy{}}\PYG{n}{C} \PYG{n}{docs}
\end{sphinxVerbatim}

\sphinxAtStartPar
To view the documenation built, go to \sphinxcode{\sphinxupquote{/local\sphinxhyphen{}docs/}} for your instance.

\sphinxAtStartPar
By default, this only generates the LaTeX needed for the PDF output, which will
be available in \sphinxcode{\sphinxupquote{docs/\_build/latex/clic\sphinxhyphen{}userguide.tex}}. To build a PDF from
this, you first need to install:

\begin{sphinxVerbatim}[commandchars=\\\{\}]
\PYG{n}{apt} \PYG{n}{install} \PYG{n}{texlive}\PYG{o}{\PYGZhy{}}\PYG{n}{latex}\PYG{o}{\PYGZhy{}}\PYG{n}{recommended} \PYGZbs{}
            \PYG{n}{texlive}\PYG{o}{\PYGZhy{}}\PYG{n}{fonts}\PYG{o}{\PYGZhy{}}\PYG{n}{recommended} \PYGZbs{}
            \PYG{n}{texlive}\PYG{o}{\PYGZhy{}}\PYG{n}{latex}\PYG{o}{\PYGZhy{}}\PYG{n}{extra} \PYGZbs{}
            \PYG{n}{latexmk}
\end{sphinxVerbatim}

\sphinxAtStartPar
…then you can build with:

\begin{sphinxVerbatim}[commandchars=\\\{\}]
\PYG{n}{make} \PYG{o}{\PYGZhy{}}\PYG{n}{C} \PYG{n}{docs} \PYG{n}{pdf}
\end{sphinxVerbatim}

\sphinxAtStartPar
The pdf will be available at \sphinxcode{\sphinxupquote{/local\sphinxhyphen{}docs/latex/clic\sphinxhyphen{}userguide.pdf}}.


\section{FlexiConc: Upgrading \& dependencies}
\label{\detokenize{advanced/development:flexiconc-upgrading-dependencies}}
\sphinxAtStartPar
FlexiConc is run client\sphinxhyphen{}side within a pyodide environment.

\sphinxAtStartPar
To upgrade the version used, or to add dependencies for a package, see \sphinxcode{\sphinxupquote{flexiclic/Makefile}}.

\sphinxstepscope


\chapter{Managing a production instance of CLiC}
\label{\detokenize{advanced/production:managing-a-production-instance-of-clic}}\label{\detokenize{advanced/production::doc}}
\sphinxAtStartPar
Installation instructions are all in \sphinxurl{https://github.com/mahlberg-lab/clic/blob/HEAD/README.rst}.


\section{Cache warming}
\label{\detokenize{advanced/production:cache-warming}}
\sphinxAtStartPar
Repeatedly used CLiC API calls will be cached by NGINX to speed up CLiC. There
is a script to ensure that the most commonly\sphinxhyphen{}used calls are in the cache. For
example:

\begin{sphinxVerbatim}[commandchars=\\\{\}]
\PYG{o}{.}\PYG{o}{/}\PYG{n}{cache}\PYG{o}{\PYGZhy{}}\PYG{n}{warm}\PYG{o}{.}\PYG{n}{sh} \PYG{n}{https}\PYG{p}{:}\PYG{o}{/}\PYG{o}{/}\PYG{n}{clic}\PYG{o}{.}\PYG{n}{bham}\PYG{o}{.}\PYG{n}{ac}\PYG{o}{.}\PYG{n}{uk}
\end{sphinxVerbatim}


\section{Start / stop the CLiC service}
\label{\detokenize{advanced/production:start-stop-the-clic-service}}
\sphinxAtStartPar
If you need to stop/start CLiC outside this for whatever reason, use systemctl,
e.g. \sphinxcode{\sphinxupquote{systemctl stop clic}}.

\sphinxAtStartPar
To see the current status, use \sphinxcode{\sphinxupquote{systemctl status clic}}.


\section{Upgrading a CLiC instance}
\label{\detokenize{advanced/production:upgrading-a-clic-instance}}
\sphinxAtStartPar
To upgrade a CLiC instance, first check out the code for the given release, e.g.:

\begin{sphinxVerbatim}[commandchars=\\\{\}]
\PYG{n}{git} \PYG{n}{fetch} \PYG{o}{\PYGZam{}}\PYG{o}{\PYGZam{}} \PYG{n}{git} \PYG{n}{checkout} \PYG{n}{v1}\PYG{l+m+mf}{.7}\PYG{l+m+mf}{.0}
\end{sphinxVerbatim}

\sphinxAtStartPar
…then run:

\begin{sphinxVerbatim}[commandchars=\\\{\}]
\PYG{n}{make}
\PYG{n}{sudo} \PYG{o}{.}\PYG{o}{/}\PYG{n}{install}\PYG{o}{.}\PYG{n}{sh}
\end{sphinxVerbatim}


\section{Troubleshooting}
\label{\detokenize{advanced/production:troubleshooting}}
\sphinxAtStartPar
Logs are available with: \sphinxcode{\sphinxupquote{journalctl \sphinxhyphen{}S\sphinxhyphen{}1d \sphinxhyphen{}uclic}}.

\sphinxAtStartPar
If you see the NGINX default “Welcome to nginx!” page when trying to use CLiC from a web browser:
\begin{itemize}
\item {} 
\sphinxAtStartPar
Make sure you used a \sphinxcode{\sphinxupquote{WWW\_SERVER\_NAME}} that resolves to the CLiC server in \sphinxcode{\sphinxupquote{local\sphinxhyphen{}conf.mk}}. Re\sphinxhyphen{}run \sphinxcode{\sphinxupquote{make}} and \sphinxcode{\sphinxupquote{./install.sh}}.

\item {} 
\sphinxAtStartPar
Make sure NGINX started without errors: \sphinxcode{\sphinxupquote{systemctl status \sphinxhyphen{}ln50 nginx}}

\end{itemize}

\sphinxAtStartPar
If any CLiC query responds with the error “The CLiC server did not respond, the query may have taken too long.”:
\begin{itemize}
\item {} 
\sphinxAtStartPar
Make sure that CLiC is started with \sphinxcode{\sphinxupquote{systemctl start clic}}.

\item {} 
\sphinxAtStartPar
Any error messages are cached by NGINX and your web browser briefly, clear your cache and try again. You can also re\sphinxhyphen{}run \sphinxcode{\sphinxupquote{./install.sh}} to clear the NGINX cache.

\end{itemize}

\sphinxAtStartPar
If you see the “CLiC is down for maintenance” page:
\begin{itemize}
\item {} 
\sphinxAtStartPar
Make sure CLiC has started without errors: \sphinxcode{\sphinxupquote{systemctl status \sphinxhyphen{}ln50 clic}}

\end{itemize}

\sphinxAtStartPar
If you see errors about missing tables, or queries are particularly slow:
\begin{itemize}
\item {} 
\sphinxAtStartPar
Postgres may need vacuuming. Run \sphinxcode{\sphinxupquote{sudo \sphinxhyphen{}upostgres psql \sphinxhyphen{}c "VACUUM ANALYSE" clic\_db}}.

\end{itemize}

\sphinxAtStartPar
If the home page is particularly slow:
\begin{itemize}
\item {} 
\sphinxAtStartPar
The homepage entries aren’t cached yet, reload the page \textasciitilde{}3 times until it is.

\end{itemize}

\sphinxAtStartPar
If the “help” link is heading to a 404 on readthedocs.io:
\begin{itemize}
\item {} 
\sphinxAtStartPar
The current branch is not generating documentation, go to \sphinxurl{https://readthedocs.org/projects/clic/} and diagnose (either adding a new branch for documentation, or diagnosing the existing one). See {\hyperref[\detokenize{advanced/development::doc}]{\sphinxcrossref{\DUrole{doc}{Developing the system}}}} for more information.

\end{itemize}


\section{Back\sphinxhyphen{}up / generating dumps}
\label{\detokenize{advanced/production:back-up-generating-dumps}}
\sphinxAtStartPar
An import process from the raw text files can take some time.
Instead, it can be more efficent to dump/restore a database from another CLiC
instance.

\sphinxAtStartPar
Dump the database with the following:

\begin{sphinxVerbatim}[commandchars=\\\{\}]
\PYG{o}{.}\PYG{o}{/}\PYG{n}{schema}\PYG{o}{/}\PYG{n+nb}{bin}\PYG{o}{/}\PYG{n}{db\PYGZus{}dump} \PYG{n}{clic}\PYG{o}{.}\PYG{n}{dump}
\end{sphinxVerbatim}

\sphinxAtStartPar
Restore with the following. NB: \sphinxstylestrong{this will destroy all existing data}:

\begin{sphinxVerbatim}[commandchars=\\\{\}]
\PYG{o}{.}\PYG{o}{/}\PYG{n}{schema}\PYG{o}{/}\PYG{n+nb}{bin}\PYG{o}{/}\PYG{n}{db\PYGZus{}restore} \PYG{n}{clic}\PYG{o}{.}\PYG{n}{dump}
\end{sphinxVerbatim}

\sphinxAtStartPar
If you do not want to destroy the existing database, you could set a new database name in \sphinxcode{\sphinxupquote{local\sphinxhyphen{}conf.mk}}:

\begin{sphinxVerbatim}[commandchars=\\\{\}]
\PYG{n}{echo} \PYG{l+s+s2}{\PYGZdq{}}\PYG{l+s+s2}{DB\PYGZus{}NAME = my\PYGZus{}new\PYGZus{}clic\PYGZus{}database}\PYG{l+s+s2}{\PYGZdq{}} \PYG{o}{\PYGZgt{}\PYGZgt{}} \PYG{n}{local}\PYG{o}{\PYGZhy{}}\PYG{n}{conf}\PYG{o}{.}\PYG{n}{mk}
\end{sphinxVerbatim}

\sphinxAtStartPar
…and then re\sphinxhyphen{}run \sphinxcode{\sphinxupquote{make}} before restoring.

\sphinxAtStartPar
There are server options to speed up the restore, see \sphinxhref{http://www.databasesoup.com/2014/09/settings-for-fast-pgrestore.html}{this blog post}.

\sphinxstepscope


\chapter{List of known issues}
\label{\detokenize{known-issues:list-of-known-issues}}\label{\detokenize{known-issues::doc}}
\sphinxAtStartPar
We are aware of the following issues and will try to fix them in future versions of CLiC.
If you notice anything else, please let us know via email at \sphinxhref{mailto:dhss-kontakt@fau.de}{dhss\sphinxhyphen{}kontakt@fau.de}
\begin{itemize}
\item {} 
\sphinxAtStartPar
\sphinxtitleref{LD} (Little Dorrit): the \sphinxhref{https://github.com/mahlberg-lab/corpora/blob/master/DNov/LD.txt}{text} is missing the title of its first part, which should read “BOOK THE FIRST POVERTY” (to match “BOOK THE SECOND RICHES” at the beginning of the second part \sphinxhyphen{} and both should be highlighted as parts)

\end{itemize}

\sphinxstepscope


\chapter{Appendices}
\label{\detokenize{appendices:appendices}}\label{\detokenize{appendices:a-list-of-texts-available-in-clic}}\label{\detokenize{appendices::doc}}
\sphinxAtStartPar
The appendices only act as an overview. For up\sphinxhyphen{}to\sphinxhyphen{}date information and
the documentation of how the corpus texts were cleaned, please refer to
the README in the \sphinxcite{references:github-corpora} repository. You can also create a
dynamic list of the books in the corpora (and their word counts) using
the {\hyperref[\detokenize{clicanalysis/counts:counts}]{\sphinxcrossref{\DUrole{std,std-ref}{Counts}}}} tab.
\subsubsection*{Texts in CLiC}


\begin{savenotes}\sphinxattablestart
\sphinxthistablewithglobalstyle
\centering
\sphinxcapstartof{table}
\sphinxthecaptionisattop
\sphinxcaption{DNov \textendash{} Dickens’s Novels (15 texts)}\label{\detokenize{appendices:id2}}
\sphinxaftertopcaption
\begin{tabular}[t]{\X{60}{100}\X{40}{100}}
\sphinxtoprule
\sphinxstyletheadfamily 
\sphinxAtStartPar
Title
&\sphinxstyletheadfamily 
\sphinxAtStartPar
Author
\\
\sphinxmidrule
\sphinxtableatstartofbodyhook
\sphinxAtStartPar
A Tale of Two Cities
&
\sphinxAtStartPar
Charles Dickens
\\
\sphinxhline
\sphinxAtStartPar
Barnaby Rudge
&
\sphinxAtStartPar
Charles Dickens
\\
\sphinxhline
\sphinxAtStartPar
Bleak House
&
\sphinxAtStartPar
Charles Dickens
\\
\sphinxhline
\sphinxAtStartPar
David Copperfield
&
\sphinxAtStartPar
Charles Dickens
\\
\sphinxhline
\sphinxAtStartPar
Dombey and Son
&
\sphinxAtStartPar
Charles Dickens
\\
\sphinxhline
\sphinxAtStartPar
Great Expectations
&
\sphinxAtStartPar
Charles Dickens
\\
\sphinxhline
\sphinxAtStartPar
Hard Times
&
\sphinxAtStartPar
Charles Dickens
\\
\sphinxhline
\sphinxAtStartPar
Little Dorrit
&
\sphinxAtStartPar
Charles Dickens
\\
\sphinxhline
\sphinxAtStartPar
Martin Chuzzlewit
&
\sphinxAtStartPar
Charles Dickens
\\
\sphinxhline
\sphinxAtStartPar
Nicholas Nickleby
&
\sphinxAtStartPar
Charles Dickens
\\
\sphinxhline
\sphinxAtStartPar
Oliver Twist
&
\sphinxAtStartPar
Charles Dickens
\\
\sphinxhline
\sphinxAtStartPar
Our Mutual Friend
&
\sphinxAtStartPar
Charles Dickens
\\
\sphinxhline
\sphinxAtStartPar
Pickwick Papers
&
\sphinxAtStartPar
Charles Dickens
\\
\sphinxhline
\sphinxAtStartPar
The Mystery of Edwin Drood
&
\sphinxAtStartPar
Charles Dickens
\\
\sphinxhline
\sphinxAtStartPar
The Old Curiosity Shop
&
\sphinxAtStartPar
Charles Dickens
\\
\sphinxbottomrule
\end{tabular}
\sphinxtableafterendhook\par
\sphinxattableend\end{savenotes}


\begin{savenotes}\sphinxattablestart
\sphinxthistablewithglobalstyle
\centering
\sphinxcapstartof{table}
\sphinxthecaptionisattop
\sphinxcaption{19C \textendash{} 19th Century Reference Corpus (29 texts)}\label{\detokenize{appendices:id3}}
\sphinxaftertopcaption
\begin{tabular}[t]{\X{60}{100}\X{40}{100}}
\sphinxtoprule
\sphinxstyletheadfamily 
\sphinxAtStartPar
Title
&\sphinxstyletheadfamily 
\sphinxAtStartPar
Author
\\
\sphinxmidrule
\sphinxtableatstartofbodyhook
\sphinxAtStartPar
Agnes Grey
&
\sphinxAtStartPar
Anne Brontë
\\
\sphinxhline
\sphinxAtStartPar
Antonina or, the Fall of Rome
&
\sphinxAtStartPar
Wilkie Collins
\\
\sphinxhline
\sphinxAtStartPar
Armadale
&
\sphinxAtStartPar
Wilkie Collins
\\
\sphinxhline
\sphinxAtStartPar
Cranford
&
\sphinxAtStartPar
Elizabeth Gaskell
\\
\sphinxhline
\sphinxAtStartPar
Daniel Deronda
&
\sphinxAtStartPar
George Eliot
\\
\sphinxhline
\sphinxAtStartPar
Dracula
&
\sphinxAtStartPar
Bram Stoker
\\
\sphinxhline
\sphinxAtStartPar
Emma
&
\sphinxAtStartPar
Jane Austen
\\
\sphinxhline
\sphinxAtStartPar
Frankenstein
&
\sphinxAtStartPar
Mary Shelley
\\
\sphinxhline
\sphinxAtStartPar
Jane Eyre
&
\sphinxAtStartPar
Charlotte Brontë
\\
\sphinxhline
\sphinxAtStartPar
Jude the Obscure
&
\sphinxAtStartPar
Thomas Hardy
\\
\sphinxhline
\sphinxAtStartPar
Lady Audley’s Secret
&
\sphinxAtStartPar
Mary Elizabeth Braddon
\\
\sphinxhline
\sphinxAtStartPar
Mary Barton
&
\sphinxAtStartPar
Elizabeth Gaskell
\\
\sphinxhline
\sphinxAtStartPar
North and South
&
\sphinxAtStartPar
Elizabeth Gaskell
\\
\sphinxhline
\sphinxAtStartPar
Persuasion
&
\sphinxAtStartPar
Jane Austen
\\
\sphinxhline
\sphinxAtStartPar
Pride and Prejudice
&
\sphinxAtStartPar
Jane Austen
\\
\sphinxhline
\sphinxAtStartPar
Sybil, or the two nations
&
\sphinxAtStartPar
Benjamin Disraeli
\\
\sphinxhline
\sphinxAtStartPar
Tess of the D’Urbervilles
&
\sphinxAtStartPar
Thomas Hardy
\\
\sphinxhline
\sphinxAtStartPar
The Hound of the Baskervilles
&
\sphinxAtStartPar
Arthur Conan Doyle
\\
\sphinxhline
\sphinxAtStartPar
The Last Days of Pompeii
&
\sphinxAtStartPar
Edward George Bulwer\sphinxhyphen{}Lytton
\\
\sphinxhline
\sphinxAtStartPar
The Mill on the Floss
&
\sphinxAtStartPar
George Eliot
\\
\sphinxhline
\sphinxAtStartPar
The Picture of Dorian Gray
&
\sphinxAtStartPar
Oscar Wilde
\\
\sphinxhline
\sphinxAtStartPar
The Professor
&
\sphinxAtStartPar
Charlotte Brontë
\\
\sphinxhline
\sphinxAtStartPar
The Return of the Native
&
\sphinxAtStartPar
Thomas Hardy
\\
\sphinxhline
\sphinxAtStartPar
The Small House at Allington
&
\sphinxAtStartPar
Anthony Trollope
\\
\sphinxhline
\sphinxAtStartPar
The Strange Case of Dr Jekyll and Mr Hyde
&
\sphinxAtStartPar
Robert Louis Stevenson
\\
\sphinxhline
\sphinxAtStartPar
The Woman in White
&
\sphinxAtStartPar
Wilkie Collins
\\
\sphinxhline
\sphinxAtStartPar
Vanity Fair
&
\sphinxAtStartPar
William Makepeace Thackeray
\\
\sphinxhline
\sphinxAtStartPar
Vivian Grey
&
\sphinxAtStartPar
Benjamin Disraeli
\\
\sphinxhline
\sphinxAtStartPar
Wuthering Heights
&
\sphinxAtStartPar
Emily Brontë
\\
\sphinxbottomrule
\end{tabular}
\sphinxtableafterendhook\par
\sphinxattableend\end{savenotes}


\begin{savenotes}
\sphinxatlongtablestart
\sphinxthistablewithglobalstyle
\makeatletter
  \LTleft \@totalleftmargin plus1fill
  \LTright\dimexpr\columnwidth-\@totalleftmargin-\linewidth\relax plus1fill
\makeatother
\begin{longtable}{\X{60}{100}\X{40}{100}}
\sphinxthelongtablecaptionisattop
\caption{ChiLit \textendash{} 19th Century Children’s Literature Corpus (71 texts)\strut}\label{\detokenize{appendices:id4}}\\*[\sphinxlongtablecapskipadjust]
\sphinxtoprule
\sphinxstyletheadfamily 
\sphinxAtStartPar
Title
&\sphinxstyletheadfamily 
\sphinxAtStartPar
Author
\\
\sphinxmidrule
\endfirsthead

\multicolumn{2}{c}{\sphinxnorowcolor
    \makebox[0pt]{\sphinxtablecontinued{\tablename\ \thetable{} \textendash{} continued from previous page}}%
}\\
\sphinxtoprule
\sphinxstyletheadfamily 
\sphinxAtStartPar
Title
&\sphinxstyletheadfamily 
\sphinxAtStartPar
Author
\\
\sphinxmidrule
\endhead

\sphinxbottomrule
\multicolumn{2}{r}{\sphinxnorowcolor
    \makebox[0pt][r]{\sphinxtablecontinued{continues on next page}}%
}\\
\endfoot

\endlastfoot
\sphinxtableatstartofbodyhook

\sphinxAtStartPar
A World of Girls: The Story of a School
&\begin{enumerate}
\sphinxsetlistlabels{\Alph}{enumi}{enumii}{}{.}%
\setcounter{enumi}{11}
\item {} \begin{enumerate}
\sphinxsetlistlabels{\Alph}{enumii}{enumiii}{}{.}%
\setcounter{enumii}{19}
\item {} 
\sphinxAtStartPar
Meade

\end{enumerate}

\end{enumerate}
\\
\sphinxhline
\sphinxAtStartPar
Adventures in Wallypug\sphinxhyphen{}Land
&\begin{enumerate}
\sphinxsetlistlabels{\Alph}{enumi}{enumii}{}{.}%
\setcounter{enumi}{6}
\item {} \begin{enumerate}
\sphinxsetlistlabels{\Alph}{enumii}{enumiii}{}{.}%
\setcounter{enumii}{4}
\item {} 
\sphinxAtStartPar
Farrow

\end{enumerate}

\end{enumerate}
\\
\sphinxhline
\sphinxAtStartPar
Alice’s Adventures in Wonderland
&
\sphinxAtStartPar
Lewis Carroll
\\
\sphinxhline
\sphinxAtStartPar
Allan Quatermain
&\begin{enumerate}
\sphinxsetlistlabels{\Alph}{enumi}{enumii}{}{.}%
\setcounter{enumi}{7}
\item {} 
\sphinxAtStartPar
Rider Haggard

\end{enumerate}
\\
\sphinxhline
\sphinxAtStartPar
Alone in London
&
\sphinxAtStartPar
Hesba Stretton
\\
\sphinxhline
\sphinxAtStartPar
At the Back of the North Wind
&
\sphinxAtStartPar
George McDonald
\\
\sphinxhline
\sphinxAtStartPar
Black Beauty
&
\sphinxAtStartPar
Anna Sewell
\\
\sphinxhline
\sphinxAtStartPar
Dream Days
&
\sphinxAtStartPar
Kenneth Grahame
\\
\sphinxhline
\sphinxAtStartPar
Eric; Or, Little by Little
&\begin{enumerate}
\sphinxsetlistlabels{\Alph}{enumi}{enumii}{}{.}%
\setcounter{enumi}{5}
\item {} \begin{enumerate}
\sphinxsetlistlabels{\Alph}{enumii}{enumiii}{}{.}%
\setcounter{enumii}{22}
\item {} 
\sphinxAtStartPar
Farrar

\end{enumerate}

\end{enumerate}
\\
\sphinxhline
\sphinxAtStartPar
Feats on the Fiord
&
\sphinxAtStartPar
Harriet Martineau
\\
\sphinxhline
\sphinxAtStartPar
Five Children and It
&\begin{enumerate}
\sphinxsetlistlabels{\Alph}{enumi}{enumii}{}{.}%
\setcounter{enumi}{4}
\item {} 
\sphinxAtStartPar
Nesbit

\end{enumerate}
\\
\sphinxhline
\sphinxAtStartPar
Holiday House: A Series of Tales
&
\sphinxAtStartPar
Catherine Sinclair
\\
\sphinxhline
\sphinxAtStartPar
Jackanapes
&
\sphinxAtStartPar
Juliana Horatia Ewing
\\
\sphinxhline
\sphinxAtStartPar
Jessica’s First Prayer
&
\sphinxAtStartPar
Hesba Stretton
\\
\sphinxhline
\sphinxAtStartPar
Kidnapped
&
\sphinxAtStartPar
Robert Louis Stevenson
\\
\sphinxhline
\sphinxAtStartPar
King Solomon’s Mines
&\begin{enumerate}
\sphinxsetlistlabels{\Alph}{enumi}{enumii}{}{.}%
\setcounter{enumi}{7}
\item {} 
\sphinxAtStartPar
Rider Haggard

\end{enumerate}
\\
\sphinxhline
\sphinxAtStartPar
Leila at home
&
\sphinxAtStartPar
Ann Fraser Tytler
\\
\sphinxhline
\sphinxAtStartPar
Little Meg’s children
&
\sphinxAtStartPar
Hesba Stretton
\\
\sphinxhline
\sphinxAtStartPar
Madam How and Lady Why; Or, First Lessons in Earth Lore for Children
&
\sphinxAtStartPar
Charles Kingsley
\\
\sphinxhline
\sphinxAtStartPar
Masterman Ready; Or, The Wreck of the Pacific
&
\sphinxAtStartPar
Frederick Marryat
\\
\sphinxhline
\sphinxAtStartPar
Moonfleet
&\begin{enumerate}
\sphinxsetlistlabels{\Alph}{enumi}{enumii}{}{.}%
\setcounter{enumi}{9}
\item {} 
\sphinxAtStartPar
Meade Falkner

\end{enumerate}
\\
\sphinxhline
\sphinxAtStartPar
Mopsa the Fairy
&
\sphinxAtStartPar
Jean Ingelow
\\
\sphinxhline
\sphinxAtStartPar
Mrs. Overtheway’s Remembrances
&
\sphinxAtStartPar
Juliana Horatia Ewing
\\
\sphinxhline
\sphinxAtStartPar
Nine Unlikely Tales
&\begin{enumerate}
\sphinxsetlistlabels{\Alph}{enumi}{enumii}{}{.}%
\setcounter{enumi}{4}
\item {} 
\sphinxAtStartPar
Nesbit

\end{enumerate}
\\
\sphinxhline
\sphinxAtStartPar
Peter Pan
&\begin{enumerate}
\sphinxsetlistlabels{\Alph}{enumi}{enumii}{}{.}%
\setcounter{enumi}{9}
\item {} \begin{enumerate}
\sphinxsetlistlabels{\Alph}{enumii}{enumiii}{}{.}%
\setcounter{enumii}{12}
\item {} 
\sphinxAtStartPar
Barrie

\end{enumerate}

\end{enumerate}
\\
\sphinxhline
\sphinxAtStartPar
Prince Prigio
&
\sphinxAtStartPar
Andrew Lang
\\
\sphinxhline
\sphinxAtStartPar
Stalky and Co
&
\sphinxAtStartPar
Rudyard Kipling
\\
\sphinxhline
\sphinxAtStartPar
The Book of Dragons
&\begin{enumerate}
\sphinxsetlistlabels{\Alph}{enumi}{enumii}{}{.}%
\setcounter{enumi}{4}
\item {} 
\sphinxAtStartPar
Nesbit

\end{enumerate}
\\
\sphinxhline
\sphinxAtStartPar
The Brass Bottle
&\begin{enumerate}
\sphinxsetlistlabels{\Alph}{enumi}{enumii}{}{.}%
\setcounter{enumi}{5}
\item {} 
\sphinxAtStartPar
Anstey

\end{enumerate}
\\
\sphinxhline
\sphinxAtStartPar
The Carved Lions
&
\sphinxAtStartPar
Mary Louisa Molesworth
\\
\sphinxhline
\sphinxAtStartPar
The Children of the New Forest
&
\sphinxAtStartPar
Frederick Marryat
\\
\sphinxhline
\sphinxAtStartPar
The Coral Island: A Tale of the Pacific Ocean
&\begin{enumerate}
\sphinxsetlistlabels{\Alph}{enumi}{enumii}{}{.}%
\setcounter{enumi}{17}
\item {} \begin{enumerate}
\sphinxsetlistlabels{\Alph}{enumii}{enumiii}{}{.}%
\setcounter{enumii}{12}
\item {} 
\sphinxAtStartPar
Ballantyne

\end{enumerate}

\end{enumerate}
\\
\sphinxhline
\sphinxAtStartPar
The Crofton Boys
&
\sphinxAtStartPar
Harriet Martineau
\\
\sphinxhline
\sphinxAtStartPar
The Cuckoo Clock
&
\sphinxAtStartPar
Mary Louisa Molesworth
\\
\sphinxhline
\sphinxAtStartPar
The Daisy Chain, or Aspirations
&
\sphinxAtStartPar
Charlotte M. Yonge
\\
\sphinxhline
\sphinxAtStartPar
The Dove in the Eagle’s Nest
&
\sphinxAtStartPar
Charlotte M. Yonge
\\
\sphinxhline
\sphinxAtStartPar
The Fifth Form at Saint Dominic’s: A School Story
&
\sphinxAtStartPar
Talbot Baines Reed
\\
\sphinxhline
\sphinxAtStartPar
The Golden Age
&
\sphinxAtStartPar
Kenneth Grahame
\\
\sphinxhline
\sphinxAtStartPar
The Happy Prince, and Other Tales
&
\sphinxAtStartPar
Oscar Wilde
\\
\sphinxhline
\sphinxAtStartPar
The Heir of Redclyffe
&
\sphinxAtStartPar
Charlotte M. Yonge
\\
\sphinxhline
\sphinxAtStartPar
The Jungle Book
&
\sphinxAtStartPar
Rudyard Kipling
\\
\sphinxhline
\sphinxAtStartPar
The King of the Golden River; or, the Black Brothers: A Legend of Stiria
&
\sphinxAtStartPar
John Ruskin
\\
\sphinxhline
\sphinxAtStartPar
The Little Duke: Richard the Fearless
&
\sphinxAtStartPar
Charlotte M. Yonge
\\
\sphinxhline
\sphinxAtStartPar
The Peasant and the Prince
&
\sphinxAtStartPar
Harriet Martineau
\\
\sphinxhline
\sphinxAtStartPar
The Princess and the Goblin
&
\sphinxAtStartPar
George McDonald
\\
\sphinxhline
\sphinxAtStartPar
The Railway Children
&\begin{enumerate}
\sphinxsetlistlabels{\Alph}{enumi}{enumii}{}{.}%
\setcounter{enumi}{4}
\item {} 
\sphinxAtStartPar
Nesbit

\end{enumerate}
\\
\sphinxhline
\sphinxAtStartPar
The Rival Crusoes; or The Shipwreck
&
\sphinxAtStartPar
Agnes Strickland
\\
\sphinxhline
\sphinxAtStartPar
The Rose and the Ring
&
\sphinxAtStartPar
William Makepeace Thackeray
\\
\sphinxhline
\sphinxAtStartPar
The Secret Garden
&
\sphinxAtStartPar
Frances Hodgson Burnett
\\
\sphinxhline
\sphinxAtStartPar
The Settlers at Home
&
\sphinxAtStartPar
Harriet Martineau
\\
\sphinxhline
\sphinxAtStartPar
The Settlers in Canada
&
\sphinxAtStartPar
Frederick Marryat
\\
\sphinxhline
\sphinxAtStartPar
The Story of the Amulet
&\begin{enumerate}
\sphinxsetlistlabels{\Alph}{enumi}{enumii}{}{.}%
\setcounter{enumi}{4}
\item {} 
\sphinxAtStartPar
Nesbit

\end{enumerate}
\\
\sphinxhline
\sphinxAtStartPar
The Story of the Treasure Seekers
&\begin{enumerate}
\sphinxsetlistlabels{\Alph}{enumi}{enumii}{}{.}%
\setcounter{enumi}{4}
\item {} 
\sphinxAtStartPar
Nesbit

\end{enumerate}
\\
\sphinxhline
\sphinxAtStartPar
The Surprising Adventures of Sir Toady Lion
&
\sphinxAtStartPar
S.R. Crockett
\\
\sphinxhline
\sphinxAtStartPar
The Tale of Benjamin Bunny
&
\sphinxAtStartPar
Beatrix Potter
\\
\sphinxhline
\sphinxAtStartPar
The Tale of Jemima Puddle\sphinxhyphen{}Duck
&
\sphinxAtStartPar
Beatrix Potter
\\
\sphinxhline
\sphinxAtStartPar
The Tale of Peter Rabbit
&
\sphinxAtStartPar
Beatrix Potter
\\
\sphinxhline
\sphinxAtStartPar
The Tale of Squirrel Nutkin
&
\sphinxAtStartPar
Beatrix Potter
\\
\sphinxhline
\sphinxAtStartPar
The Tale of the Flopsy Bunnies
&
\sphinxAtStartPar
Beatrix Potter
\\
\sphinxhline
\sphinxAtStartPar
The Tale of Two Bad Mice
&
\sphinxAtStartPar
Beatrix Potter
\\
\sphinxhline
\sphinxAtStartPar
The Tapestry Room: A Child’s Romance
&
\sphinxAtStartPar
Mary Louisa Molesworth
\\
\sphinxhline
\sphinxAtStartPar
The Three Mulla\sphinxhyphen{}mulgars
&
\sphinxAtStartPar
Walter de la Mare
\\
\sphinxhline
\sphinxAtStartPar
The Water\sphinxhyphen{}Babies
&
\sphinxAtStartPar
Charles Kingsley
\\
\sphinxhline
\sphinxAtStartPar
The Wind in the Willows
&
\sphinxAtStartPar
Kenneth Grahame
\\
\sphinxhline
\sphinxAtStartPar
Through the Looking\sphinxhyphen{}Glass
&
\sphinxAtStartPar
Lewis Carroll
\\
\sphinxhline
\sphinxAtStartPar
Tom Brown’s Schooldays
&
\sphinxAtStartPar
Thomas Hughes
\\
\sphinxhline
\sphinxAtStartPar
Treasure Island
&
\sphinxAtStartPar
Robert Louis Stevenson
\\
\sphinxhline
\sphinxAtStartPar
Vice Versa; or, A Lesson to Fathers
&\begin{enumerate}
\sphinxsetlistlabels{\Alph}{enumi}{enumii}{}{.}%
\setcounter{enumi}{5}
\item {} 
\sphinxAtStartPar
Anstey

\end{enumerate}
\\
\sphinxhline
\sphinxAtStartPar
Winning His Spurs. A Tale of the Crusades
&\begin{enumerate}
\sphinxsetlistlabels{\Alph}{enumi}{enumii}{}{.}%
\setcounter{enumi}{6}
\item {} \begin{enumerate}
\sphinxsetlistlabels{\Alph}{enumii}{enumiii}{}{.}%
\item {} 
\sphinxAtStartPar
Henty

\end{enumerate}

\end{enumerate}
\\
\sphinxhline
\sphinxAtStartPar
With Clive in India; Or, The Beginnings of an Empire
&\begin{enumerate}
\sphinxsetlistlabels{\Alph}{enumi}{enumii}{}{.}%
\setcounter{enumi}{6}
\item {} \begin{enumerate}
\sphinxsetlistlabels{\Alph}{enumii}{enumiii}{}{.}%
\item {} 
\sphinxAtStartPar
Henty

\end{enumerate}

\end{enumerate}
\\
\sphinxhline
\sphinxAtStartPar
Wood Magic, a Fable
&
\sphinxAtStartPar
Richard Jefferies
\\
\sphinxbottomrule
\end{longtable}
\sphinxtableafterendhook
\sphinxatlongtableend
\end{savenotes}


\begin{savenotes}
\sphinxatlongtablestart
\sphinxthistablewithglobalstyle
\makeatletter
  \LTleft \@totalleftmargin plus1fill
  \LTright\dimexpr\columnwidth-\@totalleftmargin-\linewidth\relax plus1fill
\makeatother
\begin{longtable}{\X{60}{100}\X{40}{100}}
\sphinxthelongtablecaptionisattop
\caption{ArTs \textendash{} Additional Requested Texts (31 texts)\strut}\label{\detokenize{appendices:id5}}\\*[\sphinxlongtablecapskipadjust]
\sphinxtoprule
\sphinxstyletheadfamily 
\sphinxAtStartPar
Title
&\sphinxstyletheadfamily 
\sphinxAtStartPar
Author
\\
\sphinxmidrule
\endfirsthead

\multicolumn{2}{c}{\sphinxnorowcolor
    \makebox[0pt]{\sphinxtablecontinued{\tablename\ \thetable{} \textendash{} continued from previous page}}%
}\\
\sphinxtoprule
\sphinxstyletheadfamily 
\sphinxAtStartPar
Title
&\sphinxstyletheadfamily 
\sphinxAtStartPar
Author
\\
\sphinxmidrule
\endhead

\sphinxbottomrule
\multicolumn{2}{r}{\sphinxnorowcolor
    \makebox[0pt][r]{\sphinxtablecontinued{continues on next page}}%
}\\
\endfoot

\endlastfoot
\sphinxtableatstartofbodyhook

\sphinxAtStartPar
A Christmas Carol: A Ghost Story of Christmas
&
\sphinxAtStartPar
Charles Dickens
\\
\sphinxhline
\sphinxAtStartPar
A Room with a View
&\begin{enumerate}
\sphinxsetlistlabels{\Alph}{enumi}{enumii}{}{.}%
\setcounter{enumi}{4}
\item {} \begin{enumerate}
\sphinxsetlistlabels{\Alph}{enumii}{enumiii}{}{.}%
\setcounter{enumii}{12}
\item {} 
\sphinxAtStartPar
Forster

\end{enumerate}

\end{enumerate}
\\
\sphinxhline
\sphinxAtStartPar
Adventures of Huckleberry Finn, Complete
&
\sphinxAtStartPar
Mark Twain
\\
\sphinxhline
\sphinxAtStartPar
American Notes for General Circulation
&
\sphinxAtStartPar
Charles Dickens
\\
\sphinxhline
\sphinxAtStartPar
Gulliver’s Travels
&
\sphinxAtStartPar
Jonathan Swift
\\
\sphinxhline
\sphinxAtStartPar
Heart of Darkness
&
\sphinxAtStartPar
Joseph Conrad
\\
\sphinxhline
\sphinxAtStartPar
Lady Susan
&
\sphinxAtStartPar
Jane Austen
\\
\sphinxhline
\sphinxAtStartPar
Mansfield Park
&
\sphinxAtStartPar
Jane Austen
\\
\sphinxhline
\sphinxAtStartPar
Middlemarch
&
\sphinxAtStartPar
George Eliot
\\
\sphinxhline
\sphinxAtStartPar
Northanger Abbey
&
\sphinxAtStartPar
Jane Austen
\\
\sphinxhline
\sphinxAtStartPar
Pictures from Italy
&
\sphinxAtStartPar
Charles Dickens
\\
\sphinxhline
\sphinxAtStartPar
Sense and Sensibility
&
\sphinxAtStartPar
Jane Austen
\\
\sphinxhline
\sphinxAtStartPar
Shirley
&
\sphinxAtStartPar
Charlotte Brontë
\\
\sphinxhline
\sphinxAtStartPar
Silas Marner: The Weaver of Raveloe
&
\sphinxAtStartPar
George Eliot
\\
\sphinxhline
\sphinxAtStartPar
Sketches by Boz: Illustrative of Every\sphinxhyphen{}Day Life and Every\sphinxhyphen{}Day People
&
\sphinxAtStartPar
Charles Dickens
\\
\sphinxhline
\sphinxAtStartPar
The Awakening and Selected Short Stories
&
\sphinxAtStartPar
Kate Chopin
\\
\sphinxhline
\sphinxAtStartPar
The Jungle
&
\sphinxAtStartPar
Upton Sinclair
\\
\sphinxhline
\sphinxAtStartPar
The Moonstone
&
\sphinxAtStartPar
Wilkie Collins
\\
\sphinxhline
\sphinxAtStartPar
The Portrait of a Lady Volume 1 (of 2)
&
\sphinxAtStartPar
Henry James
\\
\sphinxhline
\sphinxAtStartPar
The Portrait of a Lady Volume 2 (of 2)
&
\sphinxAtStartPar
Henry James
\\
\sphinxhline
\sphinxAtStartPar
The Return of the Soldier
&
\sphinxAtStartPar
Rebecca West
\\
\sphinxhline
\sphinxAtStartPar
The Sign of the Four
&
\sphinxAtStartPar
Arthur Conan Doyle
\\
\sphinxhline
\sphinxAtStartPar
The Tenant of Wildfell Hall
&
\sphinxAtStartPar
Anne Brontë
\\
\sphinxhline
\sphinxAtStartPar
The Time Machine: An Invention
&\begin{enumerate}
\sphinxsetlistlabels{\Alph}{enumi}{enumii}{}{.}%
\setcounter{enumi}{7}
\item {} \begin{enumerate}
\sphinxsetlistlabels{\Alph}{enumii}{enumiii}{}{.}%
\setcounter{enumii}{6}
\item {} 
\sphinxAtStartPar
Wells

\end{enumerate}

\end{enumerate}
\\
\sphinxhline
\sphinxAtStartPar
The Uncommercial Traveller
&
\sphinxAtStartPar
Charles Dickens
\\
\sphinxhline
\sphinxAtStartPar
The War of the Worlds
&\begin{enumerate}
\sphinxsetlistlabels{\Alph}{enumi}{enumii}{}{.}%
\setcounter{enumi}{7}
\item {} \begin{enumerate}
\sphinxsetlistlabels{\Alph}{enumii}{enumiii}{}{.}%
\setcounter{enumii}{6}
\item {} 
\sphinxAtStartPar
Wells

\end{enumerate}

\end{enumerate}
\\
\sphinxhline
\sphinxAtStartPar
The Yellow Wallpaper
&
\sphinxAtStartPar
Charlotte Perkins Gilman
\\
\sphinxhline
\sphinxAtStartPar
Twelve Years a Slave
&
\sphinxAtStartPar
Solomon Northup
\\
\sphinxhline
\sphinxAtStartPar
Villette
&
\sphinxAtStartPar
Charlotte Brontë
\\
\sphinxhline
\sphinxAtStartPar
What Maisie Knew
&
\sphinxAtStartPar
Henry James
\\
\sphinxhline
\sphinxAtStartPar
Women in Love
&\begin{enumerate}
\sphinxsetlistlabels{\Alph}{enumi}{enumii}{}{.}%
\setcounter{enumi}{3}
\item {} \begin{enumerate}
\sphinxsetlistlabels{\Alph}{enumii}{enumiii}{}{.}%
\setcounter{enumii}{7}
\item {} 
\sphinxAtStartPar
Lawrence

\end{enumerate}

\end{enumerate}
\\
\sphinxbottomrule
\end{longtable}
\sphinxtableafterendhook
\sphinxatlongtableend
\end{savenotes}


\begin{savenotes}
\sphinxatlongtablestart
\sphinxthistablewithglobalstyle
\makeatletter
  \LTleft \@totalleftmargin plus1fill
  \LTright\dimexpr\columnwidth-\@totalleftmargin-\linewidth\relax plus1fill
\makeatother
\begin{longtable}{\X{50}{100}\X{30}{100}\X{20}{100}}
\sphinxthelongtablecaptionisattop
\caption{DE19 \textendash{} German 19th Century Reference Corpus (33 texts)\strut}\label{\detokenize{appendices:id6}}\\*[\sphinxlongtablecapskipadjust]
\sphinxtoprule
\sphinxstyletheadfamily 
\sphinxAtStartPar
Title
&\sphinxstyletheadfamily 
\sphinxAtStartPar
Author
&\sphinxstyletheadfamily 
\sphinxAtStartPar
Year
\\
\sphinxmidrule
\endfirsthead

\multicolumn{3}{c}{\sphinxnorowcolor
    \makebox[0pt]{\sphinxtablecontinued{\tablename\ \thetable{} \textendash{} continued from previous page}}%
}\\
\sphinxtoprule
\sphinxstyletheadfamily 
\sphinxAtStartPar
Title
&\sphinxstyletheadfamily 
\sphinxAtStartPar
Author
&\sphinxstyletheadfamily 
\sphinxAtStartPar
Year
\\
\sphinxmidrule
\endhead

\sphinxbottomrule
\multicolumn{3}{r}{\sphinxnorowcolor
    \makebox[0pt][r]{\sphinxtablecontinued{continues on next page}}%
}\\
\endfoot

\endlastfoot
\sphinxtableatstartofbodyhook

\sphinxAtStartPar
Die Günderode
&
\sphinxAtStartPar
Arnim, B.
&
\sphinxAtStartPar
1840
\\
\sphinxhline
\sphinxAtStartPar
Clemens Brentanos Frühlingskranz
&
\sphinxAtStartPar
Arnim, B.
&
\sphinxAtStartPar
1844
\\
\sphinxhline
\sphinxAtStartPar
Fanny Förster
&
\sphinxAtStartPar
Boy\sphinxhyphen{}Ed, I.
&
\sphinxAtStartPar
1889
\\
\sphinxhline
\sphinxAtStartPar
Kampf um Rom
&
\sphinxAtStartPar
Dahn, F.
&
\sphinxAtStartPar
1876
\\
\sphinxhline
\sphinxAtStartPar
Sibilla Dalmar
&
\sphinxAtStartPar
Dohm, H.
&
\sphinxAtStartPar
1896
\\
\sphinxhline
\sphinxAtStartPar
Schicksale einer Seele
&
\sphinxAtStartPar
Dohm, H.
&
\sphinxAtStartPar
1899
\\
\sphinxhline
\sphinxAtStartPar
Bozena
&
\sphinxAtStartPar
Ebner\sphinxhyphen{}Eschenbach, M.
&
\sphinxAtStartPar
1876
\\
\sphinxhline
\sphinxAtStartPar
Vor dem Sturm
&
\sphinxAtStartPar
Fontane, T.
&
\sphinxAtStartPar
1878
\\
\sphinxhline
\sphinxAtStartPar
Der Stechlin
&
\sphinxAtStartPar
Fontane, T.
&
\sphinxAtStartPar
1897/98
\\
\sphinxhline
\sphinxAtStartPar
Die letzte Reckenburgerin
&
\sphinxAtStartPar
François, Louise.
&
\sphinxAtStartPar
1871
\\
\sphinxhline
\sphinxAtStartPar
Soll und Haben
&
\sphinxAtStartPar
Freytag, G.
&
\sphinxAtStartPar
1855
\\
\sphinxhline
\sphinxAtStartPar
Die verlorene Handschrift
&
\sphinxAtStartPar
Freytag, G.
&
\sphinxAtStartPar
1864
\\
\sphinxhline
\sphinxAtStartPar
Schloß Hubertus
&
\sphinxAtStartPar
Ganghofer, L.
&
\sphinxAtStartPar
1895
\\
\sphinxhline
\sphinxAtStartPar
Sibylle
&
\sphinxAtStartPar
Hahn\sphinxhyphen{}Hahn, I. G.
&
\sphinxAtStartPar
1846
\\
\sphinxhline
\sphinxAtStartPar
Der Sonnenwirt
&
\sphinxAtStartPar
Kurz, H.
&
\sphinxAtStartPar
1855
\\
\sphinxhline
\sphinxAtStartPar
Jenny
&
\sphinxAtStartPar
Lewald, F.
&
\sphinxAtStartPar
1843
\\
\sphinxhline
\sphinxAtStartPar
Das Heideprinzeßchen
&
\sphinxAtStartPar
Marlitt, E.
&
\sphinxAtStartPar
1871/72
\\
\sphinxhline
\sphinxAtStartPar
Goldelse
&
\sphinxAtStartPar
Marlitt, E.
&
\sphinxAtStartPar
1866/67
\\
\sphinxhline
\sphinxAtStartPar
Das Geheimnis der alten Mamsell
&
\sphinxAtStartPar
Marlitt, E.
&
\sphinxAtStartPar
1868
\\
\sphinxhline
\sphinxAtStartPar
Am Jenseits
&
\sphinxAtStartPar
May, K.
&
\sphinxAtStartPar
1898
\\
\sphinxhline
\sphinxAtStartPar
Jürg Jenatsch
&
\sphinxAtStartPar
Meyer, C. F.
&
\sphinxAtStartPar
1874
\\
\sphinxhline
\sphinxAtStartPar
Der Büttnerbauer
&
\sphinxAtStartPar
Polenz, W.
&
\sphinxAtStartPar
1895
\\
\sphinxhline
\sphinxAtStartPar
Der Schüdderump
&
\sphinxAtStartPar
Raabe, W.
&
\sphinxAtStartPar
1869/70
\\
\sphinxhline
\sphinxAtStartPar
Der Hungerpastor
&
\sphinxAtStartPar
Raabe, W.
&
\sphinxAtStartPar
1863/64
\\
\sphinxhline
\sphinxAtStartPar
Die Leute aus dem Walde, ihre Sterne, Wege und Schicksale
&
\sphinxAtStartPar
Raabe, W.
&
\sphinxAtStartPar
1863
\\
\sphinxhline
\sphinxAtStartPar
Anna
&
\sphinxAtStartPar
Schopenhauer, A.
&
\sphinxAtStartPar
1845
\\
\sphinxhline
\sphinxAtStartPar
Problematische Naturen. Erste Abtheilung
&
\sphinxAtStartPar
Spielhagen, F.
&
\sphinxAtStartPar
1861
\\
\sphinxhline
\sphinxAtStartPar
Hammer und Amboß
&
\sphinxAtStartPar
Spielhagen, F.
&
\sphinxAtStartPar
1869
\\
\sphinxhline
\sphinxAtStartPar
Die Mappe meines Urgroßvaters
&
\sphinxAtStartPar
Stifter, A.
&
\sphinxAtStartPar
1841
\\
\sphinxhline
\sphinxAtStartPar
Witiko
&
\sphinxAtStartPar
Stifter, A.
&
\sphinxAtStartPar
1855, 1865\textendash{}1867
\\
\sphinxhline
\sphinxAtStartPar
Vittoria Accorombona
&
\sphinxAtStartPar
Tieck, L.
&
\sphinxAtStartPar
1840
\\
\sphinxhline
\sphinxAtStartPar
Auch Einer 1
&
\sphinxAtStartPar
Vischer, F. T.
&
\sphinxAtStartPar
1879
\\
\sphinxhline
\sphinxAtStartPar
Auch Einer 2
&
\sphinxAtStartPar
Vischer, F. T.
&
\sphinxAtStartPar
1879
\\
\sphinxbottomrule
\end{longtable}
\sphinxtableafterendhook
\sphinxatlongtableend
\end{savenotes}
\subsubsection*{CLiC texts listed in A\sphinxhyphen{}Level and GCSE specifications}

\begin{figure}[htbp]
\centering
\capstart

\noindent\sphinxincludegraphics{{appendices-alevelgcsespecs-overview}.png}
\caption{Overview of CLiC texts listed in the AQA, Edexcel and OCR
A\sphinxhyphen{}Level / GCSE specifications}\label{\detokenize{appendices:id7}}\end{figure}

\sphinxstepscope


\chapter{What’s new}
\label{\detokenize{whatsnew:what-s-new}}\label{\detokenize{whatsnew:whatsnew}}\label{\detokenize{whatsnew::doc}}

\section{What’s new in CLiC 2.3}
\label{\detokenize{whatsnew:what-s-new-in-clic-2-3}}\begin{itemize}
\item {} 
\sphinxAtStartPar
CLiC has moved to \sphinxhref{https://clic-fiction.com}{clic\sphinxhyphen{}fiction.com}.

\item {} 
\sphinxAtStartPar
New \sphinxstylestrong{FlexiConc} tab for flexible reproducible concordance algorithms
(see {\hyperref[\detokenize{clicanalysis/flexiconc:flexiconc}]{\sphinxcrossref{\DUrole{std,std-ref}{FlexiConc}}}}).

\item {} 
\sphinxAtStartPar
New \sphinxstylestrong{DE19} corpus: German 19th Century Reference Corpus (33 texts),
compiled for the \sphinxstyleemphasis{Reading Concordances in the 21st Century} project.

\end{itemize}


\section{What’s new in CLiC 2.1}
\label{\detokenize{whatsnew:what-s-new-in-clic-2-1}}
\sphinxAtStartPar
CLiC v2.1 adds features to improve the accessibility of the website and
share more about the website’s policies with the users.


\section{What’s new in CLiC 2.0}
\label{\detokenize{whatsnew:what-s-new-in-clic-2-0}}
\sphinxAtStartPar
There have been major changes to the back\sphinxhyphen{}end as the underlying database
infrastructure has been replaced. For further technical details of this
replacement, please refer to the \sphinxcite{references:github-clic} repository. The focus of
this user guide is the CLiC interface. Due to the changes of the
database, some of the basic functions of CLiC \textendash{} like the concordance \textendash{}
will give slightly different results.

\sphinxAtStartPar
The main changes from previous versions to CLiC 2.0 that users should be
aware of are:
\begin{itemize}
\item {} 
\sphinxAtStartPar
\sphinxstylestrong{New look}:

\sphinxAtStartPar
The table of contents on the landing page has been replaced with a
carousel of summary images for the corpora (see
\hyperref[\detokenize{whatsnew:figure-landing-page-2-0-0-19c}]{Fig.\@ \ref{\detokenize{whatsnew:figure-landing-page-2-0-0-19c}}}). An interactive overview of
the corpora can now be found in the \sphinxstylestrong{Counts tab}.

\sphinxAtStartPar
Another change to the look is that in addition to the CLiC version
number (i.e. the “release” of CLiC), the interface now also displays
the version of the corpora. For a more detailed explanation see
{\hyperref[\detokenize{cliccorpora:the-clic-corpora}]{\sphinxcrossref{\DUrole{std,std-ref}{The CLiC corpora}}}}.

\begin{figure}[htbp]
\centering
\capstart

\noindent\sphinxincludegraphics{{figure-landing-page-2_0_0_19C}.png}
\caption{The CLiC home screen with the menu in the side panel on
the right}\label{\detokenize{whatsnew:id2}}\label{\detokenize{whatsnew:figure-landing-page-2-0-0-19c}}\end{figure}

\item {} 
\sphinxAtStartPar
Changes to the \sphinxstylestrong{tokenisation} (i.e. what CLiC recognises as a
“word”)

\sphinxAtStartPar
The tokenisation of texts has changed from CLiC 2.0. It is now based
on unicode standard rules, used both for queries and importing books.
By contrast to previous versions, the main changes include:
\begin{enumerate}
\sphinxsetlistlabels{\arabic}{enumi}{enumii}{}{.}%
\item {} 
\sphinxAtStartPar
Word forms ending in \sphinxtitleref{‘s} are now distinct types. So, a search for
\sphinxtitleref{mother} or \sphinxtitleref{Oliver} will no longer include \sphinxtitleref{mother’s} or
\sphinxtitleref{Oliver’s}.

\item {} 
\sphinxAtStartPar
Any surrounding punctuation is filtered from types, thus searching
for “connisseur” will return tokens “connisseur” and
“\_connisseur\_”.

\end{enumerate}

\sphinxAtStartPar
Also see the explanation of the Concordance tab
{\hyperref[\detokenize{clicanalysis/concordance:search-the-corpora}]{\sphinxcrossref{\DUrole{std,std-ref}{Search the corpora}}}}. More examples and the detailed technical
documentation of the tokenizer is available from
\sphinxcode{\sphinxupquote{clic.tokenizer}}.

\item {} 
\sphinxAtStartPar
\sphinxstylestrong{Searching all books by one author}

\sphinxAtStartPar
Across all tabs, the “search the corpora” field now allows you to
select all books from a particular author at once, e.g. all 7 books
by Jane Austen.

\item {} 
\sphinxAtStartPar
\sphinxstylestrong{Concordances tab}

\sphinxAtStartPar
In concordances, short titles of books are now always visible and full
titles can be seen hovering over the short title. Concordance queries
now follow the new tokenisation rules (and searches for exact types,
e.g. “Oliver” will no longer return “Oliver’s”) and accept wildcards
(e.g. you can now query “Oliver*” to find both “Oliver” and
“Oliver’s”). Information on \sphinxstylestrong{relative frequencies} has been added to
the output to simple concordances and to the new \sphinxstylestrong{distribution plot
view}. See the documentation of the Concordance tab for more details:
{\hyperref[\detokenize{clicanalysis/concordance:concordance}]{\sphinxcrossref{\DUrole{std,std-ref}{Concordance}}}}.

\item {} 
\sphinxAtStartPar
\sphinxstylestrong{Subset}

\sphinxAtStartPar
The new subset output shows percentage of text instead of relative
frequency (unlike the concordance output). For example when displaying
all quotes in book, the percentage indicates the proportion of the
words in quotes.

\item {} 
\sphinxAtStartPar
\sphinxstylestrong{Clusters}

\sphinxAtStartPar
The clusters search now ignores boundaries and only pays attention to
token order. i.e. clusters can span adjacent quotes, but not chapters
(as the title is in the way). You can also click on a cluster to
retrieve a \sphinxstylestrong{matching concordance search}. If only one book is
selected the frequency cutoff is 2.

\sphinxAtStartPar
In previous versions of CLiC there were minor discrepancies between
concordance and cluster counts. In CLiC 2.0, these have been resolved.

\item {} 
\sphinxAtStartPar
\sphinxstylestrong{Keywords tab}

\sphinxAtStartPar
The keyword output has been simplified (the expected target/ref columns
have been removed). You can click on a keyword or key cluster to
retrieve a \sphinxstylestrong{matching concordance search} and a new \sphinxstylestrong{swap button}
will reverse the target and reference corpus options.

\item {} 
\sphinxAtStartPar
\sphinxstylestrong{New Text tab}

\sphinxAtStartPar
The text tab shows full book contents (to replace the
chapter\sphinxhyphen{}by\sphinxhyphen{}chapter selection on CLiC 1.6 landing page). This tab is
designed to show one book at a time and allows you to navigate to
particular chapters. You can select which levels of annotation to
display on the text (e.g. Sentences, Quotes, Short suspensions, Long
suspensions, Embedded quotes). The text can now be easily copied and
pasted into Microsoft Word documents etc.

\item {} 
\sphinxAtStartPar
\sphinxstylestrong{New Counts tab}

\sphinxAtStartPar
The Counts tab shows an interactive overview of the corpora, with word
counts within corpora, books and subsets.

\end{itemize}

\sphinxstepscope


\chapter{References}
\label{\detokenize{references:references}}\label{\detokenize{references:id1}}\label{\detokenize{references::doc}}
\sphinxstepscope


\chapter{Cover Image Credits}
\label{\detokenize{image_sources:cover-image-credits}}\label{\detokenize{image_sources::doc}}
\sphinxAtStartPar
The following images were used for corpora on \sphinxhref{https://clic-fiction.com/}{clic\sphinxhyphen{}fiction.com}. All images are in the public domain.

\sphinxAtStartPar
\sphinxstylestrong{19C}: Sybil Tawse, illustration for \sphinxstyleemphasis{Cranford} by Elizabeth Gaskell (London: Adam and Charles Black, 1914), captioned “She always told them to hold out their tiny palms, into which she shook either peppermint or ginger lozenges.” Internet Archive: \sphinxhref{https://archive.org/details/cranford191400gaskuoft}{cranford191400gaskuoft}

\sphinxAtStartPar
\sphinxstylestrong{ArTS}: E.W. Kemble, \sphinxstyleemphasis{Striking for the Back Country}, in Mark Twain, \sphinxstyleemphasis{Adventures of Huckleberry Finn} (1885), p. 275. Wikimedia Commons: \sphinxhref{https://commons.wikimedia.org/wiki/File:Adventures\_of\_Huckleberry\_Finn\_1885-p275.png}{File:Adventures\_of\_Huckleberry\_Finn\_1885\sphinxhyphen{}p275.png}

\sphinxAtStartPar
\sphinxstylestrong{ChiLit}: John Tenniel, \sphinxstyleemphasis{The Mad Hatter’s Tea Party} (1865). Wikimedia Commons: \sphinxhref{https://commons.wikimedia.org/wiki/File:Alice\_par\_John\_Tenniel\_25.png}{File:Alice\_par\_John\_Tenniel\_25.png}

\sphinxAtStartPar
\sphinxstylestrong{DE19}: Ludwig Richter, \sphinxstyleemphasis{Der Brautzug im Frühling}. Wikimedia Commons: \sphinxhref{https://commons.wikimedia.org/wiki/File:Adrian\_Ludwig\_Richter\_005FXD.jpg}{File:Adrian\_Ludwig\_Richter\_005FXD.jpg}

\sphinxAtStartPar
\sphinxstylestrong{DNov}: Robert William Buss (1804\textendash{}1875), \sphinxstyleemphasis{Dickens’s Dream}. Wikimedia Commons: \sphinxhref{https://commons.wikimedia.org/wiki/File:Dickens\_dream.jpg}{File:Dickens\_dream.jpg}

\sphinxstepscope


\chapter{Policies}
\label{\detokenize{policies:policies}}\label{\detokenize{policies::doc}}
\sphinxstepscope


\section{Accessibility}
\label{\detokenize{policies/accessibility:accessibility}}\label{\detokenize{policies/accessibility::doc}}
\sphinxAtStartPar
Public authorities are obliged to make websites and mobile applications accessible in accordance with Directive (EU) 2016/2102 of the European Parliament and of the Council. The following accessibility statement documents the implementation status of this website in accordance with the applicable legislation.

\sphinxAtStartPar
We aim to make the CLiC website as easy to use and as accessible as possible for all users. Our aim is to meet the recommended government standard for web accessibility. This means that you should be able to:
\begin{itemize}
\item {} 
\sphinxAtStartPar
change the colours, contrast levels, and fonts

\item {} 
\sphinxAtStartPar
resize the text to make it easier to read

\item {} 
\sphinxAtStartPar
navigate the site using just a keyboard or using speech recognition software

\item {} 
\sphinxAtStartPar
listen to the content using a screen reader

\end{itemize}

\sphinxAtStartPar
If you wish to contact us about the accessibility of the CLiC website please email us at \sphinxhref{mailto:dhss-kontakt@fau.de}{dhss\sphinxhyphen{}kontakt@fau.de}

\sphinxAtStartPar
Enforcement procedures
If a request via the contact option remains wholly or partially unanswered within six weeks, the competent supervisory authority shall examine, at the request of the user, whether measures are necessary in the context of monitoring vis\sphinxhyphen{}à\sphinxhyphen{}vis the obligated party.

\sphinxAtStartPar
Contact information of the enforcement oversight body:

\sphinxAtStartPar
Landesamt für Digitalisierung, Breitband und Vermessung

\sphinxAtStartPar
IT\sphinxhyphen{}Dienstleistungszentrum des Freistaats Bayern \sphinxhyphen{} Durchsetzungs\sphinxhyphen{} und Überwachungsstelle für barrierefreie Informationstechnik
St.\sphinxhyphen{}Martin\sphinxhyphen{}Straße 47
81541 München
EMail: \sphinxhref{mailto:bitv@bayern.de}{bitv@bayern.de}
URL: \sphinxurl{https://www.ldbv.bayern.de/digitalisierung/bitv/}

\sphinxAtStartPar
This statement was last reviewed on 9th December 2025.

\sphinxstepscope


\section{Cookies}
\label{\detokenize{policies/cookies:cookies}}\label{\detokenize{policies/cookies::doc}}

\subsection{Cookies}
\label{\detokenize{policies/cookies:id1}}
\sphinxAtStartPar
Description and scope of data processing
Our website uses cookies. Cookies are text files that are saved in the user’s web browser or by the web browser on the user’s computer system. When a user accesses a website, a cookie can be stored in the user’s operating system. This cookie contains a character string that allows the unique identification of the browser when the website is accessed again.

\sphinxAtStartPar
We use cookies to make our website more user\sphinxhyphen{}friendly. Some parts of our website require that the requesting browser can also be identified after changing pages.
During this process, the following data are stored in the cookies and transmitted:
\begin{itemize}
\item {} 
\sphinxAtStartPar
Search preferences

\item {} 
\sphinxAtStartPar
Technical measures are taken to pseudonymise user data collected in this way. This means that the data can no longer be assigned to the user. The data are not stored together with other personal data of the user.

\item {} 
\sphinxAtStartPar
All cookies are used for technical reasons and are used only in the situations described above. If there are additional applications in the website, who need to set other cookies, they are described in the sections about the applications below.

\end{itemize}

\sphinxAtStartPar
Legal basis for data processing
The legal basis for the temporary storage of data and logfiles is Art. 6 Abs. lit. f DSGVO.

\sphinxAtStartPar
Purpose of data processing
Analysis cookies are used for the purpose of improving the quality of our website and its content. We learn through the analysis cookies how the website is used and in this way can continuously optimise our web presence.

\sphinxAtStartPar
Storage period, options for filing an objection or requesting removal
As cookies are stored on the user’s computer and are transmitted from it to our website, users have full control over the use of cookies. You can disable or restrict the transmission of cookies by changing the cookie settings in your web browser or by clicking the cookie preferences link below. Cookies that are already stored can be deleted at any time. This can also be done automatically. If cookies are deactivated for our website, it may be the case that not all of the website’s functions can be used in full.

\sphinxAtStartPar
Cookie Preferences


\begin{savenotes}\sphinxattablestart
\sphinxthistablewithglobalstyle
\centering
\begin{tabulary}{\linewidth}[t]{TTTT}
\sphinxtoprule
\sphinxtableatstartofbodyhook
\sphinxAtStartPar
Provider
&
\sphinxAtStartPar
Cookie
&
\sphinxAtStartPar
Purpose
&
\sphinxAtStartPar
Expiry
\\
\sphinxhline
\sphinxAtStartPar
Google Analytics
Google Analytics
Google Analytics
This website
&
\sphinxAtStartPar
\_ga
\_gid
\_gat
cookieMessageApprove
&
\sphinxAtStartPar
This helps us count how many people visit our website by tracking if you’ve visited before
This helps us count how many people visit our website by tracking if you’ve visited before
Used to manage the rate at which page view requests are made
Used to know if you’ve previously hidden the cookies information messages on our site
&
\sphinxAtStartPar
2 years
24 hours
10 minutes
2039\sphinxhyphen{}11\sphinxhyphen{}13
\\
\sphinxbottomrule
\end{tabulary}
\sphinxtableafterendhook\par
\sphinxattableend\end{savenotes}

\sphinxstepscope


\section{Legal}
\label{\detokenize{policies/legal:legal}}\label{\detokenize{policies/legal::doc}}

\subsection{Website owner}
\label{\detokenize{policies/legal:website-owner}}
\sphinxAtStartPar
Friedrich\sphinxhyphen{}Alexander\sphinxhyphen{}Universität Erlangen\sphinxhyphen{}Nürnberg

\sphinxAtStartPar
Freyeslebenstraße 1
91058 Erlangen
E\sphinxhyphen{}Mail: \sphinxhref{mailto:poststelle@fau.de}{poststelle@fau.de}

\sphinxAtStartPar
\#Legal form and representation
Körperschaft des Öffentlichen Rechts. Legally represented by Präsident Prof. Dr. Joachim Hornegger.
The responsibility for contributions marked by name or with a separate imprint lies with the respective authors.


\subsection{Responsibility for content}
\label{\detokenize{policies/legal:responsibility-for-content}}
\sphinxAtStartPar
Prof. Dr. Michaela Mahlberg

\sphinxAtStartPar
Nürnberger Str. 74
91052 Erlangen
EMail: \sphinxhref{mailto:michaela.mahlberg@fau.de}{michaela.mahlberg@fau.de}
Editorial and technical contacts

\sphinxAtStartPar
EMail: \sphinxhref{mailto:dhss-kontakt@fau.de}{dhss\sphinxhyphen{}kontakt@fau.de}
Responsible controlling authority
Bayerisches Staatsministerium für Wissenschaft und Kunst

\sphinxAtStartPar
Salvatorstraße 2
80327 München
URL: \sphinxurl{https://www.stmwk.bayern.de/}
Telefon: +49 89 2186 0
E\sphinxhyphen{}Mail: \sphinxhref{mailto:poststelle@ldbv.bayern.de}{poststelle@ldbv.bayern.de}
ID numbers
VAT identification number       DE 132507686
Tax number      216/114/20045 (Finanzamt Erlangen)
DUNS number     327958716
EORI number     DE4204891
Bank details:   Bayerische Landesbank München
IBAN: DE66700500000301279280
SWIFT/BIC\sphinxhyphen{}Code: BYLADEMMXXX
Report misuse of computers and network resources
Contact point for reporting misuse of computers and network resources.

\sphinxAtStartPar
IT\sphinxhyphen{}Sicherheit
URL: \sphinxurl{https://www.rrze.fau.de/abuse/}
Postanschrift:
IT\sphinxhyphen{}Sicherheit
Martensstraße 1
91058 Erlangen


\subsection{Description and scope of data processing}
\label{\detokenize{policies/legal:description-and-scope-of-data-processing}}
\sphinxAtStartPar
Contact forms are available on our website that can be used to contact us electronically. If a user makes use of this possibility, the data they enter in the input form are transmitted to us and stored. The contact forms list and explain which data is required. The contact forms indicate if there are any deviations from or additions to the principles, purpose and duration of storage as presented here.


\subsection{Legal basis for data processing}
\label{\detokenize{policies/legal:legal-basis-for-data-processing}}
\sphinxAtStartPar
The legal basis for the processing of the data transmitted in the course of sending an email is Article 6 (1) lit. e DSGVO i.V.m. Art. 4 and 5 BayDSG for the fulfilment of the tasks of \S{} 5 TMG, Art. 3 para. 1 BayEGovG and \S{} 2 BayBITV

\sphinxAtStartPar
If the email contact aims to conclude a contract, then additional legal basis for the processing is Art. 6 para. 1 lit. b DSGVO.


\subsection{Purpose of data processing}
\label{\detokenize{policies/legal:purpose-of-data-processing}}
\sphinxAtStartPar
The personal data from the input form are processed solely for the purpose of contacting the user. If the user contacts us by email, this also constitutes our legitimate interests in processing the data.
All other personal data processed during the dispatch of an email serve to prevent misuse of the contact form and to ensure that our IT systems are secure.


\subsection{Storage period}
\label{\detokenize{policies/legal:storage-period}}
\sphinxAtStartPar
Data are deleted as soon as they are no longer necessary for fulfilling the purpose for which they were collected. This is the case for the personal data from the input template of the contact form and those data sent by email when the respective conversation with the user has ended. The conversation is regarded to have ended when it can be seen from the circumstances that the subject matter in question has been conclusively settled.


\subsection{Options for filing an objection}
\label{\detokenize{policies/legal:options-for-filing-an-objection}}
\sphinxAtStartPar
For reasons that arise from your particular situation, you may also object to the processing of personal data relating to us by us at any time (Art. 21 GDPR). If the legal requirements are met, we will no longer process your personal data in the following.


\subsection{Obligation to provide}
\label{\detokenize{policies/legal:obligation-to-provide}}
\sphinxAtStartPar
Insofar as the personal data required for the performance of the contract is not specified, this is not possible for us.


\subsection{Siteimprove Analytics}
\label{\detokenize{policies/legal:siteimprove-analytics}}
\sphinxAtStartPar
Description and scope of data processing
This website uses Siteimprove Analytics, a web analytics service provided by Siteimprove. Siteimprove Analytics uses „cookies“, which are text files placed on your computer, to help the FAU analyze how visitors use the site. The information generated by the cookies about the visitors’ use of the website will be stored and processed by Siteimprove on servers in Denmark.

\sphinxAtStartPar
IP addresses are anonymized irreversibly before data is made available in the Siteimprove Analytics or Intelligence Suite for the FAU.

\sphinxAtStartPar
The FAU will use this information for evaluating the visitors’ use of the website, compiling reports on website activity, and ultimately for improving the website experience for its visitors. Siteimprove will not transmit this information to third parties or use it for any marketing or advertising purposes.


\subsection{These are the cookies used by Siteimprove Analytics on this website:}
\label{\detokenize{policies/legal:these-are-the-cookies-used-by-siteimprove-analytics-on-this-website}}
\sphinxAtStartPar
Cookie name: nmstat
Cookie type: Persistent \sphinxhyphen{} expires after 1000 days
About: This cookie is used to help record visitors’ use of the website. It is used to collect statistics about site usage such as when the visitor last visited the site. The cookie contains no personal information and is used only for web analytics.
Cookie name: AWSELB / AWSELBCOR
Cookie type: Session cookie
About: The AWSELB cookie ensures that all page views for the same visit (user session) are sent to the same endpoint. This enables us to determine the sequence of a user’s page views needed for features like Behavior Tracking and Funnels.
By using this website, the visitor consents to the processing of data about him/her by Siteimprove in the manner and for the purposes set out above.

\sphinxAtStartPar
Legal basis for the processing of personal data
The legal basis for the processing of personal data using cookies is Art. 6 (1) (e) GDPR in relation with Art. 4 BayDSG, especially the specification under \S{} 15 (3) TMG and Art. 10 BayHSchG.


\subsection{SSL encryption}
\label{\detokenize{policies/legal:ssl-encryption}}
\sphinxAtStartPar
Our website uses SSL encryption for security reasons and to protect the transmission of confidential information, for example enquiries you send to us as operators of the website. You can recognise an encrypted connection when the browser’s address line changes from \sphinxurl{http://} to \sphinxurl{https://} and a padlock appears in your web browser.

\sphinxAtStartPar
If SSL encryption is activated, the data you transmit to us cannot be read by third parties.


\subsection{Rights of the data subject}
\label{\detokenize{policies/legal:rights-of-the-data-subject}}
\sphinxAtStartPar
With regard to the processing of your personal data, you as a data subject are entitled to the following rights pursuant to Art. 15 et seq. GDPR:

\sphinxAtStartPar
You can request information as to whether we process your personal data. If this is the case, you have the right to information about this personal data as well as further information in connection with the processing (Art. 15 GDPR). Please note that this right of access may be restricted or excluded in certain cases (cf. in particular Art. 10 BayDSG).
In the event that personal data about you is (no longer) accurate or incomplete, you may request that this data be corrected and, if necessary, completed (Art. 16 GDPR).
If the legal requirements are met, you can demand that your personal data be erased (Art. 17 GDPR) or that the processing of this data be restricted (Art. 18 DSGVO). However, the right to erasure pursuant to Art. 17 (1) and (2) GDPR does not apply, inter alia, if the processing of personal data is necessary for the performance of a task carried out in the public interest or in the exercise of official authority or in the exercise of official authority vested (Art. 17 para. 3 letter b GDPR).
If you have given your consent to the processing, you have the right to withdrawal it at any time. The withdrawal will only take effect in the future; this means that the withdrawal does not affect the lawfulness of the processing operations carried out on the basis of the consent up to the withdrawal.
For reasons arising from your particular situation, you may also object to the processing of your personal data by us at any time (Art. 21 GDPR). If the legal requirements are met, we will subsequently no longer process your personal data.
Insofar as you have consented to the processing of your personal data or have agreed to the performance of the contract and the data processing is carried out automated , you may be entitled to data portability (Art. 20 GDPR).
You have the right to lodge a complaint to a supervisory authority within the meaning of Art. 51 GDPR about the processing of your personal data. The responsible supervisory authority for Bavarian public authorities is the Bavarian Data Protection Commissioner, Wagmüllerstraße 18, 80538 Munich.

\sphinxstepscope


\section{Privacy}
\label{\detokenize{policies/privacy:privacy}}\label{\detokenize{policies/privacy::doc}}

\subsection{Name and address of the Data Protection Officer}
\label{\detokenize{policies/privacy:name-and-address-of-the-data-protection-officer}}
\sphinxAtStartPar
Datenschutzbeauftragter FAU
Klaus Hoogestraat

\sphinxAtStartPar
Postal address
Datenschutzbeauftragter FAU
c/o ITM Gesellschaft für IT\sphinxhyphen{}Management mbH
Bürgerstraße 81
01127 Dresden
Phone: +49 9131 85\sphinxhyphen{}25860
EMail: \sphinxhref{mailto:datenschutzbeauftragter@fau.de}{datenschutzbeauftragter@fau.de}


\subsection{General information on data processing}
\label{\detokenize{policies/privacy:general-information-on-data-processing}}

\subsection{Scope of processing of personal data}
\label{\detokenize{policies/privacy:scope-of-processing-of-personal-data}}
\sphinxAtStartPar
We only process our users’ personal data to the extent necessary to provide services, content and a functional website. As a rule, personal data are only processed after the user gives their consent. An exception applies in those cases where it is impractical to obtain the user’s prior consent and the processing of such data is permitted by law.


\subsection{Legal basis for the processing of personal data}
\label{\detokenize{policies/privacy:legal-basis-for-the-processing-of-personal-data}}
\sphinxAtStartPar
Art. 6 (1) (a) of the EU General Data Protection Regulation (GDPR) forms the legal basis for us to obtain the consent of a data subject for their personal data to be processed.
When processing personal data required for the performance of a contract in which the contractual party is the data subject, Art. 6 (1) (b) GDPR forms the legal basis. This also applies if data has to be processed in order to carry out pre\sphinxhyphen{}contractual activities.
Art. 6 (1) (c) GDPR forms the legal basis if personal data has to be processed in order to fulfil a legal obligation on the part of our organisation.
Art. 6 (1) (d) GDPR forms the legal basis in the case that vital interests of the data subject or another natural person make the processing of personal data necessary.
If data processing is necessary in order to protect the legitimate interests of our organisation or of a third party and if the interests, basic rights and fundamental freedoms of the data subject do not outweigh the interests mentioned above, Art. 6 (1) (f) GDPR forms the legal basis for such data processing.


\subsection{Deletion of data and storage period}
\label{\detokenize{policies/privacy:deletion-of-data-and-storage-period}}
\sphinxAtStartPar
The personal data of the data subject are deleted or blocked as soon as the reason for storing them ceases to exist. Storage beyond this time period may occur if provided for by European or national legislators in directives under Union legislation, laws or other regulations to which the data controller is subject. Such data are also blocked or deleted if a storage period prescribed by one of the above\sphinxhyphen{}named rules expires, unless further storage of the data is necessary for entering into or performing a contract.


\subsection{Provision of the website and generation of log files}
\label{\detokenize{policies/privacy:provision-of-the-website-and-generation-of-log-files}}

\subsection{Description and scope of data processing}
\label{\detokenize{policies/privacy:description-and-scope-of-data-processing}}
\sphinxAtStartPar
Each time our website is accessed, our system automatically collects data and information from the user’s computer system.
In this context, the following data are collected:
\begin{itemize}
\item {} 
\sphinxAtStartPar
Address (URL) of the website from which the file was requested

\item {} 
\sphinxAtStartPar
Name of the retrieved file

\item {} 
\sphinxAtStartPar
Date and time of the request

\item {} 
\sphinxAtStartPar
Data volume transmitted

\item {} 
\sphinxAtStartPar
Access status (file transferred, file not found, etc.)

\item {} 
\sphinxAtStartPar
Description of the type of web browser and/or operating system used

\item {} 
\sphinxAtStartPar
Anonymised IP address of the requesting computer

\end{itemize}

\sphinxAtStartPar
The data stored are required exclusively for technical or statistical purposes; no comparison with other data or disclosure to third parties occurs, not even in part. The data are stored in our system’s log files. This is not the case for the user’s IP addresses or other data that make it possible to assign the data to a specific user: before data are stored, each dataset is anonymised by changing the IP address. These data are not stored together with other personal data .


\subsection{Legal basis for data processing}
\label{\detokenize{policies/privacy:legal-basis-for-data-processing}}
\sphinxAtStartPar
The legal basis for the temporary storage of data and logfiles is Art. 6 Abs. lit. f DSGVO.


\subsection{Purpose of data processing}
\label{\detokenize{policies/privacy:purpose-of-data-processing}}
\sphinxAtStartPar
The temporary storage of the IP address by the system is necessary in order to deliver the website to the user’s computer. For this purpose, the user’s IP address must remain stored for the duration of the session.
The storage of such data in log files takes place in order to ensure the website’s functionality. These data also serve to help us optimise the website and ensure that our IT systems are secure. They are not evaluated for marketing purposes in this respect.


\subsection{Storage period}
\label{\detokenize{policies/privacy:storage-period}}
\sphinxAtStartPar
Data are deleted as soon as they are no longer necessary for fulfilling the purpose for which they were collected. If data have been collected for the purpose of providing the website, they are deleted at the end of the respective session.
If data are stored in log files, they are deleted at the latest after seven days. A longer storage period is possible. In this case, the users’ IP addresses are deleted or masked so that they can no longer be assigned to the client accessing the website.


\subsection{Options for filing an objection or requesting removal}
\label{\detokenize{policies/privacy:options-for-filing-an-objection-or-requesting-removal}}
\sphinxAtStartPar
The collection of data for the purpose of providing the website and the storage of such data in log files is essential to the website’s operation. As a consequence, the user has no possibility to object.


\subsection{Cookies}
\label{\detokenize{policies/privacy:cookies}}

\subsection{Description and scope of data processing}
\label{\detokenize{policies/privacy:id1}}
\sphinxAtStartPar
Our website uses cookies. Cookies are text files that are saved in the user’s web browser or by the web browser on the user’s computer system. When a user accesses a website, a cookie can be stored in the user’s operating system. This cookie contains a character string that allows the unique identification of the browser when the website is accessed again.

\sphinxAtStartPar
We use cookies to make our website more user\sphinxhyphen{}friendly. Some parts of our website require that the requesting browser can also be identified after changing pages.
During this process, the following data are stored in the cookies and transmitted:
\begin{itemize}
\item {} 
\sphinxAtStartPar
Log\sphinxhyphen{}in information (only in the case of protected information that is made available exclusively to FAU members)

\item {} 
\sphinxAtStartPar
Search preferences (from October 2018)

\item {} 
\sphinxAtStartPar
Technical measures are taken to pseudonymise user data collected in this way. This means that the data can no longer be assigned to the user. The data are not stored together with other personal data of the user.

\end{itemize}

\sphinxAtStartPar
All cookies are used for technical reasons and are used only in the situations described above. If there are additional applications in the website, who need to set other cookies, they are described in the sections about the applications below.


\subsection{Legal basis for data processing}
\label{\detokenize{policies/privacy:id2}}
\sphinxAtStartPar
The legal basis for the temporary storage of data and logfiles is Art. 6 Abs. lit. f DSGVO.


\subsection{Purpose of data processing}
\label{\detokenize{policies/privacy:id3}}
\sphinxAtStartPar
Analysis cookies are used for the purpose of improving the quality of our website and its content. We learn through the analysis cookies how the website is used and in this way can continuously optimise our web presence.


\subsection{Storage period, options for filing an objection or requesting removal}
\label{\detokenize{policies/privacy:storage-period-options-for-filing-an-objection-or-requesting-removal}}
\sphinxAtStartPar
As cookies are stored on the user’s computer and are transmitted from it to our website, users have full control over the use of cookies. You can disable or restrict the transmission of cookies by changing the cookie settings in your web browser or by clicking the cookie preferences link below. Cookies that are already stored can be deleted at any time. This can also be done automatically. If cookies are deactivated for our website, it may be the case that not all of the website’s functions can be used in full.


\subsection{Contact form and contact by email}
\label{\detokenize{policies/privacy:contact-form-and-contact-by-email}}

\subsection{Description and scope of data processing}
\label{\detokenize{policies/privacy:id4}}
\sphinxAtStartPar
Contact forms are available on our website that can be used to contact us electronically. If a user makes use of this possibility, the data they enter in the input form are transmitted to us and stored. The contact forms list and explain which data is required. The contact forms indicate if there are any deviations from or additions to the principles, purpose and duration of storage as presented here.


\subsection{Legal basis for data processing}
\label{\detokenize{policies/privacy:id5}}
\sphinxAtStartPar
The legal basis for the processing of the data transmitted in the course of sending an email is Article 6 (1) lit. e DSGVO i.V.m. Art. 4 and 5 BayDSG for the fulfilment of the tasks of \S{} 5 TMG, Art. 3 para. 1 BayEGovG and \S{} 2 BayBITV

\sphinxAtStartPar
If the email contact aims to conclude a contract, then additional legal basis for the processing is Art. 6 para. 1 lit. b DSGVO.


\subsection{Purpose of data processing}
\label{\detokenize{policies/privacy:id6}}
\sphinxAtStartPar
The personal data from the input form are processed solely for the purpose of contacting the user. If the user contacts us by email, this also constitutes our legitimate interests in processing the data.
All other personal data processed during the dispatch of an email serve to prevent misuse of the contact form and to ensure that our IT systems are secure.


\subsection{Storage period}
\label{\detokenize{policies/privacy:id7}}
\sphinxAtStartPar
Data are deleted as soon as they are no longer necessary for fulfilling the purpose for which they were collected. This is the case for the personal data from the input template of the contact form and those data sent by email when the respective conversation with the user has ended. The conversation is regarded to have ended when it can be seen from the circumstances that the subject matter in question has been conclusively settled.


\subsection{Options for filing an objection}
\label{\detokenize{policies/privacy:options-for-filing-an-objection}}
\sphinxAtStartPar
For reasons that arise from your particular situation, you may also object to the processing of personal data relating to us by us at any time (Art. 21 GDPR). If the legal requirements are met, we will no longer process your personal data in the following.


\subsection{Obligation to provide}
\label{\detokenize{policies/privacy:obligation-to-provide}}
\sphinxAtStartPar
Insofar as the personal data required for the performance of the contract is not specified, this is not possible for us.


\subsection{External Service Providers}
\label{\detokenize{policies/privacy:external-service-providers}}

\subsection{Siteimprove Analytics}
\label{\detokenize{policies/privacy:siteimprove-analytics}}
\sphinxAtStartPar
Usage of Siteimprove Analytics


\subsection{Description and scope of data processing}
\label{\detokenize{policies/privacy:id8}}
\sphinxAtStartPar
This website uses Siteimprove Analytics, a web analytics service provided by Siteimprove. Siteimprove Analytics uses „cookies“, which are text files placed on your computer, to help the FAU analyze how visitors use the site. The information generated by the cookies about the visitors’ use of the website will be stored and processed by Siteimprove on servers in Denmark.

\sphinxAtStartPar
IP addresses are anonymized irreversibly before data is made available in the Siteimprove Analytics or Intelligence Suite for the FAU.

\sphinxAtStartPar
The FAU will use this information for evaluating the visitors’ use of the website, compiling reports on website activity, and ultimately for improving the website experience for its visitors. Siteimprove will not transmit this information to third parties or use it for any marketing or advertising purposes.

\sphinxAtStartPar
These are the cookies used by Siteimprove Analytics on this website:

\sphinxAtStartPar
Cookie name: nmstat
Cookie type: Persistent \sphinxhyphen{} expires after 1000 days
About: This cookie is used to help record visitors’ use of the website. It is used to collect statistics about site usage such as when the visitor last visited the site. The cookie contains no personal information and is used only for web analytics.
Cookie name: AWSELB / AWSELBCOR
Cookie type: Session cookie
About: The AWSELB cookie ensures that all page views for the same visit (user session) are sent to the same endpoint. This enables us to determine the sequence of a user’s page views needed for features like Behavior Tracking and Funnels.
By using this website, the visitor consents to the processing of data about him/her by Siteimprove in the manner and for the purposes set out above.


\subsection{Legal basis for the processing of personal data}
\label{\detokenize{policies/privacy:id9}}
\sphinxAtStartPar
The legal basis for the processing of personal data using cookies is Art. 6 (1) (e) GDPR in relation with Art. 4 BayDSG, especially the specification under \S{} 15 (3) TMG and Art. 10 BayHSchG.


\subsection{Objection and removal option}
\label{\detokenize{policies/privacy:objection-and-removal-option}}

\subsection{SSL encryption}
\label{\detokenize{policies/privacy:ssl-encryption}}
\sphinxAtStartPar
Our website uses SSL encryption for security reasons and to protect the transmission of confidential information, for example enquiries you send to us as operators of the website. You can recognise an encrypted connection when the browser’s address line changes from \sphinxurl{http://} to \sphinxurl{https://} and a padlock appears in your web browser.

\sphinxAtStartPar
If SSL encryption is activated, the data you transmit to us cannot be read by third parties.


\subsection{Rights of the data subject}
\label{\detokenize{policies/privacy:rights-of-the-data-subject}}
\sphinxAtStartPar
With regard to the processing of your personal data, you as a data subject are entitled to the following rights pursuant to Art. 15 et seq. GDPR:
\begin{itemize}
\item {} 
\sphinxAtStartPar
You can request information as to whether we process your personal data. If this is the case, you have the right to information about this personal data as well as further information in connection with the processing (Art. 15 GDPR). Please note that this right of access may be restricted or excluded in certain cases (cf. in particular Art. 10 BayDSG).

\item {} 
\sphinxAtStartPar
In the event that personal data about you is (no longer) accurate or incomplete, you may request that this data be corrected and, if necessary, completed (Art. 16 GDPR).

\item {} 
\sphinxAtStartPar
If the legal requirements are met, you can demand that your personal data be erased (Art. 17 GDPR) or that the processing of this data be restricted (Art. 18 DSGVO). However, the right to erasure pursuant to Art. 17 (1) and (2) GDPR does not apply, inter alia, if the processing of personal data is necessary for the performance of a task carried out in the public interest or in the exercise of official authority or in the exercise of official authority vested (Art. 17 para. 3 letter b GDPR).

\item {} 
\sphinxAtStartPar
If you have given your consent to the processing, you have the right to withdrawal it at any time. The withdrawal will only take effect in the future; this means that the withdrawal does not affect the lawfulness of the processing operations carried out on the basis of the consent up to the withdrawal.

\item {} 
\sphinxAtStartPar
For reasons arising from your particular situation, you may also object to the processing of your personal data by us at any time (Art. 21 GDPR). If the legal requirements are met, we will subsequently no longer process your personal data.

\item {} 
\sphinxAtStartPar
Insofar as you have consented to the processing of your personal data or have agreed to the performance of the contract and the data processing is carried out automated , you may be entitled to data portability (Art. 20 GDPR).

\item {} 
\sphinxAtStartPar
You have the right to lodge a complaint to a supervisory authority within the meaning of Art. 51 GDPR about the processing of your personal data. The responsible supervisory authority for Bavarian public authorities is the Bavarian Data Protection Commissioner, Wagmüllerstraße 18, 80538 Munich.

\end{itemize}

\sphinxstepscope


\section{Freedom of Information}
\label{\detokenize{policies/freedomofinformation:freedom-of-information}}\label{\detokenize{policies/freedomofinformation::doc}}
\sphinxAtStartPar
CLiC follows the \sphinxhref{https://www.birmingham.ac.uk/university/governance/policies-regs/information/index.aspx}{University of Birmingham’s freedom on information policy}

\sphinxstepscope


\chapter{Glossary of terms used}
\label{\detokenize{glossary:glossary-of-terms-used}}\label{\detokenize{glossary::doc}}\begin{description}
\sphinxlineitem{book:}
\sphinxAtStartPar
A UTF\sphinxhyphen{}8 encoded text file that has been imported into CLiC, invariably from the corpora repository.

\sphinxlineitem{region:}
\sphinxAtStartPar
A labelled portion of a book. Each region will have:
\begin{itemize}
\item {} 
\sphinxAtStartPar
A name (or rclass), for example ‘chapter.sentence’. For all possible names, see \sphinxurl{https://github.com/mahlberg-lab/clic/blob/HEAD//schema/10-rclass.sql}.

\item {} 
\sphinxAtStartPar
A start and end character position within the full book text

\item {} 
\sphinxAtStartPar
(optionally) a number (or rvalue), for example it’s position within a chapter.

\end{itemize}

\sphinxAtStartPar
Regions are added by \sphinxstyleemphasis{region tagger} scripts, which are in \sphinxcode{\sphinxupquote{clic.region}}.

\sphinxlineitem{token:}
\sphinxAtStartPar
A token is a labelled portion of a book that contains a single word.
In the phrase \sphinxcode{\sphinxupquote{\textquotesingle{}For \_more\_!\textquotesingle{} said Mr. Limbkins.}}, \sphinxcode{\sphinxupquote{For}}, \sphinxcode{\sphinxupquote{more}}, \sphinxcode{\sphinxupquote{said}}, \sphinxcode{\sphinxupquote{Mr}}, \sphinxcode{\sphinxupquote{Limbkins}} would be tokens.

\sphinxAtStartPar
See \sphinxcode{\sphinxupquote{clic.tokenizer}}.

\sphinxlineitem{type/ttype:}
\sphinxAtStartPar
A token has a type (thus ttype). This is a normalised form of the token.
The tokens in the phrase \sphinxcode{\sphinxupquote{\textquotesingle{}For \_more\_!\textquotesingle{} said Mr. Limbkins.}} would have types \sphinxcode{\sphinxupquote{for}}, \sphinxcode{\sphinxupquote{more}}, \sphinxcode{\sphinxupquote{said}}, \sphinxcode{\sphinxupquote{mr}}, \sphinxcode{\sphinxupquote{limbkins}}.

\sphinxAtStartPar
See \sphinxcode{\sphinxupquote{clic.tokenizer}}.

\end{description}

\appendix
\phantomsection\label{\detokenize{appendices::doc}}

\chapter{Appendices}
\label{\detokenize{appendices:appendices}}\label{\detokenize{appendices:a-list-of-texts-available-in-clic}}
\sphinxAtStartPar
The appendices only act as an overview. For up\sphinxhyphen{}to\sphinxhyphen{}date information and
the documentation of how the corpus texts were cleaned, please refer to
the README in the \sphinxcite{references:github-corpora} repository. You can also create a
dynamic list of the books in the corpora (and their word counts) using
the {\hyperref[\detokenize{clicanalysis/counts:counts}]{\sphinxcrossref{\DUrole{std,std-ref}{Counts}}}} tab.
\subsubsection*{Texts in CLiC}


\begin{savenotes}\sphinxattablestart
\sphinxthistablewithglobalstyle
\centering
\sphinxcapstartof{table}
\sphinxthecaptionisattop
\sphinxcaption{DNov \textendash{} Dickens’s Novels (15 texts)}\label{\detokenize{appendices:id2}}
\sphinxaftertopcaption
\begin{tabular}[t]{\X{60}{100}\X{40}{100}}
\sphinxtoprule
\sphinxstyletheadfamily 
\sphinxAtStartPar
Title
&\sphinxstyletheadfamily 
\sphinxAtStartPar
Author
\\
\sphinxmidrule
\sphinxtableatstartofbodyhook
\sphinxAtStartPar
A Tale of Two Cities
&
\sphinxAtStartPar
Charles Dickens
\\
\sphinxhline
\sphinxAtStartPar
Barnaby Rudge
&
\sphinxAtStartPar
Charles Dickens
\\
\sphinxhline
\sphinxAtStartPar
Bleak House
&
\sphinxAtStartPar
Charles Dickens
\\
\sphinxhline
\sphinxAtStartPar
David Copperfield
&
\sphinxAtStartPar
Charles Dickens
\\
\sphinxhline
\sphinxAtStartPar
Dombey and Son
&
\sphinxAtStartPar
Charles Dickens
\\
\sphinxhline
\sphinxAtStartPar
Great Expectations
&
\sphinxAtStartPar
Charles Dickens
\\
\sphinxhline
\sphinxAtStartPar
Hard Times
&
\sphinxAtStartPar
Charles Dickens
\\
\sphinxhline
\sphinxAtStartPar
Little Dorrit
&
\sphinxAtStartPar
Charles Dickens
\\
\sphinxhline
\sphinxAtStartPar
Martin Chuzzlewit
&
\sphinxAtStartPar
Charles Dickens
\\
\sphinxhline
\sphinxAtStartPar
Nicholas Nickleby
&
\sphinxAtStartPar
Charles Dickens
\\
\sphinxhline
\sphinxAtStartPar
Oliver Twist
&
\sphinxAtStartPar
Charles Dickens
\\
\sphinxhline
\sphinxAtStartPar
Our Mutual Friend
&
\sphinxAtStartPar
Charles Dickens
\\
\sphinxhline
\sphinxAtStartPar
Pickwick Papers
&
\sphinxAtStartPar
Charles Dickens
\\
\sphinxhline
\sphinxAtStartPar
The Mystery of Edwin Drood
&
\sphinxAtStartPar
Charles Dickens
\\
\sphinxhline
\sphinxAtStartPar
The Old Curiosity Shop
&
\sphinxAtStartPar
Charles Dickens
\\
\sphinxbottomrule
\end{tabular}
\sphinxtableafterendhook\par
\sphinxattableend\end{savenotes}


\begin{savenotes}\sphinxattablestart
\sphinxthistablewithglobalstyle
\centering
\sphinxcapstartof{table}
\sphinxthecaptionisattop
\sphinxcaption{19C \textendash{} 19th Century Reference Corpus (29 texts)}\label{\detokenize{appendices:id3}}
\sphinxaftertopcaption
\begin{tabular}[t]{\X{60}{100}\X{40}{100}}
\sphinxtoprule
\sphinxstyletheadfamily 
\sphinxAtStartPar
Title
&\sphinxstyletheadfamily 
\sphinxAtStartPar
Author
\\
\sphinxmidrule
\sphinxtableatstartofbodyhook
\sphinxAtStartPar
Agnes Grey
&
\sphinxAtStartPar
Anne Brontë
\\
\sphinxhline
\sphinxAtStartPar
Antonina or, the Fall of Rome
&
\sphinxAtStartPar
Wilkie Collins
\\
\sphinxhline
\sphinxAtStartPar
Armadale
&
\sphinxAtStartPar
Wilkie Collins
\\
\sphinxhline
\sphinxAtStartPar
Cranford
&
\sphinxAtStartPar
Elizabeth Gaskell
\\
\sphinxhline
\sphinxAtStartPar
Daniel Deronda
&
\sphinxAtStartPar
George Eliot
\\
\sphinxhline
\sphinxAtStartPar
Dracula
&
\sphinxAtStartPar
Bram Stoker
\\
\sphinxhline
\sphinxAtStartPar
Emma
&
\sphinxAtStartPar
Jane Austen
\\
\sphinxhline
\sphinxAtStartPar
Frankenstein
&
\sphinxAtStartPar
Mary Shelley
\\
\sphinxhline
\sphinxAtStartPar
Jane Eyre
&
\sphinxAtStartPar
Charlotte Brontë
\\
\sphinxhline
\sphinxAtStartPar
Jude the Obscure
&
\sphinxAtStartPar
Thomas Hardy
\\
\sphinxhline
\sphinxAtStartPar
Lady Audley’s Secret
&
\sphinxAtStartPar
Mary Elizabeth Braddon
\\
\sphinxhline
\sphinxAtStartPar
Mary Barton
&
\sphinxAtStartPar
Elizabeth Gaskell
\\
\sphinxhline
\sphinxAtStartPar
North and South
&
\sphinxAtStartPar
Elizabeth Gaskell
\\
\sphinxhline
\sphinxAtStartPar
Persuasion
&
\sphinxAtStartPar
Jane Austen
\\
\sphinxhline
\sphinxAtStartPar
Pride and Prejudice
&
\sphinxAtStartPar
Jane Austen
\\
\sphinxhline
\sphinxAtStartPar
Sybil, or the two nations
&
\sphinxAtStartPar
Benjamin Disraeli
\\
\sphinxhline
\sphinxAtStartPar
Tess of the D’Urbervilles
&
\sphinxAtStartPar
Thomas Hardy
\\
\sphinxhline
\sphinxAtStartPar
The Hound of the Baskervilles
&
\sphinxAtStartPar
Arthur Conan Doyle
\\
\sphinxhline
\sphinxAtStartPar
The Last Days of Pompeii
&
\sphinxAtStartPar
Edward George Bulwer\sphinxhyphen{}Lytton
\\
\sphinxhline
\sphinxAtStartPar
The Mill on the Floss
&
\sphinxAtStartPar
George Eliot
\\
\sphinxhline
\sphinxAtStartPar
The Picture of Dorian Gray
&
\sphinxAtStartPar
Oscar Wilde
\\
\sphinxhline
\sphinxAtStartPar
The Professor
&
\sphinxAtStartPar
Charlotte Brontë
\\
\sphinxhline
\sphinxAtStartPar
The Return of the Native
&
\sphinxAtStartPar
Thomas Hardy
\\
\sphinxhline
\sphinxAtStartPar
The Small House at Allington
&
\sphinxAtStartPar
Anthony Trollope
\\
\sphinxhline
\sphinxAtStartPar
The Strange Case of Dr Jekyll and Mr Hyde
&
\sphinxAtStartPar
Robert Louis Stevenson
\\
\sphinxhline
\sphinxAtStartPar
The Woman in White
&
\sphinxAtStartPar
Wilkie Collins
\\
\sphinxhline
\sphinxAtStartPar
Vanity Fair
&
\sphinxAtStartPar
William Makepeace Thackeray
\\
\sphinxhline
\sphinxAtStartPar
Vivian Grey
&
\sphinxAtStartPar
Benjamin Disraeli
\\
\sphinxhline
\sphinxAtStartPar
Wuthering Heights
&
\sphinxAtStartPar
Emily Brontë
\\
\sphinxbottomrule
\end{tabular}
\sphinxtableafterendhook\par
\sphinxattableend\end{savenotes}


\begin{savenotes}
\sphinxatlongtablestart
\sphinxthistablewithglobalstyle
\makeatletter
  \LTleft \@totalleftmargin plus1fill
  \LTright\dimexpr\columnwidth-\@totalleftmargin-\linewidth\relax plus1fill
\makeatother
\begin{longtable}{\X{60}{100}\X{40}{100}}
\sphinxthelongtablecaptionisattop
\caption{ChiLit \textendash{} 19th Century Children’s Literature Corpus (71 texts)\strut}\label{\detokenize{appendices:id4}}\\*[\sphinxlongtablecapskipadjust]
\sphinxtoprule
\sphinxstyletheadfamily 
\sphinxAtStartPar
Title
&\sphinxstyletheadfamily 
\sphinxAtStartPar
Author
\\
\sphinxmidrule
\endfirsthead

\multicolumn{2}{c}{\sphinxnorowcolor
    \makebox[0pt]{\sphinxtablecontinued{\tablename\ \thetable{} \textendash{} continued from previous page}}%
}\\
\sphinxtoprule
\sphinxstyletheadfamily 
\sphinxAtStartPar
Title
&\sphinxstyletheadfamily 
\sphinxAtStartPar
Author
\\
\sphinxmidrule
\endhead

\sphinxbottomrule
\multicolumn{2}{r}{\sphinxnorowcolor
    \makebox[0pt][r]{\sphinxtablecontinued{continues on next page}}%
}\\
\endfoot

\endlastfoot
\sphinxtableatstartofbodyhook

\sphinxAtStartPar
A World of Girls: The Story of a School
&\begin{enumerate}
\sphinxsetlistlabels{\Alph}{enumi}{enumii}{}{.}%
\setcounter{enumi}{11}
\item {} \begin{enumerate}
\sphinxsetlistlabels{\Alph}{enumii}{enumiii}{}{.}%
\setcounter{enumii}{19}
\item {} 
\sphinxAtStartPar
Meade

\end{enumerate}

\end{enumerate}
\\
\sphinxhline
\sphinxAtStartPar
Adventures in Wallypug\sphinxhyphen{}Land
&\begin{enumerate}
\sphinxsetlistlabels{\Alph}{enumi}{enumii}{}{.}%
\setcounter{enumi}{6}
\item {} \begin{enumerate}
\sphinxsetlistlabels{\Alph}{enumii}{enumiii}{}{.}%
\setcounter{enumii}{4}
\item {} 
\sphinxAtStartPar
Farrow

\end{enumerate}

\end{enumerate}
\\
\sphinxhline
\sphinxAtStartPar
Alice’s Adventures in Wonderland
&
\sphinxAtStartPar
Lewis Carroll
\\
\sphinxhline
\sphinxAtStartPar
Allan Quatermain
&\begin{enumerate}
\sphinxsetlistlabels{\Alph}{enumi}{enumii}{}{.}%
\setcounter{enumi}{7}
\item {} 
\sphinxAtStartPar
Rider Haggard

\end{enumerate}
\\
\sphinxhline
\sphinxAtStartPar
Alone in London
&
\sphinxAtStartPar
Hesba Stretton
\\
\sphinxhline
\sphinxAtStartPar
At the Back of the North Wind
&
\sphinxAtStartPar
George McDonald
\\
\sphinxhline
\sphinxAtStartPar
Black Beauty
&
\sphinxAtStartPar
Anna Sewell
\\
\sphinxhline
\sphinxAtStartPar
Dream Days
&
\sphinxAtStartPar
Kenneth Grahame
\\
\sphinxhline
\sphinxAtStartPar
Eric; Or, Little by Little
&\begin{enumerate}
\sphinxsetlistlabels{\Alph}{enumi}{enumii}{}{.}%
\setcounter{enumi}{5}
\item {} \begin{enumerate}
\sphinxsetlistlabels{\Alph}{enumii}{enumiii}{}{.}%
\setcounter{enumii}{22}
\item {} 
\sphinxAtStartPar
Farrar

\end{enumerate}

\end{enumerate}
\\
\sphinxhline
\sphinxAtStartPar
Feats on the Fiord
&
\sphinxAtStartPar
Harriet Martineau
\\
\sphinxhline
\sphinxAtStartPar
Five Children and It
&\begin{enumerate}
\sphinxsetlistlabels{\Alph}{enumi}{enumii}{}{.}%
\setcounter{enumi}{4}
\item {} 
\sphinxAtStartPar
Nesbit

\end{enumerate}
\\
\sphinxhline
\sphinxAtStartPar
Holiday House: A Series of Tales
&
\sphinxAtStartPar
Catherine Sinclair
\\
\sphinxhline
\sphinxAtStartPar
Jackanapes
&
\sphinxAtStartPar
Juliana Horatia Ewing
\\
\sphinxhline
\sphinxAtStartPar
Jessica’s First Prayer
&
\sphinxAtStartPar
Hesba Stretton
\\
\sphinxhline
\sphinxAtStartPar
Kidnapped
&
\sphinxAtStartPar
Robert Louis Stevenson
\\
\sphinxhline
\sphinxAtStartPar
King Solomon’s Mines
&\begin{enumerate}
\sphinxsetlistlabels{\Alph}{enumi}{enumii}{}{.}%
\setcounter{enumi}{7}
\item {} 
\sphinxAtStartPar
Rider Haggard

\end{enumerate}
\\
\sphinxhline
\sphinxAtStartPar
Leila at home
&
\sphinxAtStartPar
Ann Fraser Tytler
\\
\sphinxhline
\sphinxAtStartPar
Little Meg’s children
&
\sphinxAtStartPar
Hesba Stretton
\\
\sphinxhline
\sphinxAtStartPar
Madam How and Lady Why; Or, First Lessons in Earth Lore for Children
&
\sphinxAtStartPar
Charles Kingsley
\\
\sphinxhline
\sphinxAtStartPar
Masterman Ready; Or, The Wreck of the Pacific
&
\sphinxAtStartPar
Frederick Marryat
\\
\sphinxhline
\sphinxAtStartPar
Moonfleet
&\begin{enumerate}
\sphinxsetlistlabels{\Alph}{enumi}{enumii}{}{.}%
\setcounter{enumi}{9}
\item {} 
\sphinxAtStartPar
Meade Falkner

\end{enumerate}
\\
\sphinxhline
\sphinxAtStartPar
Mopsa the Fairy
&
\sphinxAtStartPar
Jean Ingelow
\\
\sphinxhline
\sphinxAtStartPar
Mrs. Overtheway’s Remembrances
&
\sphinxAtStartPar
Juliana Horatia Ewing
\\
\sphinxhline
\sphinxAtStartPar
Nine Unlikely Tales
&\begin{enumerate}
\sphinxsetlistlabels{\Alph}{enumi}{enumii}{}{.}%
\setcounter{enumi}{4}
\item {} 
\sphinxAtStartPar
Nesbit

\end{enumerate}
\\
\sphinxhline
\sphinxAtStartPar
Peter Pan
&\begin{enumerate}
\sphinxsetlistlabels{\Alph}{enumi}{enumii}{}{.}%
\setcounter{enumi}{9}
\item {} \begin{enumerate}
\sphinxsetlistlabels{\Alph}{enumii}{enumiii}{}{.}%
\setcounter{enumii}{12}
\item {} 
\sphinxAtStartPar
Barrie

\end{enumerate}

\end{enumerate}
\\
\sphinxhline
\sphinxAtStartPar
Prince Prigio
&
\sphinxAtStartPar
Andrew Lang
\\
\sphinxhline
\sphinxAtStartPar
Stalky and Co
&
\sphinxAtStartPar
Rudyard Kipling
\\
\sphinxhline
\sphinxAtStartPar
The Book of Dragons
&\begin{enumerate}
\sphinxsetlistlabels{\Alph}{enumi}{enumii}{}{.}%
\setcounter{enumi}{4}
\item {} 
\sphinxAtStartPar
Nesbit

\end{enumerate}
\\
\sphinxhline
\sphinxAtStartPar
The Brass Bottle
&\begin{enumerate}
\sphinxsetlistlabels{\Alph}{enumi}{enumii}{}{.}%
\setcounter{enumi}{5}
\item {} 
\sphinxAtStartPar
Anstey

\end{enumerate}
\\
\sphinxhline
\sphinxAtStartPar
The Carved Lions
&
\sphinxAtStartPar
Mary Louisa Molesworth
\\
\sphinxhline
\sphinxAtStartPar
The Children of the New Forest
&
\sphinxAtStartPar
Frederick Marryat
\\
\sphinxhline
\sphinxAtStartPar
The Coral Island: A Tale of the Pacific Ocean
&\begin{enumerate}
\sphinxsetlistlabels{\Alph}{enumi}{enumii}{}{.}%
\setcounter{enumi}{17}
\item {} \begin{enumerate}
\sphinxsetlistlabels{\Alph}{enumii}{enumiii}{}{.}%
\setcounter{enumii}{12}
\item {} 
\sphinxAtStartPar
Ballantyne

\end{enumerate}

\end{enumerate}
\\
\sphinxhline
\sphinxAtStartPar
The Crofton Boys
&
\sphinxAtStartPar
Harriet Martineau
\\
\sphinxhline
\sphinxAtStartPar
The Cuckoo Clock
&
\sphinxAtStartPar
Mary Louisa Molesworth
\\
\sphinxhline
\sphinxAtStartPar
The Daisy Chain, or Aspirations
&
\sphinxAtStartPar
Charlotte M. Yonge
\\
\sphinxhline
\sphinxAtStartPar
The Dove in the Eagle’s Nest
&
\sphinxAtStartPar
Charlotte M. Yonge
\\
\sphinxhline
\sphinxAtStartPar
The Fifth Form at Saint Dominic’s: A School Story
&
\sphinxAtStartPar
Talbot Baines Reed
\\
\sphinxhline
\sphinxAtStartPar
The Golden Age
&
\sphinxAtStartPar
Kenneth Grahame
\\
\sphinxhline
\sphinxAtStartPar
The Happy Prince, and Other Tales
&
\sphinxAtStartPar
Oscar Wilde
\\
\sphinxhline
\sphinxAtStartPar
The Heir of Redclyffe
&
\sphinxAtStartPar
Charlotte M. Yonge
\\
\sphinxhline
\sphinxAtStartPar
The Jungle Book
&
\sphinxAtStartPar
Rudyard Kipling
\\
\sphinxhline
\sphinxAtStartPar
The King of the Golden River; or, the Black Brothers: A Legend of Stiria
&
\sphinxAtStartPar
John Ruskin
\\
\sphinxhline
\sphinxAtStartPar
The Little Duke: Richard the Fearless
&
\sphinxAtStartPar
Charlotte M. Yonge
\\
\sphinxhline
\sphinxAtStartPar
The Peasant and the Prince
&
\sphinxAtStartPar
Harriet Martineau
\\
\sphinxhline
\sphinxAtStartPar
The Princess and the Goblin
&
\sphinxAtStartPar
George McDonald
\\
\sphinxhline
\sphinxAtStartPar
The Railway Children
&\begin{enumerate}
\sphinxsetlistlabels{\Alph}{enumi}{enumii}{}{.}%
\setcounter{enumi}{4}
\item {} 
\sphinxAtStartPar
Nesbit

\end{enumerate}
\\
\sphinxhline
\sphinxAtStartPar
The Rival Crusoes; or The Shipwreck
&
\sphinxAtStartPar
Agnes Strickland
\\
\sphinxhline
\sphinxAtStartPar
The Rose and the Ring
&
\sphinxAtStartPar
William Makepeace Thackeray
\\
\sphinxhline
\sphinxAtStartPar
The Secret Garden
&
\sphinxAtStartPar
Frances Hodgson Burnett
\\
\sphinxhline
\sphinxAtStartPar
The Settlers at Home
&
\sphinxAtStartPar
Harriet Martineau
\\
\sphinxhline
\sphinxAtStartPar
The Settlers in Canada
&
\sphinxAtStartPar
Frederick Marryat
\\
\sphinxhline
\sphinxAtStartPar
The Story of the Amulet
&\begin{enumerate}
\sphinxsetlistlabels{\Alph}{enumi}{enumii}{}{.}%
\setcounter{enumi}{4}
\item {} 
\sphinxAtStartPar
Nesbit

\end{enumerate}
\\
\sphinxhline
\sphinxAtStartPar
The Story of the Treasure Seekers
&\begin{enumerate}
\sphinxsetlistlabels{\Alph}{enumi}{enumii}{}{.}%
\setcounter{enumi}{4}
\item {} 
\sphinxAtStartPar
Nesbit

\end{enumerate}
\\
\sphinxhline
\sphinxAtStartPar
The Surprising Adventures of Sir Toady Lion
&
\sphinxAtStartPar
S.R. Crockett
\\
\sphinxhline
\sphinxAtStartPar
The Tale of Benjamin Bunny
&
\sphinxAtStartPar
Beatrix Potter
\\
\sphinxhline
\sphinxAtStartPar
The Tale of Jemima Puddle\sphinxhyphen{}Duck
&
\sphinxAtStartPar
Beatrix Potter
\\
\sphinxhline
\sphinxAtStartPar
The Tale of Peter Rabbit
&
\sphinxAtStartPar
Beatrix Potter
\\
\sphinxhline
\sphinxAtStartPar
The Tale of Squirrel Nutkin
&
\sphinxAtStartPar
Beatrix Potter
\\
\sphinxhline
\sphinxAtStartPar
The Tale of the Flopsy Bunnies
&
\sphinxAtStartPar
Beatrix Potter
\\
\sphinxhline
\sphinxAtStartPar
The Tale of Two Bad Mice
&
\sphinxAtStartPar
Beatrix Potter
\\
\sphinxhline
\sphinxAtStartPar
The Tapestry Room: A Child’s Romance
&
\sphinxAtStartPar
Mary Louisa Molesworth
\\
\sphinxhline
\sphinxAtStartPar
The Three Mulla\sphinxhyphen{}mulgars
&
\sphinxAtStartPar
Walter de la Mare
\\
\sphinxhline
\sphinxAtStartPar
The Water\sphinxhyphen{}Babies
&
\sphinxAtStartPar
Charles Kingsley
\\
\sphinxhline
\sphinxAtStartPar
The Wind in the Willows
&
\sphinxAtStartPar
Kenneth Grahame
\\
\sphinxhline
\sphinxAtStartPar
Through the Looking\sphinxhyphen{}Glass
&
\sphinxAtStartPar
Lewis Carroll
\\
\sphinxhline
\sphinxAtStartPar
Tom Brown’s Schooldays
&
\sphinxAtStartPar
Thomas Hughes
\\
\sphinxhline
\sphinxAtStartPar
Treasure Island
&
\sphinxAtStartPar
Robert Louis Stevenson
\\
\sphinxhline
\sphinxAtStartPar
Vice Versa; or, A Lesson to Fathers
&\begin{enumerate}
\sphinxsetlistlabels{\Alph}{enumi}{enumii}{}{.}%
\setcounter{enumi}{5}
\item {} 
\sphinxAtStartPar
Anstey

\end{enumerate}
\\
\sphinxhline
\sphinxAtStartPar
Winning His Spurs. A Tale of the Crusades
&\begin{enumerate}
\sphinxsetlistlabels{\Alph}{enumi}{enumii}{}{.}%
\setcounter{enumi}{6}
\item {} \begin{enumerate}
\sphinxsetlistlabels{\Alph}{enumii}{enumiii}{}{.}%
\item {} 
\sphinxAtStartPar
Henty

\end{enumerate}

\end{enumerate}
\\
\sphinxhline
\sphinxAtStartPar
With Clive in India; Or, The Beginnings of an Empire
&\begin{enumerate}
\sphinxsetlistlabels{\Alph}{enumi}{enumii}{}{.}%
\setcounter{enumi}{6}
\item {} \begin{enumerate}
\sphinxsetlistlabels{\Alph}{enumii}{enumiii}{}{.}%
\item {} 
\sphinxAtStartPar
Henty

\end{enumerate}

\end{enumerate}
\\
\sphinxhline
\sphinxAtStartPar
Wood Magic, a Fable
&
\sphinxAtStartPar
Richard Jefferies
\\
\sphinxbottomrule
\end{longtable}
\sphinxtableafterendhook
\sphinxatlongtableend
\end{savenotes}


\begin{savenotes}
\sphinxatlongtablestart
\sphinxthistablewithglobalstyle
\makeatletter
  \LTleft \@totalleftmargin plus1fill
  \LTright\dimexpr\columnwidth-\@totalleftmargin-\linewidth\relax plus1fill
\makeatother
\begin{longtable}{\X{60}{100}\X{40}{100}}
\sphinxthelongtablecaptionisattop
\caption{ArTs \textendash{} Additional Requested Texts (31 texts)\strut}\label{\detokenize{appendices:id5}}\\*[\sphinxlongtablecapskipadjust]
\sphinxtoprule
\sphinxstyletheadfamily 
\sphinxAtStartPar
Title
&\sphinxstyletheadfamily 
\sphinxAtStartPar
Author
\\
\sphinxmidrule
\endfirsthead

\multicolumn{2}{c}{\sphinxnorowcolor
    \makebox[0pt]{\sphinxtablecontinued{\tablename\ \thetable{} \textendash{} continued from previous page}}%
}\\
\sphinxtoprule
\sphinxstyletheadfamily 
\sphinxAtStartPar
Title
&\sphinxstyletheadfamily 
\sphinxAtStartPar
Author
\\
\sphinxmidrule
\endhead

\sphinxbottomrule
\multicolumn{2}{r}{\sphinxnorowcolor
    \makebox[0pt][r]{\sphinxtablecontinued{continues on next page}}%
}\\
\endfoot

\endlastfoot
\sphinxtableatstartofbodyhook

\sphinxAtStartPar
A Christmas Carol: A Ghost Story of Christmas
&
\sphinxAtStartPar
Charles Dickens
\\
\sphinxhline
\sphinxAtStartPar
A Room with a View
&\begin{enumerate}
\sphinxsetlistlabels{\Alph}{enumi}{enumii}{}{.}%
\setcounter{enumi}{4}
\item {} \begin{enumerate}
\sphinxsetlistlabels{\Alph}{enumii}{enumiii}{}{.}%
\setcounter{enumii}{12}
\item {} 
\sphinxAtStartPar
Forster

\end{enumerate}

\end{enumerate}
\\
\sphinxhline
\sphinxAtStartPar
Adventures of Huckleberry Finn, Complete
&
\sphinxAtStartPar
Mark Twain
\\
\sphinxhline
\sphinxAtStartPar
American Notes for General Circulation
&
\sphinxAtStartPar
Charles Dickens
\\
\sphinxhline
\sphinxAtStartPar
Gulliver’s Travels
&
\sphinxAtStartPar
Jonathan Swift
\\
\sphinxhline
\sphinxAtStartPar
Heart of Darkness
&
\sphinxAtStartPar
Joseph Conrad
\\
\sphinxhline
\sphinxAtStartPar
Lady Susan
&
\sphinxAtStartPar
Jane Austen
\\
\sphinxhline
\sphinxAtStartPar
Mansfield Park
&
\sphinxAtStartPar
Jane Austen
\\
\sphinxhline
\sphinxAtStartPar
Middlemarch
&
\sphinxAtStartPar
George Eliot
\\
\sphinxhline
\sphinxAtStartPar
Northanger Abbey
&
\sphinxAtStartPar
Jane Austen
\\
\sphinxhline
\sphinxAtStartPar
Pictures from Italy
&
\sphinxAtStartPar
Charles Dickens
\\
\sphinxhline
\sphinxAtStartPar
Sense and Sensibility
&
\sphinxAtStartPar
Jane Austen
\\
\sphinxhline
\sphinxAtStartPar
Shirley
&
\sphinxAtStartPar
Charlotte Brontë
\\
\sphinxhline
\sphinxAtStartPar
Silas Marner: The Weaver of Raveloe
&
\sphinxAtStartPar
George Eliot
\\
\sphinxhline
\sphinxAtStartPar
Sketches by Boz: Illustrative of Every\sphinxhyphen{}Day Life and Every\sphinxhyphen{}Day People
&
\sphinxAtStartPar
Charles Dickens
\\
\sphinxhline
\sphinxAtStartPar
The Awakening and Selected Short Stories
&
\sphinxAtStartPar
Kate Chopin
\\
\sphinxhline
\sphinxAtStartPar
The Jungle
&
\sphinxAtStartPar
Upton Sinclair
\\
\sphinxhline
\sphinxAtStartPar
The Moonstone
&
\sphinxAtStartPar
Wilkie Collins
\\
\sphinxhline
\sphinxAtStartPar
The Portrait of a Lady Volume 1 (of 2)
&
\sphinxAtStartPar
Henry James
\\
\sphinxhline
\sphinxAtStartPar
The Portrait of a Lady Volume 2 (of 2)
&
\sphinxAtStartPar
Henry James
\\
\sphinxhline
\sphinxAtStartPar
The Return of the Soldier
&
\sphinxAtStartPar
Rebecca West
\\
\sphinxhline
\sphinxAtStartPar
The Sign of the Four
&
\sphinxAtStartPar
Arthur Conan Doyle
\\
\sphinxhline
\sphinxAtStartPar
The Tenant of Wildfell Hall
&
\sphinxAtStartPar
Anne Brontë
\\
\sphinxhline
\sphinxAtStartPar
The Time Machine: An Invention
&\begin{enumerate}
\sphinxsetlistlabels{\Alph}{enumi}{enumii}{}{.}%
\setcounter{enumi}{7}
\item {} \begin{enumerate}
\sphinxsetlistlabels{\Alph}{enumii}{enumiii}{}{.}%
\setcounter{enumii}{6}
\item {} 
\sphinxAtStartPar
Wells

\end{enumerate}

\end{enumerate}
\\
\sphinxhline
\sphinxAtStartPar
The Uncommercial Traveller
&
\sphinxAtStartPar
Charles Dickens
\\
\sphinxhline
\sphinxAtStartPar
The War of the Worlds
&\begin{enumerate}
\sphinxsetlistlabels{\Alph}{enumi}{enumii}{}{.}%
\setcounter{enumi}{7}
\item {} \begin{enumerate}
\sphinxsetlistlabels{\Alph}{enumii}{enumiii}{}{.}%
\setcounter{enumii}{6}
\item {} 
\sphinxAtStartPar
Wells

\end{enumerate}

\end{enumerate}
\\
\sphinxhline
\sphinxAtStartPar
The Yellow Wallpaper
&
\sphinxAtStartPar
Charlotte Perkins Gilman
\\
\sphinxhline
\sphinxAtStartPar
Twelve Years a Slave
&
\sphinxAtStartPar
Solomon Northup
\\
\sphinxhline
\sphinxAtStartPar
Villette
&
\sphinxAtStartPar
Charlotte Brontë
\\
\sphinxhline
\sphinxAtStartPar
What Maisie Knew
&
\sphinxAtStartPar
Henry James
\\
\sphinxhline
\sphinxAtStartPar
Women in Love
&\begin{enumerate}
\sphinxsetlistlabels{\Alph}{enumi}{enumii}{}{.}%
\setcounter{enumi}{3}
\item {} \begin{enumerate}
\sphinxsetlistlabels{\Alph}{enumii}{enumiii}{}{.}%
\setcounter{enumii}{7}
\item {} 
\sphinxAtStartPar
Lawrence

\end{enumerate}

\end{enumerate}
\\
\sphinxbottomrule
\end{longtable}
\sphinxtableafterendhook
\sphinxatlongtableend
\end{savenotes}


\begin{savenotes}
\sphinxatlongtablestart
\sphinxthistablewithglobalstyle
\makeatletter
  \LTleft \@totalleftmargin plus1fill
  \LTright\dimexpr\columnwidth-\@totalleftmargin-\linewidth\relax plus1fill
\makeatother
\begin{longtable}{\X{50}{100}\X{30}{100}\X{20}{100}}
\sphinxthelongtablecaptionisattop
\caption{DE19 \textendash{} German 19th Century Reference Corpus (33 texts)\strut}\label{\detokenize{appendices:id6}}\\*[\sphinxlongtablecapskipadjust]
\sphinxtoprule
\sphinxstyletheadfamily 
\sphinxAtStartPar
Title
&\sphinxstyletheadfamily 
\sphinxAtStartPar
Author
&\sphinxstyletheadfamily 
\sphinxAtStartPar
Year
\\
\sphinxmidrule
\endfirsthead

\multicolumn{3}{c}{\sphinxnorowcolor
    \makebox[0pt]{\sphinxtablecontinued{\tablename\ \thetable{} \textendash{} continued from previous page}}%
}\\
\sphinxtoprule
\sphinxstyletheadfamily 
\sphinxAtStartPar
Title
&\sphinxstyletheadfamily 
\sphinxAtStartPar
Author
&\sphinxstyletheadfamily 
\sphinxAtStartPar
Year
\\
\sphinxmidrule
\endhead

\sphinxbottomrule
\multicolumn{3}{r}{\sphinxnorowcolor
    \makebox[0pt][r]{\sphinxtablecontinued{continues on next page}}%
}\\
\endfoot

\endlastfoot
\sphinxtableatstartofbodyhook

\sphinxAtStartPar
Die Günderode
&
\sphinxAtStartPar
Arnim, B.
&
\sphinxAtStartPar
1840
\\
\sphinxhline
\sphinxAtStartPar
Clemens Brentanos Frühlingskranz
&
\sphinxAtStartPar
Arnim, B.
&
\sphinxAtStartPar
1844
\\
\sphinxhline
\sphinxAtStartPar
Fanny Förster
&
\sphinxAtStartPar
Boy\sphinxhyphen{}Ed, I.
&
\sphinxAtStartPar
1889
\\
\sphinxhline
\sphinxAtStartPar
Kampf um Rom
&
\sphinxAtStartPar
Dahn, F.
&
\sphinxAtStartPar
1876
\\
\sphinxhline
\sphinxAtStartPar
Sibilla Dalmar
&
\sphinxAtStartPar
Dohm, H.
&
\sphinxAtStartPar
1896
\\
\sphinxhline
\sphinxAtStartPar
Schicksale einer Seele
&
\sphinxAtStartPar
Dohm, H.
&
\sphinxAtStartPar
1899
\\
\sphinxhline
\sphinxAtStartPar
Bozena
&
\sphinxAtStartPar
Ebner\sphinxhyphen{}Eschenbach, M.
&
\sphinxAtStartPar
1876
\\
\sphinxhline
\sphinxAtStartPar
Vor dem Sturm
&
\sphinxAtStartPar
Fontane, T.
&
\sphinxAtStartPar
1878
\\
\sphinxhline
\sphinxAtStartPar
Der Stechlin
&
\sphinxAtStartPar
Fontane, T.
&
\sphinxAtStartPar
1897/98
\\
\sphinxhline
\sphinxAtStartPar
Die letzte Reckenburgerin
&
\sphinxAtStartPar
François, Louise.
&
\sphinxAtStartPar
1871
\\
\sphinxhline
\sphinxAtStartPar
Soll und Haben
&
\sphinxAtStartPar
Freytag, G.
&
\sphinxAtStartPar
1855
\\
\sphinxhline
\sphinxAtStartPar
Die verlorene Handschrift
&
\sphinxAtStartPar
Freytag, G.
&
\sphinxAtStartPar
1864
\\
\sphinxhline
\sphinxAtStartPar
Schloß Hubertus
&
\sphinxAtStartPar
Ganghofer, L.
&
\sphinxAtStartPar
1895
\\
\sphinxhline
\sphinxAtStartPar
Sibylle
&
\sphinxAtStartPar
Hahn\sphinxhyphen{}Hahn, I. G.
&
\sphinxAtStartPar
1846
\\
\sphinxhline
\sphinxAtStartPar
Der Sonnenwirt
&
\sphinxAtStartPar
Kurz, H.
&
\sphinxAtStartPar
1855
\\
\sphinxhline
\sphinxAtStartPar
Jenny
&
\sphinxAtStartPar
Lewald, F.
&
\sphinxAtStartPar
1843
\\
\sphinxhline
\sphinxAtStartPar
Das Heideprinzeßchen
&
\sphinxAtStartPar
Marlitt, E.
&
\sphinxAtStartPar
1871/72
\\
\sphinxhline
\sphinxAtStartPar
Goldelse
&
\sphinxAtStartPar
Marlitt, E.
&
\sphinxAtStartPar
1866/67
\\
\sphinxhline
\sphinxAtStartPar
Das Geheimnis der alten Mamsell
&
\sphinxAtStartPar
Marlitt, E.
&
\sphinxAtStartPar
1868
\\
\sphinxhline
\sphinxAtStartPar
Am Jenseits
&
\sphinxAtStartPar
May, K.
&
\sphinxAtStartPar
1898
\\
\sphinxhline
\sphinxAtStartPar
Jürg Jenatsch
&
\sphinxAtStartPar
Meyer, C. F.
&
\sphinxAtStartPar
1874
\\
\sphinxhline
\sphinxAtStartPar
Der Büttnerbauer
&
\sphinxAtStartPar
Polenz, W.
&
\sphinxAtStartPar
1895
\\
\sphinxhline
\sphinxAtStartPar
Der Schüdderump
&
\sphinxAtStartPar
Raabe, W.
&
\sphinxAtStartPar
1869/70
\\
\sphinxhline
\sphinxAtStartPar
Der Hungerpastor
&
\sphinxAtStartPar
Raabe, W.
&
\sphinxAtStartPar
1863/64
\\
\sphinxhline
\sphinxAtStartPar
Die Leute aus dem Walde, ihre Sterne, Wege und Schicksale
&
\sphinxAtStartPar
Raabe, W.
&
\sphinxAtStartPar
1863
\\
\sphinxhline
\sphinxAtStartPar
Anna
&
\sphinxAtStartPar
Schopenhauer, A.
&
\sphinxAtStartPar
1845
\\
\sphinxhline
\sphinxAtStartPar
Problematische Naturen. Erste Abtheilung
&
\sphinxAtStartPar
Spielhagen, F.
&
\sphinxAtStartPar
1861
\\
\sphinxhline
\sphinxAtStartPar
Hammer und Amboß
&
\sphinxAtStartPar
Spielhagen, F.
&
\sphinxAtStartPar
1869
\\
\sphinxhline
\sphinxAtStartPar
Die Mappe meines Urgroßvaters
&
\sphinxAtStartPar
Stifter, A.
&
\sphinxAtStartPar
1841
\\
\sphinxhline
\sphinxAtStartPar
Witiko
&
\sphinxAtStartPar
Stifter, A.
&
\sphinxAtStartPar
1855, 1865\textendash{}1867
\\
\sphinxhline
\sphinxAtStartPar
Vittoria Accorombona
&
\sphinxAtStartPar
Tieck, L.
&
\sphinxAtStartPar
1840
\\
\sphinxhline
\sphinxAtStartPar
Auch Einer 1
&
\sphinxAtStartPar
Vischer, F. T.
&
\sphinxAtStartPar
1879
\\
\sphinxhline
\sphinxAtStartPar
Auch Einer 2
&
\sphinxAtStartPar
Vischer, F. T.
&
\sphinxAtStartPar
1879
\\
\sphinxbottomrule
\end{longtable}
\sphinxtableafterendhook
\sphinxatlongtableend
\end{savenotes}
\subsubsection*{CLiC texts listed in A\sphinxhyphen{}Level and GCSE specifications}

\begin{figure}[htbp]
\centering
\capstart

\noindent\sphinxincludegraphics{{appendices-alevelgcsespecs-overview}.png}
\caption{Overview of CLiC texts listed in the AQA, Edexcel and OCR
A\sphinxhyphen{}Level / GCSE specifications}\label{\detokenize{appendices:id7}}\end{figure}
\phantomsection\label{\detokenize{glossary::doc}}

\chapter{Glossary of terms used}
\label{\detokenize{glossary:glossary-of-terms-used}}\begin{description}
\sphinxlineitem{book:}
\sphinxAtStartPar
A UTF\sphinxhyphen{}8 encoded text file that has been imported into CLiC, invariably from the corpora repository.

\sphinxlineitem{region:}
\sphinxAtStartPar
A labelled portion of a book. Each region will have:
\begin{itemize}
\item {} 
\sphinxAtStartPar
A name (or rclass), for example ‘chapter.sentence’. For all possible names, see \sphinxurl{https://github.com/mahlberg-lab/clic/blob/HEAD//schema/10-rclass.sql}.

\item {} 
\sphinxAtStartPar
A start and end character position within the full book text

\item {} 
\sphinxAtStartPar
(optionally) a number (or rvalue), for example it’s position within a chapter.

\end{itemize}

\sphinxAtStartPar
Regions are added by \sphinxstyleemphasis{region tagger} scripts, which are in \sphinxcode{\sphinxupquote{clic.region}}.

\sphinxlineitem{token:}
\sphinxAtStartPar
A token is a labelled portion of a book that contains a single word.
In the phrase \sphinxcode{\sphinxupquote{\textquotesingle{}For \_more\_!\textquotesingle{} said Mr. Limbkins.}}, \sphinxcode{\sphinxupquote{For}}, \sphinxcode{\sphinxupquote{more}}, \sphinxcode{\sphinxupquote{said}}, \sphinxcode{\sphinxupquote{Mr}}, \sphinxcode{\sphinxupquote{Limbkins}} would be tokens.

\sphinxAtStartPar
See \sphinxcode{\sphinxupquote{clic.tokenizer}}.

\sphinxlineitem{type/ttype:}
\sphinxAtStartPar
A token has a type (thus ttype). This is a normalised form of the token.
The tokens in the phrase \sphinxcode{\sphinxupquote{\textquotesingle{}For \_more\_!\textquotesingle{} said Mr. Limbkins.}} would have types \sphinxcode{\sphinxupquote{for}}, \sphinxcode{\sphinxupquote{more}}, \sphinxcode{\sphinxupquote{said}}, \sphinxcode{\sphinxupquote{mr}}, \sphinxcode{\sphinxupquote{limbkins}}.

\sphinxAtStartPar
See \sphinxcode{\sphinxupquote{clic.tokenizer}}.

\end{description}

\begin{sphinxthebibliography}{GitHub\_c}
\bibitem[CLiC\_KWICGrouper\_video]{references:clic-kwicgrouper-video}
\sphinxAtStartPar
Wiegand, V. \& Guglielmi, A. (2017, June 22). Introduction to KWICGrouper video. \sphinxurl{https://blog.bham.ac.uk/clic-dickens/2017/06/22/video-introducing-the-clic-kwicgrouper-function-to-group-concordance-lines/}
\bibitem[CLiC\_project\_page]{references:clic-project-page}
\sphinxAtStartPar
\sphinxurl{https://clic-fiction.com}
\bibitem[GitHub\_CLiC]{references:github-clic}
\sphinxAtStartPar
\sphinxurl{https://github.com/mahlberg-lab/clic}
\bibitem[GitHub\_corpora]{references:github-corpora}
\sphinxAtStartPar
\sphinxurl{https://github.com/mahlberg-lab/corpora}
\bibitem[GitHub\_corpora\_initial\_ArTs]{references:github-corpora-initial-arts}
\sphinxAtStartPar
\sphinxurl{https://github.com/mahlberg-lab/corpora/tree/026a8436bf9ea3282d283a05725c0153e023d74c/Other}
\bibitem[GitHub\_corpora\_initial\_ChiLit]{references:github-corpora-initial-chilit}
\sphinxAtStartPar
\sphinxurl{https://github.com/mahlberg-lab/corpora/tree/a020b2a7153baf8849056be833861ecb3d77e7a1/ChiLit}
\bibitem[GLARE\_project\_page]{references:glare-project-page}
\sphinxAtStartPar
\sphinxurl{http://birmingham.ac.uk/GLARE}
\bibitem[ICU]{references:icu}
\sphinxAtStartPar
\sphinxurl{http://userguide.icu-project.org/boundaryanalysis}
\bibitem[ICU\_RSV]{references:icu-rsv}
\sphinxAtStartPar
\sphinxurl{http://userguide.icu-project.org/boundaryanalysis\#TOC-Rule-Status-Values}
\bibitem[Mahlberg\_2013]{references:mahlberg-2013}
\sphinxAtStartPar
Mahlberg, M. (2013). Corpus Stylistics and Dickens’s Fiction. London: Routledge.
\bibitem[Mahlberg\_et\_al\_2016]{references:mahlberg-et-al-2016}
\sphinxAtStartPar
Mahlberg, M., Stockwell, P., de Joode, J., Smith, C., \& O’Donnell, M. B. (2016). CLiC Dickens: novel uses of concordances for the integration of corpus stylistics and cognitive poetics. Corpora, 11(3), 433\textendash{}463 {[}Open access, available from \sphinxurl{https://doi.org/10.3366/cor.2016.0102}{]}
\bibitem[Rayson\_Garside\_2000]{references:rayson-garside-2000}
\sphinxAtStartPar
Rayson, P. and Garside, R. (2000). Comparing corpora using frequency profiling. In proceedings of the workshop on Comparing Corpora, held in conjunction with the 38th annual meeting of the Association for Computational Linguistics (ACL 2000). 1\sphinxhyphen{}8 October 2000, Hong Kong, pp. 1\sphinxhyphen{}6, retrieved from \sphinxurl{http://ucrel.lancs.ac.uk/people/paul/publications/rg\_acl2000.pdf}
\bibitem[UAX29]{references:uax29}
\sphinxAtStartPar
\sphinxurl{https://www.unicode.org/reports/tr29/tr29-33.html\#Word\_Boundaries}
\bibitem[UNIDECODE]{references:unidecode}
\sphinxAtStartPar
\sphinxurl{https://pypi.org/project/Unidecode/}
\end{sphinxthebibliography}



\renewcommand{\indexname}{Index}
\printindex
\end{document}